\documentclass{amsart}
\usepackage[utf8]{inputenc}
\usepackage{lmodern}
\usepackage{amssymb,mathtools,array,longtable,tabularx,booktabs}
\usepackage{mathrsfs,enumitem,graphicx}
\usepackage[hidelinks]{hyperref}
\hypersetup{pdftitle={Shifted primes in arithmetic progressions to large moduli with linear sieve remainder weights},pdfauthor={Ruiyi Zhang}}
\allowdisplaybreaks[3]
\setlength{\emergencystretch}{3em}
\setcounter{tocdepth}{2}
\providecommand{\tightlist}{\setlength{\itemsep}{0pt}\setlength{\parskip}{0pt}}
\providecommand{\passthrough}[1]{#1}
\providecommand{\texorpdfstring}[2]{#1}
\providecommand{\PP}{\mathbb P}
\providecommand{\UC}{U_C}
\providecommand{\NN}{\mathbb N}
\providecommand{\ZZ}{\mathbb Z}
\providecommand{\QQ}{\mathbb Q}
\providecommand{\RR}{\mathbb R}
\providecommand{\CC}{\mathbb C}
\DeclareMathOperator{\Li}{Li}
\DeclareMathOperator{\lcm}{lcm}
\DeclareMathOperator{\rad}{rad}
\DeclareMathOperator{\supp}{supp}
\DeclareUnicodeCharacter{2212}{\ensuremath{-}}
\DeclareUnicodeCharacter{2264}{\ensuremath{\le}}
\DeclareUnicodeCharacter{2265}{\ensuremath{\ge}}
\DeclareUnicodeCharacter{2260}{\ensuremath{\ne}}
\DeclareUnicodeCharacter{221E}{\ensuremath{\infty}}
\DeclareUnicodeCharacter{2208}{\ensuremath{\in}}
\DeclareUnicodeCharacter{2286}{\ensuremath{\subseteq}}
\DeclareUnicodeCharacter{00D7}{\ensuremath{\times}}
\DeclareUnicodeCharacter{2211}{\ensuremath{\sum}}
\DeclareUnicodeCharacter{220F}{\ensuremath{\prod}}
\DeclareUnicodeCharacter{220E}{\ensuremath{\square}}
\DeclareUnicodeCharacter{2032}{\ensuremath{\prime}}
\DeclareUnicodeCharacter{2011}{-}
\DeclareUnicodeCharacter{2013}{--}
\DeclareUnicodeCharacter{2014}{---}
\newtheorem{theorem}{Theorem}[section]
\newtheorem{lemma}[theorem]{Lemma}
\newtheorem{proposition}[theorem]{Proposition}
\newtheorem{corollary}[theorem]{Corollary}
\numberwithin{equation}{section}
\theoremstyle{definition}
\newtheorem{definition}[theorem]{Definition}
\newtheorem{hypothesis}[theorem]{Hypothesis}
\theoremstyle{plain}
\newtheorem*{introproposition}{Proposition}
\title[Shifted primes with linear sieve remainder weights]{Shifted primes in arithmetic progressions to large moduli with linear sieve remainder weights}
\author{Ruiyi Zhang}
\address{Independent Researcher, Columbia, Maryland, USA}
\email{ryzfsu@gmail.com}
\subjclass[2020]{Primary 11N13; Secondary 11N36, 11L05, 11P32}
\keywords{shifted primes, linear sieve, level of distribution, dispersion method, Kloosterman sums, primes in arithmetic progressions}
\begin{document}
\begin{abstract}
Let $N$ be a sufficiently large even integer, and fix $B>0$ and
$0<\varepsilon\le10^{-4}$. Take Iwaniec's upper or lower
linear-sieve remainder weights. We estimate a signed sum of
prime-counting discrepancies on $N/(\log N)^B<p<(1-\varepsilon)N$. Fix $0<u\le1/500$ and a sieve exponent
$1/2<\theta<(5-7/32-3u)/8$. Choose
$0<\sigma\le\min\{10^{-10},((5-7/32-3u)/8-\theta)/200,
\sqrt{u/(2\theta)}\}$. Then the complete modulus support lies below
$N^{(5-7/32-3u)/8}$. For every fixed saving exponent $A>0$, the error
is $O(N(\log N)^{-A})$. For each of the two signs $\nu\in\{+,-\}$, the
result holds uniformly for integers $a\in[N,N\log N/3]$ satisfying
$(a,q)=1$ for every modulus
$q$ in the support of the corresponding coefficient $w_N^\nu$. Call these shifts
admissible; they form a subset of the interval that includes $a=N$.
When $u=1/500$ and $\theta=19101/32000-10^{-4}$, the ratio
$\log D/\log z$ is $298.403125$.
The proof keeps the smooth-divisor and prime-tuple factors of the
linear-sieve weights. It combines dispersion with a Heath--Brown
decomposition. The boundary, Type-I and Type-II terms are estimated
separately.
\end{abstract}
\maketitle

\section{Introduction}
The Bombieri--Vinogradov theorem \cite{Bombieri65,Vinogradov65}
gives level $1/2$ for primes in arithmetic progressions on average.
Related early weighted work includes Fouvry--Iwaniec
\cite{FouvryIwaniec80,FouvryIwaniec83}.
Consider weighted sums with a fixed residue.
Bombieri--Friedlander--Iwaniec reached level $4/7$ with well-factorable
weights \cite[Theorem~10]{BFI86}.
With a fixed residue, Maynard \cite[Theorems~1.1--1.2]{MaynardII}
proved level $3/5$ for triply well-factorable weights and level $7/12$
for the upper linear-sieve weights. Parts I--III address fixed residue
classes, well-factorable estimates, and uniform residue classes, respectively
\cite{MaynardI,MaynardII,MaynardIII}.
Lichtman \cite[Theorem~1.7, (1.6)]{Lichtman23} obtained level $66/107$ for
triply well-factorable weights with a fixed residue. For such weights
he also obtained level $153/256-\delta$, uniformly over residues
$|a|<N^{1+\delta}$ for each fixed $\delta>0$
\cite[Theorem~1.7, (1.7)]{Lichtman23}. His Proposition~6.6 treats
Iwaniec's modified upper and lower linear-sieve weights. There the
largest prime factors are restricted to prescribed short intervals.
Its admissible level is
$\vartheta\le\boldsymbol{\vartheta}(t_1)-\delta$ with
$\boldsymbol{\vartheta}(t_1)$ defined in \cite[(6.2)]{Lichtman23}.
Take at least three prescribed factors and $t_1$ below the threshold
in that proposition. Then the level becomes
$\vartheta\le\boldsymbol{\vartheta}(t_1,t_2,t_3)-\delta$, with the
function of \cite[(6.4)]{Lichtman23} and fixed $\delta>0$.
The error constant in that proposition depends on the
fixed residue $a$ \cite[Proposition~6.6]{Lichtman23}. Lichtman's
modified linear-sieve weights reach level $10/17$ \cite{Lichtman21}.

\pagebreak
Pascadi \cite[Theorem~1.3]{Pascadi25} gives levels approaching $5/8$
for triply well-factorable weights, with a fixed residue. For suitable
upper linear-sieve weights his levels approach $3/5$.
Pascadi's Theorem~5.4 also treats both signs of the restricted
well-factorable linear-sieve weights and refines Proposition~6.6
\cite[Theorem~5.4]{Pascadi25}. His Lemma~5.6 records a correction to a
small error in Section~6 of \cite{Lichtman23}
\cite[Lemma~5.6]{Pascadi25}.

Here $z=N^u$ with $u\le1/500$ and $\theta>1/2$, so
$s=\theta/u>250$, in the fundamental-lemma range. The cited
linear-sieve levels use sifting ranges with $s\approx1$--$2$ and are
not directly comparable.

The congruence $p\equiv a\pmod q$ says that $q$ divides the
shifted prime $a-p$. So the sums below count shifted primes
$a-p$ in the class $0$ modulo $q$.
The modulus coefficients are Iwaniec's explicit lower and upper
remainder weights \cite[Theorem~4]{AA37}.
We retain their smooth-divisor factors, prime-tuple factors, and
distinctness conditions. We estimate the complete signed sum for
the admissible shifts in the stated interval. Theorem~\ref{thm:aa37-interior}
covers the fixed parameters $u,\theta,A,B,\varepsilon$ below.
Its support exponent tends to $153/256$ as $u\downarrow0$. This numerical
endpoint matches Lichtman's moving-residue level for triply well-factorable
weights \cite[Theorem~1.7, (1.7)]{Lichtman23}; the residue ranges differ.
At $u=1/500$ the exponent is smaller.
The coefficient family here is the full remainder in
\cite[Theorem~4]{AA37}. It includes the smooth-only term and the
distinctness masks. Lichtman's Proposition~6.6 estimates restricted
factor-pattern components of different modified sieve weights. Its
implied constant depends on the fixed residue. Here uniformity is only over
the support-coprime shifts specified below, an admissible subset of
$[N,N\log N/3]$ that includes $a=N$; no uniformity over the full interval
is asserted.

The local statement uses the hypotheses of Section~\ref{sec:pi-psdt}.
Let $B$ be a Heath--Brown box at scale $X$. Write $w_B(n)$ for its
finite source coefficient and $C_B(q)$ for its modulus coefficient.
Set
\[
 \Delta_B(q;a)=\sum_{n}w_B(n)
 \left(\mathbf1_{n\equiv a\, (q)}
       -\frac{\mathbf1_{(n,q)=1}}{\varphi(q)}\right).
\]
In Section~\ref{sec:pi-psdt}, $w_B$ is written as a sum over the
variables of $B$; this is the same discrepancy.

\begin{introproposition}[Dispersion for a source box]
Fix $0<\sigma\le10^{-10}$ and constants $K_{\rm scale}>0$ and
$C_{\rm box}>0$. Here $K_{\rm scale}$ bounds the logarithmic
difference between the prime scale $X$ and $N$. The source boxes have
relative width $(\log X)^{-C_{\rm box}}$.
Let $N$ be sufficiently large and even,
$N/(\log N)^{K_{\rm scale}}\le X\le N$, and put
$L=\log X$ and $A_0=X^{1+\sigma/20}$.
Let $a$ be a positive integer with
$a\le X(\log N)^{K_{\rm scale}}\le A_0$.
Suppose every supported modulus is coprime to $2a$.
Suppose $B$ and $C_B$ satisfy (H1)--(H5) of
Section~\ref{sec:pi-psdt}. These include the support condition
$q\le X^{\vartheta_0}$, where
\[
 \eta_L=10^{-4},\quad
 \vartheta_0=\frac{19101}{32000}-\eta_L=0.59680625,
 \quad \ell_0=\frac{53}{412}-\frac{\eta_L}2,
 \quad c_\sigma=\frac{\ell_0\sigma^2}{4}>0.
\]
For each fixed $A'>0$ there is a finite
$C=C(\sigma,K_{\rm scale},C_{\rm box},A',\mathcal C_{\rm src})$
such that
\begin{equation}\label{eq:intro-dispersion}
 \left|\sum_{q\le X^{\vartheta_0}}
       C_B(q)\Delta_B(q;a)\right|
 \ll_{\sigma,K_{\rm scale},C_{\rm box},A',\mathcal C_{\rm src}}
       XL^{-A'}+X^{1-c_\sigma}L^C.
\end{equation}
Here $\mathcal C_{\rm src}$ denotes the fixed uniform constants in
(H1)--(H5). The estimate is uniform in $B$, $N$, and all $a$ as above.
The implied constant and the threshold for $N$ are ineffective.
\end{introproposition}

Theorem~\ref{thm:standalone} in Section~\ref{sec:pi-psdt}
is the abstract dispersion estimate under the support and packing
hypotheses in (H3).
Theorem~\ref{thm:main} is the dispersion engine used for
Theorem~\ref{thm:aa37-interior}. It allows support up to $X^{0.6}$.

\begin{theorem}[Signed shifted-prime mean value at $u=1/500$]
\label{thm:aa37-fixed-u}
Fix $A,B>0$, $0<\varepsilon\le10^{-4}$, and
\[
 \frac12<\theta<\Theta_*:=\frac{19101}{32000}.
\]
Set
\[
\begin{gathered}
 u=\frac1{500},\qquad \eta_\Theta=\Theta_*-\theta,\qquad
 0<\sigma\le\min\left\{10^{-10},\frac{\eta_\Theta}{200},
                         \sqrt{\frac1{1000\theta}}\right\},\\
 \tau=\sigma^9,\quad D=N^\theta,\quad z=N^{1/500},
 \quad Y=D^{\sigma^2},\quad K_{\rm scale}=B+1,
\end{gathered}
\]
and choose an integer
$C_{\rm box}\ge\lceil2A+32\rceil$; constants may depend on this choice,
as in Theorem~\ref{thm:aa37-interior}.
Let $N$ be sufficiently large and even. For either sign
$\nu\in\{+,-\}$, let $w_N^\nu(q)$ be the complete Iwaniec remainder
coefficient of formula~\textup{(W)} in
Section~\ref{sec:linear-sieve-application}. Its prime set is
$\mathcal P_N=\{p\text{ prime}:p\nmid2N\}$, with the grid and
distinctness conditions there. Let $a$ be an integer with
$N\le a\le N\log N/3$ and $(a,q)=1$ for every
$q\in\operatorname{supp}w_N^\nu$.
Call these shifts admissible; they form a subset of the interval that includes
$a=N$.
Put $I_N=(N/(\log N)^B,(1-\varepsilon)N)$. Let $\pi_{I_N}$ count
its primes, and $\pi_{I_N}(q,a)$ those congruent to $a$ modulo $q$.
Then
\[
 \operatorname{supp}w_N^\nu\subseteq
 \{q<D^{1+\sigma+\tau}:(q,2N)=1\},\qquad
 D^{1+\sigma+\tau}<N^{\Theta_*},
 \qquad |w_N^\nu(q)|\ll_{\sigma,\theta}1,
\]
and
\[
 \left|\sum_q w_N^\nu(q)
  \left\{\pi_{I_N}(q,a)-\frac{\pi_{I_N}}{\varphi(q)}\right\}\right|
 \ll_{\sigma,A,\theta,B,\varepsilon,C_{\rm box}}N(\log N)^{-A}.
\]
The sum covers the complete support of formula~\textup{(W)}.
For the admissible choice
$\theta=19101/32000-10^{-4}$, the sieve ratio
$\log D/\log z=500\theta$ equals $298.403125$.
The implied constant and the threshold for $N$ are ineffective.
\end{theorem}

\medskip
\noindent\textit{Remark (parameter conventions).}
In Theorem~\ref{thm:aa37-fixed-u} one takes $A_{\rm agg}=A$ and
$C_{\rm box}\ge\lceil2A+32\rceil$.  The shifted choice
$A_{\rm agg}=A+3$ and the corresponding
$C_{\rm box}\ge\lceil2A+38\rceil$ occur only in the auxiliary
source-wide estimate (3.95), which is not used in the proofs of
Theorems~\ref{thm:aa37-fixed-u}--\ref{thm:aa37-interior} or
Corollary~\ref{cor:aa37-interior-li}, nor in Section~\ref{sec:linear-sieve-application}.

\begin{theorem}[Signed shifted-prime mean value]
\label{thm:aa37-interior}
Fix $A>0$, $B>0$, $0<\varepsilon\le10^{-4}$, $0<u\le1/500$,
and
\[
 \frac12<\theta<\Theta(u):=\frac{5-7/32-3u}{8}.
\]
Put
\[
 \begin{gathered}
 \eta_\Theta=\Theta(u)-\theta,\\
 0<\sigma\le\min\{10^{-10},\eta_\Theta/200,
                   \sqrt{u/(2\theta)}\},\\
 \tau=\sigma^9,\quad D=N^\theta,\quad z=N^u,
 \quad Y=D^{\sigma^2}.
 \end{gathered}
\]
Set $K_{\rm scale}=B+1$. Choose an integer box-width exponent
\[
 C_{\rm box}\ge\lceil2A+32\rceil.
\]
The implied constant and the threshold for $N$ below may also depend
on the chosen $C_{\rm box}$. For the canonical choice
$C_{\rm box}=\lceil2A+32\rceil$ this is a dependence on $A$ alone.
Equivalently $C_{\rm box}\ge2(A+c_R)$, with the boundary exponent
$c_R=16$ of (3.24) and (B.19). This also gives the logarithm-freezing
condition $C_{\rm box}\ge A+18$ from (B.18).
Let $N$ be sufficiently large and even. For each sign
$\nu\in\{+,-\}$, let $w_N^\nu(q)$ be the complete Iwaniec
remainder coefficient defined by formula~\textup{(W)} in
Section~\ref{sec:linear-sieve-application}. Its prime set is
$\mathcal P_N=\{p\text{ prime}:p\nmid2N\}$, with the
grid and distinctness conditions specified there. Let $a$ be any
integer satisfying
\[
 N\le a\le N\log N/3,\qquad
 (a,q)=1\quad\text{for every }q\in\operatorname{supp}w_N^\nu.
\]
Call these shifts admissible; they form a subset of the interval that includes
$a=N$.
Put
\[
 I_N=\left(\frac{N}{(\log N)^B},(1-\varepsilon)N\right).
\]
Write $\pi_{I_N}$ for the number of primes in $I_N$.
Write $\pi_{I_N}(q,a)$ for the number of these primes with
$p\equiv a\pmod q$.
For every fixed $\sigma$ in the displayed range, these coefficients
satisfy
\[
 \begin{gathered}
 \operatorname{supp}w_N^\nu\subseteq
 \{q<D^{1+\sigma+\tau}:(q,2N)=1\},
 \\ D^{1+\sigma+\tau}<N^{\Theta(u)},
 \qquad |w_N^\nu(q)|\ll_{\sigma,u,\theta}1,
 \end{gathered}
\]
and
\begin{equation}\label{eq:aa37-interior-mean}
 \left|\sum_q w_N^\nu(q)
  \left\{\pi_{I_N}(q,a)-\frac{\pi_{I_N}}{\varphi(q)}\right\}\right|
 \ll_{\sigma,A,u,\theta,B,\varepsilon,C_{\rm box}} N(\log N)^{-A}.
\end{equation}
The sum includes the entire support of formula~\textup{(W)}.
The sieve ratio is $\log D/\log z=\theta/u$. The condition on
$\sigma$ gives $Y=D^{\sigma^2}=N^{\theta\sigma^2}\le N^{u/2}<z$,
as required for the rough-factor bounds.
The implied constant and the threshold for $N$ are ineffective.
The proof uses the Siegel--Walfisz theorem.
\end{theorem}

\begin{proof}[Proof of Theorem~\ref{thm:aa37-fixed-u}]
Apply Theorem~\ref{thm:aa37-interior} with $u=1/500$.
Then $\Theta(u)=19101/32000$, and its condition on $\sigma$ becomes
the displayed condition above. All the support, coefficient,
and discrepancy bounds specialize as stated. The numerical sieve
ratio follows from $\log D/\log z=\theta/u$.
At $\theta=19101/32000-10^{-4}$ the proof uses Theorem~\ref{thm:main}
through Theorem~\ref{thm:aa37-interior}. That theorem allows support
up to $X^{0.6}$.
It does not use the introductory proposition. At $\theta=\vartheta_0$
the support bound $D^{1+\sigma+\tau}$ of \textup{(W)} exceeds $X^{\vartheta_0}$.
\end{proof}

In Theorem~\ref{thm:aa37-interior}, $D=N^\theta$ and
$D^{1+\sigma+\tau}<N^{\Theta(u)}$ by the choice of $\sigma$.
The introductory source-box estimate has the fixed local level
$\vartheta_0=19101/32000-10^{-4}$. It assumes its support hypothesis. The packing ceiling for Theorem~\ref{thm:aa37-interior}
is $\Theta(u)$. The latter applies the $X^{0.6}$
dispersion theorem~\ref{thm:main} together with the source estimates
below. Here $\Theta(u)<153/256<0.6$. Fix $B>0$. Each source shell has
$X>N/(2(\log N)^B)$. So, for sufficiently large $N$,
\[
 N^{\Theta(u)}\le\left(\frac{N}{2(\log N)^B}\right)^{0.6}<X^{0.6}.
\]
Thus the full support of Iwaniec's remainder weights lies in the range of
Theorem~\ref{thm:main}. Section~\ref{sec:pack} proves the modulus
packing. Appendix~\ref{app:proof-aa37} estimates the Heath--Brown,
boundary, and other source terms. Section~\ref{sec:linear-sieve-application}
assembles these estimates.

\begin{corollary}[Logarithmic-integral comparator]
\label{cor:aa37-interior-li}
Under the hypotheses of Theorem~\ref{thm:aa37-interior}, for each
$\nu\in\{+,-\}$,
\[
 \begin{aligned}
 &\left|\sum_q w_N^\nu(q)
  \left\{\pi_{I_N}(q,a)-
  \frac{\substack{\Li((1-\varepsilon)N)\\
                   {}-\Li(N/(\log N)^B)}}
       {\varphi(q)}\right\}\right|\\
 &\ll_{\sigma,A,u,\theta,B,\varepsilon,C_{\rm box}} N(\log N)^{-A}.
 \end{aligned}
\]
The implied constant and the threshold for $N$ are ineffective.
\end{corollary}

The proof uses the spectral large sieves of Deshouillers and Iwaniec
\cite{DI82}. It uses the Kim--Sarnak exponent $7/64$ \cite{KS03}.
It also uses Deligne's bound for holomorphic coefficients
\cite{Deligne74}. Section~\ref{p1:sec:classical-inputs} states these
inputs. The dispersion argument and its application are in
Sections~\ref{sec:pi-psdt}--\ref{sec:linear-sieve-application}.

\section{Classical inputs and Theorem 2.8}\label{p1:sec:classical-inputs}

Put $e(t)=\exp(2\pi i t)$. Write $n\sim N$ for $N<n\leq2N$. Unspecified arithmetic indices are positive integers. Norms of sequences are counting-measure $\ell^2$ norms. A star denotes reduced residue classes in residue sums. In character sums it denotes primitive characters. These two uses are distinguished explicitly below. Constants are uniform within each sequence family. Put
\[
 \theta_{\mathrm{KS}}=\frac7{32}.
\tag{2.1}\label{p1:eq:2.1}
\]
Thus $\theta_{\mathrm{KS}}=2\cdot7/64$, where $7/64$ is the unramified Ramanujan exponent. This symbol is distinct from the level parameter $\theta$ and the exceptional parameters $\vartheta_j$.

\subsection{Spectral normalization and large sieves}

The spectral large sieves are due to Deshouillers and Iwaniec \cite{DI82}.

Let $q\geq1$. Use Petersson-orthonormal bases at level $\Gamma_0(q)$. The measure is $dx\,dy/y^2$. Take $\sigma_\infty=I$. For a Maass cusp form,
\[
 u_j(z)=\sum_{n\ne0}\rho_{j\infty}(n)\sqrt y\,
 K_{i\kappa_j}(2\pi|n|y)e(nx),\qquad
 \lambda_j=\tfrac14+\kappa_j^2.
\tag{2.2}\label{p1:eq:2.2}
\]
This is the Fourier expansion of \cite[(1.15)]{DI82}, as printed there.
It has no factor $2\sqrt{|n|}$; that factor enters only through the dictionary (2.11).
For a holomorphic form $f$ of even weight $k\geq2$, write
\[
 f(z)=\sum_{n\geq1}\psi_f(\infty,n)e(nz).
\tag{2.3}\label{p1:eq:2.3}
\]
The Eisenstein coefficients $\varphi_{\mathfrak c\infty n}(s)$ have the normalization of \cite[(1.17)]{DI82}. Their relation to the Whittaker coefficients $\rho^{\rm Dr}$ of Drappeau, whose full dictionary is (2.11) in \S2.2 below, is
\[
 \rho^{\rm Dr}_{\mathfrak c\infty}(n,t)
 =\frac{\pi^{1/2+it}|n|^{it-1/2}}{\Gamma(1/2+it)}
   \varphi_{\mathfrak c\infty n}(1/2+it).
\tag{2.4}\label{p1:eq:2.4}
\]
Let $\mathfrak c$ run over inequivalent cusps. Use width-one scaling matrices. The continuous-spectrum measure in the DI large sieve is plain $dt$.

\textbf{Theorem 2.1 (consequence of \cite[Theorem~2]{DI82}).} Let $\epsilon>0$, $q\geq1$, $K\geq1$ and $N\geq\tfrac12$. Let $a_n\in\mathbb C$ be supported on $N<n\leq2N$. Then each of
\[
\begin{split}
 H(K,a)&=\sum_{\substack{2\leq k\leq K\\k\in{\rm even}}}
 \frac{(k-1)!}{(4\pi)^{k-1}}
 \sum_{f\in\mathcal B_k(q)}
 \left|\sum_{n\sim N}a_n n^{-(k-1)/2}\psi_f(\infty,n)\right|^2,\\
 M(K,a)&=\sum_{|\kappa_j|\leq K}
 \frac1{\cosh(\pi\kappa_j)}
 \left|\sum_{n\sim N}a_n\rho_{j\infty}(n)\right|^2,\\
 E(K,a)&=\sum_{\mathfrak c}\int_{-K}^{K}
 \left|\sum_{n\sim N}a_n n^{it}
 \varphi_{\mathfrak c\infty n}(1/2+it)\right|^2dt
\end{split}
\tag{2.5}\label{p1:eq:2.5}
\]
is at most
\[
 C_\epsilon\left(K^2+\frac{N^{1+\epsilon}}q\right)
 \sum_{n\sim N}|a_n|^2,
\tag{2.6}\label{p1:eq:2.6}
\]
uniformly in $q,K,N,a$. By \cite[Theorem~2, (1.28)--(1.30)]{DI82},
each expression is bounded by
$C_\epsilon\{K^2+\mu(\mathfrak a)N^{1+\epsilon}\}\|a_N\|_2^2$,
with the constant depending only on $\epsilon$. At
$\mathfrak a=\infty$, $\mu(\infty)=1/q$, giving (2.6); the source includes
the endpoint $N\geq\tfrac12$. The third line of (2.5) has no factor
$1/\cosh(\pi t)$. The cusp transport for Theorem 2.8 is in
\S\ref{p1:sec:t-cusp-transport}.

For an exceptional eigenvalue $0<\lambda_j<\tfrac14$, choose $i\kappa_j>0$, and put
\[
 \vartheta_j=2i\kappa_j=2\sqrt{\tfrac14-\lambda_j}.
\tag{2.7}\label{p1:eq:2.7}
\]

\textbf{Theorem 2.2 (\cite[Theorem~5]{DI82}, the specialization used here).} Let $\epsilon>0$, $q\geq1$, $N\geq\tfrac12$ and $Y\geq1$. Let $a_n\in\mathbb C$ be supported on $n\sim N$. Then
\[
 \sum_{j\ {\rm exceptional}}Y^{\vartheta_j}
 \left|\sum_{n\sim N}a_n\rho_{j\infty}(n)\right|^2
 \ll_\epsilon
 \left(1+\sqrt{\frac{NY}q}\right)
 \left(1+\sqrt{\frac{N^{1+\epsilon}}q}\right)\|a\|_2^2.
\tag{2.8}\label{p1:eq:2.8}
\]
This is \cite[Theorem~5]{DI82}. It uses the notation of Theorem 2, so $N\geq\tfrac12$ is allowed. The implied constant depends on $\epsilon$ alone.

\textbf{Theorem 2.3 (from \cite[proof of Theorem~7, (8.17)]{DI82}).} Let $\epsilon>0$. Let $a_n$ be the characteristic sequence of an integer interval
$I\subset(N,2N]$. Then, for some $c_\epsilon$ and uniformly in $Q,N,Y\ge1$ and $I$,
\[
\sum_{2\le q\le Q}\ \sum_{j\ {\rm exceptional}}^{(q)}Y^{\vartheta_j}
\Bigl|\sum_{n\sim N}a_n\rho_{j\infty}(n)\Bigr|^2\le c_\epsilon(QN)^\epsilon\bigl(Q+N+\sqrt{NY}\bigr)N.
\tag{2.9}\label{p1:eq:2.9}
\]
The range starts at $q=2$ because Theorem 2.8 uses level $q=rs$ with $r>R\ge1$, so $q\ge2$ there.

\emph{Proof.} By \cite[(8.17) and (8.4)]{DI82},
\[
\begin{aligned}
&\sum_{Q_0<q\le16Q_0}\sum_{\rm exc}Y^{\vartheta_j}
 \left|\sum_{N<n\le2N}a_n\rho_{j\infty}(n)\right|^2\\
&\qquad\le c^{\rm DI}_{\epsilon'}(Q_0N)^{5\epsilon'}
                  (Q_0+N+\sqrt{NY})N
\end{aligned}
\]
for $Q_0\ge1$. In \cite{DI82} the interval is printed as $[N,N_1]$, $N\le N_1\le2N$. The sum runs over $N<n\le2N$, so $a_n$ is the characteristic sequence of the initial interval
$(N,N_1]\cap\mathbb Z$. The source estimate (8.17) is uniform for
$Q_0,N,Y\ge1$ and every real $N_1\in[N,2N]$; its final statement imposes
no relation between $Y$ and $Q_0,N$. In the proof, the induction first treats
$1\le N\le Q_0$ at $Y_*=Q_0^{2-2\epsilon'}/N$, while (8.18) handles the
remaining large-$N$ range, and the final comparison extends the result to
every $Y\ge1$. We write all intervals in this half-open form.

(i) \emph{$q$-blocks.} Let $I_Q=\max\{i\ge0:16^{-i}Q\ge1\}$. The blocks $(16^{-i}Q,16^{1-i}Q]$,
$1\le i\le I_Q$, together with $(1,16]$, cover $2\le q\le Q$. There are at most $2+\log_{16}Q$ of
them. Every summand is nonnegative. So each block may be enlarged to its full DI range. Each block has
$Q_0\le Q$.

(ii) \emph{Intervals.} Write $I=(u,v]\cap\mathbb Z$ with $N\le u\le v\le2N$. Then
$\mathbf1_I=\mathbf1_{(N,v]}-\mathbf1_{(N,u]}$, both initial intervals in DI's sense, and
$|A-B|^2\le2|A|^2+2|B|^2$. This costs a factor 4.

Hence the left side of (2.9) is at most
$4c^{\rm DI}_{\epsilon'}(2+\log_{16}Q)(QN)^{5\epsilon'}(Q+N+\sqrt{NY})N$. Take $\epsilon'=\epsilon/10$
and $c_\epsilon:=4c^{\rm DI}_{\epsilon/10}\sup_{Q\ge1}(2+\log_{16}Q)Q^{-\epsilon/2}<\infty$. \hfill$\square$

The interval hypothesis is the one in \cite[(8.17)]{DI82}, and the induction closes in the proof of Theorem~7. Theorem 2.3 comes from this step of that proof. \cite[Theorem~7]{DI82} treats $a_n=1$ on $n\le N$. Theorem 2.3 needs $N\ge1$. The one-integer block $N=\tfrac12$ is
treated in \S\ref{p1:sec:t-exceptional}.

\subsection{Hecke exponents and the adopted trace normalization}

\textbf{Kim--Sarnak input.} The bound $7/64$ at an unramified place is \cite[Appendix~2, Proposition~2]{KS03}; the Satake normalization is stated in \cite[p.~175]{KS03}. Independently, the archimedean input
$\lambda_j\geq975/4096$ is the lower bound in \cite[Appendix~2]{KS03}; it is not a consequence of the finite-place estimate below. Let $f$ be a weight-zero Maass newform of congruence level. Each exceptional parameter of $f$ satisfies $0\leq i\kappa\leq7/64$. For a prime $p$ not dividing the level,
% Locator retained: the Satake normalization passage on printed p. 175 has no theorem or equation number.
\[
 |\alpha_p|,|\beta_p|\leq p^{7/64},\qquad
 |\lambda_f(p^\nu)|\leq(\nu+1)p^{7\nu/64}
 \quad(\nu\geq0).
\tag{2.10}\label{p1:eq:2.10}
\]
Thus the preceding archimedean input gives $\vartheta_j\leq\theta_{\mathrm{KS}}$, equivalently $\lambda_j\geq975/4096$ by (2.7). At the finite places, (2.10) gives $|\lambda_f(k)|\leq\tau(k)k^{\theta_{\mathrm{KS}}/2}$ for $k$ coprime to the level. We use these bounds for newforms of level dividing $q$. These newforms diagonalize the unramified Hecke operators on the full level-$q$ space. For holomorphic newforms we use Deligne's $ |\lambda_f(k)|\leq\tau(k)$, again for $k$ coprime to the level. Both factorizations use primes coprime to the level.

\textbf{Deligne input.} The holomorphic bound is \cite[Th\'eor\`eme~8.2]{Deligne74}. Both inputs enter at primes coprime to the level, through the prime-power recurrences and multiplicativity above.

Bruggeman and Kuznetsov independently introduced the relevant trace
formulas; see \cite[p.~220]{DI82} and their original papers
\cite{Bruggeman78,Kuznetsov81}. For the normalization and formulas used
here, see \cite[(4.7), (4.9)--(4.19), Lemma~4.5]{Drappeau}, in the PLMS version.
% Locator retained: DI82's historical attribution on printed p. 220 has no theorem or equation number.
The coefficient dictionary is
\[
\begin{aligned}
 \rho^{\rm Dr}_{f\mathfrak a}(n)&=\frac{\rho^{\rm DI}_{f\mathfrak a}(n)}{2\sqrt{|n|}}
 &&\hbox{(Maass)},\\
 \rho^{\rm Dr}_{f\mathfrak a}(n)&=(4\pi n)^{-k/2}\psi_f(\mathfrak a,n)
 &&\hbox{(holomorphic)},
\end{aligned}
\tag{2.11}\label{p1:eq:2.11}
\]
with (2.4) for Eisenstein series. The Maass identity follows from $W_{0,it}(4\pi|n|y)=2\sqrt{|n|y}K_{it}(2\pi|n|y)$.

For $\phi\in C_c^\infty((0,\infty))$, define the DI transforms
\[
\begin{split}
 \widehat\phi(t)&=\frac\pi{\sinh\pi t}
 \int_0^\infty\frac{J_{2it}(x)-J_{-2it}(x)}{2i}\phi(x)\frac{dx}x,\\
 \check\phi(t)&=\frac4\pi\cosh\pi t
 \int_0^\infty K_{2it}(x)\phi(x)\frac{dx}x,\\
 \widetilde\phi(k-1)&=\int_0^\infty J_{k-1}(x)\phi(x)\frac{dx}x.
\end{split}
\tag{2.12}\label{p1:eq:2.12}
\]
Values at removable singularities are limits. In Drappeau's notation the corresponding transforms are $-4\widehat\phi$, $2\pi\check\phi$, and $4i^k\widetilde\phi(k-1)$.

Theorem 2.8 and Section~3 use the following identity. Let $q=rs$, $(r,s)=1$, and $r>1$; take
\[
 \sigma_{1/s}=\begin{pmatrix}\sqrt r&0\\s\sqrt r&1/\sqrt r\end{pmatrix},
 \qquad
 \mathcal C(\infty,1/s)=\{\gamma=cs\sqrt r:c\geq1, (c,r)=1\}.
\tag{2.13}\label{p1:eq:2.13}
\]
The scaling matrix and admissible moduli in (2.13) are also given in
\cite[(1.6) and (2.4)]{DI82} and \cite[Lemma~4.3]{Drappeau}.
Since $r>1$, the cusps $\infty$ and $1/s$ are inequivalent. Hence the same-sign identity has no diagonal term. If a matrix in $\Gamma_0(rs)$ maps $\infty$ to the reduced fraction $1/s$, its lower-left entry is $\pm s$. Then $rs\mid s$, so $r=1$. Every application here has $r>R\geq1$. For $m,n\geq1$, put $\mathcal K_\pm(m,n)=\sum_{\gamma\in\mathcal C}\gamma^{-1}S_{\infty,1/s}(m,\pm n;\gamma)\phi(4\pi\sqrt{mn}/\gamma)$. In the DI coefficients, Drappeau's formulas read
\[
\begin{split}
 \mathcal K_+(m,n)
 &=\frac1\pi\sum_{\substack{k\geq2\\k\ {\rm even}}}
 i^k(k-1)!(4\pi\sqrt{mn})^{1-k}\widetilde\phi(k-1)
 \sum_f\overline{\psi_f(\infty,m)}\psi_f(1/s,n)\\
 &\quad-\sum_j\frac{\widehat\phi(\kappa_j)}{\cosh\pi\kappa_j}
 \overline{\rho_{j\infty}(m)}\rho_{j,1/s}(n)\\
 &\quad-\frac1\pi\sum_{\mathfrak c}\int_{\mathbb R}
 \widehat\phi(t)(n/m)^{it}
 \overline{\varphi_{\mathfrak c\infty m}(1/2+it)}
 \varphi_{\mathfrak c,1/s,n}(1/2+it)\,dt,\\[2mm]
 \mathcal K_-(m,n)
 &=\frac\pi2\sum_j\frac{\check\phi(\kappa_j)}{\cosh\pi\kappa_j}
 \overline{\rho_{j\infty}(m)}\rho_{j,1/s}(-n)\\
 &\quad+\frac12\sum_{\mathfrak c}\int_{\mathbb R}
 \check\phi(t)(n/m)^{it}
 \overline{\varphi_{\mathfrak c\infty m}(1/2+it)}
 \varphi_{\mathfrak c,1/s,-n}(1/2+it)\,dt.
\end{split}
\tag{2.14}\label{p1:eq:2.14}
\]
The proof uses only the identities (2.7.T16)--(2.7.T18), derived in \S\ref{p1:sec:t-cusp-pair}. The continuous integrals carry plain $dt$. Let $\alpha_m,\beta_n$ be finite sequences. Multiply (2.14) by $\alpha_m\beta_n$, sum in $(m,n)$, then take spectral absolute values and apply Cauchy--Schwarz. The first functional is $L_1(u)=\sum_m\overline{\alpha_m}\rho_{u\infty}(m)$. The spectral product is $\overline{L_1(u)}\sum_n\beta_n\rho_{u,1/s}(\pm n)$.

\subsection{BFI and the classical progression estimates}

For a finitely supported sequence $f$, define
\[
 \Delta_f(q,a)=\sum_{n\equiv a\ (q)}f(n)
 -\frac1{\varphi(q)}\sum_{(n,q)=1}f(n),\qquad (a,q)=1.
\tag{2.15}\label{p1:eq:2.15}
\]
The notation of \cite[(1.1)]{BFI86} writes $\Delta_f(x;q,a)$ for (2.15) with both sums cut at $n\le x$.

The progression estimates are from Bombieri, Friedlander and Iwaniec
\cite{BFI86}. We use \cite[(A2), Theorem~0(a), and proof of Theorem~0(a)]{BFI86}.

\textbf{BFI hypothesis (A2) \cite[(A2)]{BFI86}.} Fix $B_\beta>0$ for the sequence family. Let $\beta$ be in the family, with support parameter $N\geq2$. For every $d,k\geq1$ and integer $\ell\ne0$ with $(\ell,k)=1$, and for every $A>0$, there is $C_A$ with
\[
 \left|\sum_{\substack{n\equiv\ell\ (k)\\(n,d)=1}}\beta_n
 -\frac1{\varphi(k)}\sum_{(n,dk)=1}\beta_n\right|
 \leq C_A\|\beta\|_2N^{1/2}\tau(d)^{B_\beta}(\log(2N))^{-A}.
\tag{2.16}\label{p1:eq:2.16}
\]
The $C_A$ are independent of $N,d,k,\ell$ and of $\beta$ in the family. As in \cite{BFI86}, $d$ and $k$ run over all positive integers. If $\|\beta\|_2^2\geq N(\log 2N)^{-B'}$ for a fixed $B'$, then a bound $C_AN\tau(d)^{B_\beta}(\log 2N)^{-(A+B'/2)}$ implies (2.16), after adjusting the constant.

\textbf{Theorem 2.4 (\cite[Theorem~0(a), (2.2)]{BFI86}).} Let $\beta_n$ be supported on $n\leq N$, $N\geq2$, and satisfy (2.16). For every $A>0$, there is $B_1>0$ such that, uniformly for $Q\leq N(\log N)^{-B_1}$,
\[
 \sum_{q\leq Q}\ \sum_{\substack{a\bmod q\\(a,q)=1}}
 |\Delta_\beta(q,a)|^2
 \ll \|\beta\|_2^2N(\log N)^{-A}.
\tag{2.17}\label{p1:eq:2.17}
\]
The constants depend on $A,B_\beta$ and the constants in (2.16). For a sequence on $N<n\leq2N$, replace $N$ by $2N$ in (2.17).

\cite[Theorem~0(b)]{BFI86} is not used; its discrepancy has the cutoff $n\leq x$ of \cite[(1.1)]{BFI86}, while (2.15) is a full-support discrepancy. The switched bilinear form uses the full-support bound (3.49), proved from (3.25), (3.48) and Theorem 2.4.

That bound uses \cite[(1.6)]{BFI86}. Let $c_h\in\mathbb C$ be supported on $1\leq h\leq H$. The large sieve gives
\[
 \sum_{q\leq Q}
 \sum_{\chi\bmod q}^{\rm primitive}
 \left|\sum_hc_h\chi(h)\right|^2
 \ll(Q^2+H)\|c\|_2^2.
\tag{2.19}\label{p1:eq:2.19}
\]
We use this initial-interval form with $H$ the upper support bound.

\textbf{Bombieri--Vinogradov input.} Write $\psi(y;q,a)=\sum_{n\leq y,n\equiv a\ (q)}\Lambda(n)$. For every $A>0$ there is $B=B(A)>0$ such that, for $y\geq3$ and $Q\leq y^{1/2}(\log y)^{-B}$,
\[
 \sum_{q\leq Q}\max_{(a,q)=1}
 \left|\psi(y;q,a)-\frac y{\varphi(q)}\right|
 \ll_A y(\log y)^{-A}.
\tag{2.20}\label{p1:eq:2.20}
\]
This is BFI's formulation of the Bombieri--Vinogradov theorem
\cite[(1.4)]{BFI86}; compare the original results
\cite{Bombieri65,Vinogradov65}. The maximum allows the varying reduced
residue $a=N\bmod q$. We apply (2.20) at each endpoint separately, with
its own level inequality. Thus (2.20) is used at fixed $y$ only.

\textbf{Siegel--Walfisz input.} We use the power-log form of
\cite[Corollary~5.29]{IK04}. Remove prime powers and sum by parts.
Let $C>0$ and $A>0$ be fixed. For $y\geq3$, $q\leq(\log y)^C$
and $(a,q)=1$,
\[
 \pi(y;q,a)=\frac{\operatorname{Li}(y)}{\varphi(q)}
 +O_{A,C}\bigl(y(\log y)^{-A}\bigr).
\tag{2.21}\label{p1:eq:2.21}
\]
The constant is uniform in reduced $a$ and is ineffective. Subtraction at two comparable endpoints keeps any fixed logarithmic saving. We use (2.21) with this power-log error. Three further forms need their own proofs. The first saves a power of the interval length. The second keeps the tensor and coprimality structure of the structured sequences. The third is localized to one dyadic box of the Heath--Brown decomposition. Section~3 proves the short-interval estimates and those for its structured sequences.

\subsection{The linear sieve and its structured remainder}

Let $\mathcal A$ be a finite integer sequence with nonnegative multiplicities. Let $\mathcal P$ be a prime set. Put
\[
 P(z)=\prod_{\substack{p<z\\p\in\mathcal P}}p,\quad
 S(\mathcal A,\mathcal P,z)=\sum_{a\in\mathcal A}1_{(a,P(z))=1},
 \quad |\mathcal A_d|=\frac{\omega(d)}dX+r(\mathcal A,d).
\tag{2.22}\label{p1:eq:2.22}
\]
Here $X>0$ and $\omega$ is multiplicative. Require $0<\omega(p)<p$ for $p\in\mathcal P$. Set $\omega(p)=0$ for excluded primes. Set $V(z)=\prod_{p<z,p\in\mathcal P}(1-\omega(p)/p)$. Use the continuous linear-sieve solutions
\[
\begin{gathered}
 sF(s)=2e^\gamma\quad(0<s\leq3),\qquad
 sf(s)=0\quad(0<s\leq2),\\
 (sF(s))'=f(s-1)\quad(s>3),\qquad
 (sf(s))'=F(s-1)\quad(s>2).
\end{gathered}
\tag{2.23}\label{p1:eq:2.23}
\]
On $s\geq1$, $F$ decreases to 1. Also $f$ increases to 1, with $f=0$ through $s=2$. In particular $F(1)=2e^\gamma$. We use this endpoint when the level is at least the sifting limit.

Let $K_{\mathrm M}\geq2$ be one fixed absolute constant such that, for all
$z>w\geq2$,
\[
 \prod_{w\leq p<z}\left(1-\frac1p\right)^{-1}
 \leq\frac{\log z}{\log w}\left(1+\frac{K_{\mathrm M}}{\log w}\right).
\]
Such a constant exists by Mertens' product estimate, after increasing it
to cover bounded $w$.

\textbf{Theorem 2.5 (\cite[Theorem~1]{AA36}, dimension one).} Suppose, for some fixed $K\geq2$ and every $z>w\geq2$,
\[
 \prod_{\substack{w\leq p<z\\p\in\mathcal P}}
 (1-\omega(p)/p)^{-1}
 \leq\frac{\log z}{\log w}\left(1+\frac K{\log w}\right).
\tag{2.24}\label{p1:eq:2.24}
\]
Let $y\geq z\geq2$ and $s=\log y/\log z$. Then
\[
\begin{split}
 S(\mathcal A,\mathcal P,z)
 &\leq XV(z)\{F(s)+e^{\sqrt K}Q_1(s)(\log y)^{-1/3}\}
 +\sum_{\substack{d<y\\d\mid P(z)}}|r(\mathcal A,d)|,\\
 S(\mathcal A,\mathcal P,z)
 &\geq XV(z)\{f(s)-e^{\sqrt K}Q_1(s)(\log y)^{-1/3}\}
 -\sum_{\substack{d<y\\d\mid P(z)}}|r(\mathcal A,d)|,
\end{split}
\tag{2.25}\label{p1:eq:2.25}
\]
where, for $s\geq1$,
\[
 Q_1(s)\leq\exp\{-s\log s+s\log\log(3s)+O(s)\}.
\tag{2.26}\label{p1:eq:2.26}
\]
See \cite{AA36}. The hypotheses are (1.1)--(1.3), and Theorem 1 is (1.4)--(1.6). The dimension-one constants $\beta=2$ and $A=2e^\gamma$ appear in \cite[p.~176]{AA36}; the source states $B=0$ in \cite[p.~173]{AA36}. For $s\leq2$ the lower main function in (2.25) vanishes. Hence its lower main term is nonpositive.
% Locators retained: the dimension-one table on p. 176 and the B=0 statement on p. 173 are unnumbered.

For the later prime set $\mathcal P_N=\{p:p\nmid2N\}$, take the
multiplicative function with $\omega(p)=1$ on $\mathcal P_N$,
$\omega(p)=0$ off it, and $\omega(p^a)=0$ for $a\geq2$. Its product
in (2.24) is over a subset of all primes in $[w,z)$, so (2.24) holds
uniformly in $N$ with the fixed value $K=K_{\mathrm M}$. The same data
satisfy (2.27) with $L_\omega=2$. Thus, whenever (2.25) is specialized
to this set, its factor is $e^{\sqrt{K_{\mathrm M}}}$. This is fixed;
combining it with (2.26) changes only the absolute constant in the
$O(s)$ term there.
This $\omega$ is a nominal sieve density: by Theorem~2.6 the constructed
coefficients are independent of $\mathcal A$, $X$, and $\omega$.

The structured remainder is from \cite[Theorem~4]{AA37}.

Assume (2.24) and nonnegativity of $\omega$ on prime powers. Fix $L_\omega>1$. Assume, for all $z>w\geq2$,
\[
 \sum_{\substack{w\leq p<z\\p\in\mathcal P}}
       \sum_{a\geq2}\frac{\omega(p^a)}{p^a}
 \leq\frac{L_\omega}{\log(3w)}.
\tag{2.27}\label{p1:eq:2.27}
\]
These are \cite[(1)--(2)]{AA37} in dimension one, with $L=L_\omega$.

Fix $0<\epsilon_s<1/3$ and $D\geq2$. Put $\delta_s=\epsilon_s^9$ and
\[
 \mathcal G(D)=\{D^{\epsilon_s^2(1+\delta_s)^n}:n\geq0\}.
\tag{2.28}\label{p1:eq:2.28}
\]
Let $\mathcal H(D)$ consist of tuples $\mathbf D=(D_1,\ldots,D_\rho)$ from this grid with $D_\rho\leq\cdots\leq D_1<D^{1/2}$. Define
\[
\begin{split}
 \mathcal D^+(D)&=\{\mathbf D\in\mathcal H(D):
 D_1\cdots D_{2l}D_{2l+1}^3<D
 \quad(0\leq l\leq(\rho-1)/2)\},\\
 \mathcal D^-(D)&=\{\mathbf D\in\mathcal H(D):
 D_1\cdots D_{2l-1}D_{2l}^3<D
 \quad(1\leq l\leq\rho/2)\}.
\end{split}
\tag{2.29}\label{p1:eq:2.29}
\]
\cite[Theorem~4]{AA37} gives $|\mathcal D^+|,|\mathcal D^-|<\exp(8\epsilon_s^{-3})$. This bound follows the definitions of $\mathcal D^+$ and $\mathcal D^-$. The grid is built from the relative level of each application.

\textbf{Theorem 2.6 (from \cite[Theorem~4]{AA37}, with $R=1$).} Assume (2.22), (2.24), (2.27), with the parameters above. Then there are coefficients $\phi_d^\pm$ and $\lambda_d^\pm(\mathbf D)$ of modulus at most 1. They are independent of $\mathcal A,X,\omega$. They satisfy
\[
\begin{aligned}
 S(\mathcal A,\mathcal P,z)&\leq XV(z)(F(s)+E_s)+R_1^++R^+
 &&(2\leq z\leq D),\\
 S(\mathcal A,\mathcal P,z)&\geq XV(z)(f(s)-E_s)-R_1^--R^-
 &&(2\leq z\leq D^{1/2}),
\end{aligned}
\tag{2.30}\label{p1:eq:2.30}
\]
where $s=\log D/\log z$ and
\[
 E_s\ll\epsilon_s+\epsilon_s^{-8}e^{K+L_\omega}(\log D)^{-1/3},
\tag{2.31}\label{p1:eq:2.31}
\]
with an absolute implied constant. \cite[Theorem~4]{AA37} contains $e^{K+L}$. For the coefficient construction used here, the prime set and multiplicative function just specified give $K=K_{\mathrm M}$ and $L_\omega=2$. Hence (2.31) has the fixed factor $e^{K_{\mathrm M}+2}$; it is retained in the formula and may be absorbed into the implied constant. The prime-power dependence comes from its (2), namely (2.27). In its estimate (31), put $u=D^{\epsilon_s^2}$. Then
\[
 U=\sum_{\substack{u\leq p<z\\p\in\mathcal P}}
       \sum_{a\geq2}\frac{\omega(p^a)}{p^a}\leq L_\omega,
 \qquad
 \prod_{\substack{u\leq p<z\\p\in\mathcal P}}
 \left\{\left(1-\frac{\omega(p)}p\right)
 \sum_{a\geq0}\frac{\omega(p^a)}{p^a}\right\}\leq e^U.
\]
For the first bound use $\max\{2,u\}$ in (2.27). An empty prime interval contributes zero. The second bound follows factor by factor. Thus $e^U\leq e^{L_\omega}$ bounds the extra prime-power factors in the nonsquarefree terms. Formula (2.31) keeps the factor $e^{K+L_\omega}$ of the source. The signed remainders are
\[
 R_1^\nu=\sum_{\substack{d<D^{\epsilon_s}\\d\mid P(D^{\epsilon_s^2})}}
 \phi_d^\nu r(\mathcal A,d),
\tag{2.32}\label{p1:eq:2.32}
\]
\[
 R^\nu=\sum_{\mathbf D\in\mathcal D^\nu(D)}
 \sum_{\substack{d<D^{\epsilon_s}\\d\mid P(D^{\epsilon_s^2})}}
 \lambda_d^\nu(\mathbf D)
 \sum_{\substack{D_i\leq p_i<D_i^{1+\delta_s}\\p_i\mid P(z)\\
                         p_i\ne p_j\ ((i,j)\in I(\mathbf D))}}
 r(\mathcal A,dp_1\cdots p_\rho).
\tag{2.33}\label{p1:eq:2.33}
\]
The sets $I(\mathbf D)$ impose distinctness on equal-label tail positions. Unequal grid intervals are disjoint. Coefficients may be zero on unused tuples. They depend only on $\mathcal P,D,z,\epsilon_s,\mathbf D$ and the distinctness conditions. In particular they are independent of $\mathcal A,X,\omega$. They may change with the level or the cutoff.

\cite[Theorem~4]{AA37} holds for general $R\geq1$, $0<\epsilon_s<3^{-R}$, with the respective ranges in (2.30). Its preceding combinatorial identities give $S\leq W_R^+$ and $S\geq W_R^-$. The coefficients come from \cite[(23)--(25), (26), and Lemma~4]{AA37}; these formulas also show their independence of the sequence.

\textbf{Relative-level grid.} Choose $0<\sigma\leq10^{-10}$ first, and put $\tau=\sigma^9$. For any level $L_c\ge2$, substitute
\[
 \epsilon_s=\sigma,\qquad D=L_c,\qquad
 \delta_s=\epsilon_s^9=\sigma^9=\tau.
\tag{2.33a}\label{p1:eq:2.33a}
\]
Then, term for term,
\[
 \mathcal G(D)
 =\{L_c^{\sigma^2(1+\tau)^n}:n\geq0\},\qquad
 D^{\epsilon_s^2}=L_c^{\sigma^2},\qquad
 D^{\epsilon_s}=L_c^\sigma.
\tag{2.33b}\label{p1:eq:2.33b}
\]
Thus the labels, half-open prime intervals, smooth-prime cutoff, and smooth-divisor support agree term by term with (2.32)--(2.33). These are the Iwaniec parameters at level $L_c$, and the tuple conditions (2.29) are unchanged. With this $\tau$, Theorem 2.6 applies as stated. In the application to (W), Section~5 takes $L_c=D=N^\theta$, so the grid there is exactly the global grid in Theorem~1.2. The coefficients are the values supplied by Iwaniec's construction for $(\mathcal P,L_c,\sigma,\mathbf D)$ at the cutoff $z$ of the application, here $z=N^u$. The cutoff also restricts the prime variables in (2.33). The proof uses only the bound $|\cdot|\le1$ and the support of these coefficients. It never compares coefficient values at different cutoffs.

\cite[Lemma~3]{AA37} gives $(F,f)$ with error
$O(e^{\sqrt K-s}(\log D)^{-1/3})$. In its use here the sieve
constant is $K=K_{\mathrm M}$, fixed independently of $N,D,z$.
Therefore $e^{\sqrt{K_{\mathrm M}}}$ is absorbed into the implied
constant, and the error is $O(e^{-s}(\log D)^{-1/3})$ with that fixed
constant. Likewise, in the profile (2.26), this multiplier changes only
its absolute $O(s)$ constant. The lemma cites
\cite[Lemmas~18 and 20]{AA36}. Theorem 4 includes the
closed endpoints $z=D$ and $z=D^{1/2}$ of (2.30), although
\cite[Lemma~20]{AA36} has strict inequalities.

\subsection{The fundamental lemma, Weil bound, and Heath--Brown identity}

\textbf{Fundamental lemma (cf. \cite[Lemma~4]{BFI86}; the
prime-support condition is explicit in
\cite[Lemma~5]{FI78}).} Let $D\geq2$, $s\geq3$ and
$z=D^{1/s}$. Let $P(z)$ be the product of all primes below $z$.
There are $\lambda_d^\pm$, supported on $d\leq D$, $d\mid P(z)$,
with $\lambda_1^\pm=1$, $|\lambda_d^\pm|\leq1$ and
\[
 (\lambda^-*1)(n)=(\lambda^+*1)(n)=1
 \quad\text{if }(n,P(z))=1,
\]
\[
 (\lambda^-*1)(n)\leq0\leq(\lambda^+*1)(n)
 \quad\text{otherwise},
\]
\[
\begin{aligned}
 \sum_{d\leq D}\frac{\lambda_d^\pm}d
 &=\prod_{p<z}(1-1/p)\\
 &\quad{}\times\{1+O(e^{-s})\}.
\end{aligned}
\tag{2.34}\label{p1:eq:2.34}
\]
The support $d\mid P(z)$ means that every supported $d$ is squarefree
and uses only primes below $z$. In (2.34), this makes
$(\lambda^\pm*1)(n)=\sum_{d\mid(n,P(z))}\lambda_d^\pm$ a sum over
exactly the sieving divisors: if $(n,P(z))=1$, only $d=1$ contributes,
and otherwise FI78's divisor-sum signs apply to $g=(n,P(z))$. The
support is also used in the remainder sum in (2.25). The explicit
Bonferroni weights in (3.102)--(3.104) use the same support to sieve
only primes in $\mathcal P$ below $z$ and to obtain the stated divisor
count. For the AA37 weights in (W), the corresponding condition
$d\mid P_N(Y)$ separates the smooth divisor, whose primes are below
$Y$, from tail primes $p_i\geq D_i\geq Y$. This is used in the unique
recovery of $d$ in the proof of the slot identity (4.13) and in
Appendix~\ref{app:proof-aa37}, Lemma~\ref{lem:aa37-support}.

The error $O(e^{-s})$ suffices below. Removing the primes of an auxiliary modulus changes the local Euler product and needs its own argument. The rough-integer estimates below use explicit Brun--Bonferroni weights; see also \cite{OperaCribro}. The Iwaniec endpoint $s=1$ follows from Theorem 2.6.

\textbf{Weil input.} For every $c\geq1$ and $u,v\in\mathbb Z$, define
\[
 S(u,v;c)=\sum_{x\bmod c,(x,c)=1}e((ux+v\bar x)/c).
\]
The required bound is
\[
 |S(u,v;c)|\leq\tau(c)(u,v,c)^{1/2}c^{1/2}.
\tag{2.35}\label{p1:eq:2.35}
\]
This is the Weil--Estermann bound with full gcd. Its gcd factor covers nonprimitive indices; see \cite[Corollary~11.12]{IK04}. In (2.14), (2.35) serves for absolute convergence. Each quantitative application keeps the gcd and modulus factors in its sums. For prime moduli this is Weil's bound; prime powers use Estermann's bound, and the general modulus follows by the Chinese remainder theorem.

\textbf{Lemma 2.7 (the exact Heath--Brown identity used here).} Let $X\geq2$, $K\geq1$ be an integer, and let $z\geq1$ satisfy $z^K>X$. Then for every integer $1\leq n\leq X$,
\[
 \Lambda(n)=\sum_{j=1}^{K}(-1)^{j-1}\binom Kj
 \sum_{\substack{m_1\cdots m_jn_1\cdots n_j=n\\m_i\leq z}}
 \mu(m_1)\cdots\mu(m_j)\log n_1.
\tag{2.36}\label{p1:eq:2.36}
\]
The identity is \cite[Lemma~1, (6)]{HB82}.

Equation (2.36) is algebraic. For $\Re s>1$, set $M_z(s)=\sum_{m\leq z}\mu(m)m^{-s}$, $Y(s)=\zeta(s)M_z(s)$. The coefficients of $Y-1$ vanish at indices $1\leq n\leq z$. The polynomial identity $1-(1-Y)^K=\sum_{j=1}^K(-1)^{j-1}\binom KjY^j$ gives
\[
 -\frac{\zeta'}\zeta
 =\sum_{j=1}^K(-1)^{j-1}\binom Kj
       (-\zeta')\zeta^{j-1}M_z^j
   +\left(-\frac{\zeta'}\zeta\right)(1-Y)^K.
\tag{2.37}\label{p1:eq:2.37}
\]
Each nonzero factor from $Y-1$ has index $>z$. Hence the last term has no coefficient of index $n\leq X$. Coefficient extraction proves (2.36). For the order-three decomposition on $n\leq2X_m$, take $K=3$, $z=2X_m^{1/3}$. Then $z^3=8X_m>2X_m$. The Möbius variables have support $m_i\leq2X_m^{1/3}$. Dyadic restrictions are imposed after this identity. Coefficient extraction preserves both product constraints and prime ordering.

\subsection{Theorem 2.8}

Theorem 2.8 belongs to the line of incomplete-Kloosterman estimates of
\cite[Theorem~12]{DI82} and
\cite[Theorem~2.1]{Drappeau}; Drappeau states that his theorem extends
DI82 Theorem~12. The closest refined quintilinear estimate is
\cite[Theorem~1.8, (1.8)--(1.9)]{Lichtman23}. Lichtman identifies
\cite{DPR23,ABL21} as antecedents for the exceptional-spectrum and
residue-dependence refinements. The statements are not identical: the
hypotheses (2.38)--(2.40) and bound (2.42) are the variant proved here
from (2.14) in \S\ref{p1:sec:theorem-2-8-proof}. Section~3 uses the full
bound (2.42) in (3.14)--(3.15).

\textbf{Theorem 2.8.}\label{p1:thm:T} Let $C,D,N,R,S\geq1$, $A\geq1$, and $0<\epsilon_0<1/100$. Let $b_{n,r,s}\in\mathbb C$ be supported on
\[
 0<n\leq N,\qquad R<r\leq2R,\qquad S<s\leq2S.
\tag{2.38}\label{p1:eq:2.38}
\]
Let $g:\mathbb R_{>0}^5\longrightarrow\mathbb C$ be smooth, supported in
\[
 (C,2C]\times(D,2D]\times\mathbb R_{>0}^3.
\]
For each multi-index $\nu=(\nu_1,\ldots,\nu_5)\in\mathbb Z_{\geq0}^5$, assume
\[
 |\partial^\nu g(c,d,n,r,s)|
 \leq G_\nu
 \bigl(c^{-\nu_1}d^{-\nu_2}n^{-\nu_3}r^{-\nu_4}s^{-\nu_5}\bigr)^{1-\epsilon_0}.
\tag{2.39}\label{p1:eq:2.39}
\]
The $G_\nu$ are uniform in the family. Let $a\in\mathbb Z$ satisfy $0<|a|\leq A$, and require
\[
 b_{n,r,s}=0\quad\text{unless }(a,rs)=1.
\tag{2.40}\label{p1:eq:2.40}
\]
Set $\mathcal X=2+ACDNRS$. For every $\epsilon>0$,
\[
\begin{split}
 &\left|\sum_{(r,s)=1}\sum_{0<n\leq N}b_{n,r,s}
 \sum_{\substack{c,d\geq1\\(rd,sc)=1}}
 g(c,d,n,r,s)e\!\left(\frac{an\overline{rd}}{sc}\right)\right|\\
 &\qquad\ll_{\epsilon,\epsilon_0,\{G_\nu\}}
 \mathcal X^{\epsilon+18\epsilon_0}J_*(a),
\end{split}
\tag{2.41}\label{p1:eq:2.41}
\]
where the inverse is modulo $sc$, and
\[
\begin{aligned}
 J_*(a)^2={}&
 \sum_{\substack{n''\mid a^\infty\\n''\leq2N}}
 A^{\theta_{\mathrm{KS}}}(n'')^{\theta_{\mathrm{KS}}}\\
 &\quad\times\Bigl[CS\Bigl(\frac N{n''}+RS\Bigr)(C+DR)
 +|a|NRS\Bigr]\|\widetilde b(n'')\|_2^2\\
 &+\Bigl[CS(CD\sqrt R)^{2\theta_{\mathrm{KS}}}
 (N+RS)^{1-\theta_{\mathrm{KS}}}\\
 &\qquad\times(C+DR)^{1-2\theta_{\mathrm{KS}}}
 +D^2NR\Bigr]\|b\|_2^2.
\end{aligned}
\tag{2.42}\label{p1:eq:2.42}
\]
Here $n''\mid a^\infty$ if each prime factor of $n''$ divides $a$. This includes $n''=1$. Also
\[
 \|b\|_2^2=\sum_{n,r,s}|b_{n,r,s}|^2,\qquad
 \|\widetilde b(n'')\|_2^2
 =\sum_{\substack{n',r,s\\(n',a)=1}}|b_{n'n'',r,s}|^2,
\tag{2.43}\label{p1:eq:2.43}
\]
with $b$ extended by zero outside (2.38). When $a=\pm1$, only $n''=1$ occurs in (2.42). The implied constant depends only on the displayed parameters and the seminorm bounds $G_\nu$. It is uniform in $(a,A,C,D,N,R,S,b)$.

For $(r,s)=1$, the condition $(rd,sc)=1$ means $(c,r)=1$ and $(d,sc)=1$. Pairs with $(c,s)>1$ are included, as are both signs of $a$. The sequences and weights may be complex.

\subsection{Proof of Theorem 2.8}\label{p1:sec:theorem-2-8-proof}

The proof uses the large-sieve bounds above, the Hecke bounds of Kim--Sarnak and Deligne, and the trace formula (2.14). We derive the opposite-sign trace formula below.

In this subsection $C,D,N_T,R_T,S_T,A$ are the parameters of Theorem 2.8. Write $R,S,N_1$ for $R_T,S_T,N_T$; Section~3 uses $R,S$ for other scales. Put
\[
\begin{aligned}
&\mathcal X=2+ACDN_1RS,\quad A_a=|a_T|,\quad q=rs,\\
&\quad(r,s)=1,
\quad\theta_{\mathrm{KS}}=7/32,\quad\vartheta_j=2i\kappa_j\in(0,\theta_{\mathrm{KS}}]
\end{aligned}
\tag{2.7.T1}\label{p1:eq:2.7.T1}
\]
on an exceptional eigenspace. All norms below use counting measure. Spectral bases are Petersson-orthonormal.

\subsubsection{Completion and separation}

Fix $n,r,s,c$, put $k=sc$, and define
$\widehat g_c(\xi)=\int g(c,t,n,r,s)e(-\xi t)\,dt$.
Group $d$ in reduced classes modulo $k$ and apply Poisson summation.
This gives the exact identity
\[
\sum_{(d,k)=1}g(c,d,n,r,s)e(a_Tn\overline{rd}/k)
 =k^{-1}\sum_{h\in\mathbb Z}\widehat g_c(h/k)
                       S(h\bar r,a_Tn;k).
\tag{2.7.T2}\label{p1:eq:2.7.T2}
\]
The zero frequency is bounded using
$|c_k(a_Tn)|\le(a_Tn,s)(a_Tn,c)$. Since
$\sum_{c\le2C}(u,c)\le2C\tau(u)$ and
$\sum_{s\le2S}(u,s)^2\le4S^2\tau(u)$, Cauchy--Schwarz gives its
absolute contribution
\[
\ll_\epsilon\mathcal X^\epsilon D(N_1R)^{1/2}\|b\|_2.
\tag{2.7.T3}\label{p1:eq:2.7.T3}
\]
This is the fourth term of (2.42).

Choose a small auxiliary exponent $\epsilon_1>0$. Set
\[
H=4SC D^{\epsilon_0-1}\mathcal X^{\epsilon_1},
\qquad\epsilon_2=\epsilon_1+\epsilon_0.
\tag{2.7.T4}\label{p1:eq:2.7.T4}
\]
Integrate by parts in $d$. Use (2.39) in $d$ on $d\in(D,2D]$. For every $\xi\ne0$ and $j\ge0$, this gives
$|\widehat g_c(\xi)|\ll_jD(|\xi|D^{1-\epsilon_0})^{-j}$. Every $h\ne0$ has $|h|/k\ge|h|/(4SC)$,
because $k=sc\le4SC$. By (2.7.T4) this gives $(|h|/k)D^{1-\epsilon_0}\ge(|h|/H)\mathcal X^{\epsilon_1}$.
Hence, for every $h$ with $|h|>H$,
\[
|\widehat g_c(h/k)|\ll_jD\,\mathcal X^{-j\epsilon_1}(H/|h|)^j .
\tag{2.7.T5}\label{p1:eq:2.7.T5}
\]
For $j\ge2$ and every $H>0$,
$\sum_{|h|>H}(H/|h|)^j\le2+2H$. If $H<1$, the sum runs over all $h\ne0$ and is at most
$2\zeta(2)H^2\le2+2H$. If $H\ge1$ it is at most $2(1+H/(j-1))$. With $|S(h\bar r,a_Tn;k)|\le k$,
the tail $|h|>H$ of (2.7.T2) is $\ll_j(1+H)D\mathcal X^{-j\epsilon_1}$ for each $(c,n,r,s)$. Summing
against $|b_{n,r,s}|$ over $c\le2C$ gives $\ll_j\mathcal X^{7-j\epsilon_1}\|b\|_2$, using
$\sum|b|\le(N_1RS)^{1/2}\|b\|_2$, $C,D,N_1,R,S\le\mathcal X$ and $1+H\le5\mathcal X^{1+\epsilon_1}$. Taking $j>20/\epsilon_1$ makes it smaller than (2.7.T3). If $H<1$, the retained
range $1\le|h|\le\lfloor H\rfloor$ is empty. Every nonzero frequency lies in this tail, including $h=\pm1$. Thus (2.7.T2) contains only the zero frequency and this tail. Equation (2.7.T3) completes this case.

For $1\le|h|\le\lfloor H\rfloor$, write
$h=\varsigma_hm$ and
$\varsigma=\varsigma_h\operatorname{sgn}(a_T)$, with signs
$\varsigma_h,\varsigma\in\{\pm1\}$.
The identity $S(-u,v;k)=S(u,-v;k)$ makes the second Kloosterman
argument $\varsigma A_an$. Partition the integers $m$ into
\[
\mathcal I_j=(M_j,\min(2M_j,\lfloor H\rfloor)]\cap\mathbb Z,
\qquad M_j=2^j/2.
\tag{2.7.T6}\label{p1:eq:2.7.T6}
\]
Thus the first nonempty block is exactly $\{1\}$, with $M=1/2$.
Partition $n$ into $(N',2N']$, $N'\ge1/2$. Shortened final
blocks remain integer intervals. There are $O(\log^2\mathcal X)$
pairs of blocks; each restricted coefficient norm is bounded by the
original norm.
For the current block define
\[
 b^\flat(n'')_{n',r,s}
 =b_{n'n'',r,s}\,1_{\{N'<n'n''\leq2N'\}}\,
  1_{\{(n',a_T)=1\}}.
\]
The slice satisfies $\|b^\flat(n'')\|_2\leq\|\widetilde b(n'')\|_2$.
Restrict $b$ to the current block. By support, $b^\flat(n'')=\widetilde b(n'')$ term by term. Thus the block estimate is exactly (2.42)--(2.43).

For one block write
\[
\begin{aligned}
\frac{\widehat g_c(\varsigma_hm/(sc))}{sc}
 &=\frac D{SC}W_0(c,m,n,r,s),\\
\gamma&=cs\sqrt r,\qquad W_1=W_0\gamma/(CS\sqrt R).
\end{aligned}
\tag{2.7.T7}\label{p1:eq:2.7.T7}
\]
Differentiate under the integral and use (2.7.T4). Then
\[
|\partial^\nu W_1|
 \ll_\nu\mathcal X^{|\nu|\epsilon_2}
            c^{-\nu_1}m^{-\nu_2}n^{-\nu_3}r^{-\nu_4}s^{-\nu_5}.
\tag{2.7.T8}\label{p1:eq:2.7.T8}
\]
The block sum is $D\sqrt R$ times a sum with weight
$\gamma^{-1}W_1S(m\bar r,\varsigma A_an;sc)$.
Put
\[
Z=\frac{4\pi\sqrt{A_aMN'}}{CS\sqrt R},
\qquad x=\frac{4\pi\sqrt{A_amn}}{cs\sqrt r}.
\tag{2.7.T9}\label{p1:eq:2.7.T9}
\]
Invert $c=4\pi\sqrt{A_amn}/(xs\sqrt r)$ and multiply by fixed smooth
cutoffs in $m,n,r,s$, each equal to 1 on that variable's range
$(y,2y]$ ($y=M,N',R,S$) and supported in $[y/2,3y]$. Since
$x/Z=\sqrt{(m/M)(n/N')(R/r)}\,(C/c)(S/s)$ and $C<c\le2C$, the
resulting smooth function $F(x,m,n,r,s)$ is supported on
$1/(12\sqrt3)\le x/Z\le6\sqrt2$, hence on $Z/21\le x\le9Z$.
Mellin inversion in its four log coordinates gives
\[
W_1=(2\pi)^{-4}\int_{\mathbb R^4}\phi_{\mathbf t}(x)
 m^{it_1}n^{it_2}r^{it_3}s^{it_4}\,d\mathbf t,
\tag{2.7.T10}\label{p1:eq:2.7.T10}
\]
\[
\max_{j\le5}\sup_x|x^j\phi_{\mathbf t}^{(j)}(x)|
 \ll\mathcal X^{17\epsilon_2}
                  \prod_{\nu=1}^4(1+|t_\nu|)^{-3}
 =:\mathfrak B(\mathbf t).
\tag{2.7.T11}\label{p1:eq:2.7.T11}
\]
For (2.7.T11), integrate three times by parts in each Mellin variable. There are four variables. Allow five $x$-derivatives. This uses at most 17 derivatives in (2.7.T8). The chain rule uses
$\partial_xc=-c/x,\partial_mc=c/(2m),\partial_nc=c/(2n)$,
$\partial_rc=-c/(2r),\partial_sc=-c/s$.
Thus (2.7.T10) and all ensuing integral interchanges are absolutely
convergent. For fixed $\mathbf t$ set
$\alpha_m=m^{it_1}1_{\mathcal I_j}(m)$ and
$\beta_{n,r,s}=b_{n,r,s}n^{it_2}r^{it_3}s^{it_4}$.
Estimate the following sum.
\[
\mathcal L_{\mathbf t}
 =\sum_{r,s,m,n}\alpha_m\beta_{n,r,s}
  \sum_{(c,r)=1}\gamma^{-1}\phi_{\mathbf t}
       (4\pi\sqrt{mA_an}/\gamma)S(m\bar r,\varsigma A_an;sc).
\tag{2.7.T12}\label{p1:eq:2.7.T12}
\]

\subsubsection{The cusp pair, both signs, and transform bounds}\label{p1:sec:t-cusp-pair}

Use $\sigma_\infty=I$ and
\[
\sigma_{1/s}=
 \begin{pmatrix}\sqrt r&0\\s\sqrt r&1/\sqrt r\end{pmatrix}.
\tag{2.7.T13}\label{p1:eq:2.7.T13}
\]
Multiply $\begin{pmatrix}a&b\\rsc_0&d\end{pmatrix}$ by (2.7.T13). With $c=rc_0+d$, the positive cusp moduli are $\gamma=cs\sqrt r$, $(c,r)=1$, as in (2.13).
Conversely choose $d\equiv c\ (r)$, $(d,sc)=1$, and
$c_0=(c-d)/r$; this supplies every allowed modulus. Double cosets
correspond to the reduced residues modulo $sc$. Insert their top and bottom entries in the cusp Kloosterman definition. This gives
\[
S_{\infty,1/s}(m,\nu;\gamma)
 =e(\nu\bar s/r)S(m\bar r,\nu;sc),
\qquad \nu\ne0.
\tag{2.7.T14}\label{p1:eq:2.7.T14}
\]
The numerator modulo $r$ is $\nu c$. Modulo $sc$ it is $m\bar r\,u+\nu\bar u$, with $u=r\bar d$.
This proves (2.7.T14) for either sign and for complex Fourier
coefficients.

We use the trace formula in DI's coefficient normalization, with the transforms $\widehat\phi,\check\phi,\widetilde\phi$ of (2.12).
Let $m,\nu>0$ and $r>R\ge1$. Use the same-sign Drappeau formula of Section 2. Apply the exact dictionary below.
\[
\rho^{\rm Dr}_{u,\mathfrak a}(n)=\frac{\rho^{\rm DI}_{u,\mathfrak a}(n)}
                                      {2\sqrt{|n|}},\quad
\rho^{\rm Dr}_{\mathfrak c\mathfrak a}(n,t)=
 \frac{\pi^{1/2+it}|n|^{it-1/2}}{\Gamma(1/2+it)}
                    \varphi_{\mathfrak c\mathfrak a n}(1/2+it).
\tag{2.7.T15}\label{p1:eq:2.7.T15}
\]
The holomorphic dictionary is
$\rho^{\rm Dr}_{f,\mathfrak a}(n)=(4\pi n)^{-\ell/2}
\psi_f(\mathfrak a,n)$.
For signs $+,-$, the Maass weights are as follows.
\[
h_+(\kappa)=-\widehat\phi(\kappa)/\cosh\pi\kappa,\quad
h_-(\kappa)=\frac\pi2\check\phi(\kappa)/\cosh\pi\kappa;
\tag{2.7.T16}\label{p1:eq:2.7.T16}
\]
For measure $dt$, the continuous weights are as follows.
\[
e_+(t)=-\widehat\phi(t)/\pi,\qquad e_-(t)=\check\phi(t)/2.
\tag{2.7.T17}\label{p1:eq:2.7.T17}
\]
Only sign $+$ has a holomorphic coefficient. It is
\[
\frac{i^\ell(\ell-1)!}{\pi(4\pi)^{\ell-1}}\,
                                 \widetilde\phi(\ell-1).
\tag{2.7.T18}\label{p1:eq:2.7.T18}
\]
There is no $1/\cosh(\pi t)$ in (2.7.T17).

For the opposite sign the normalization follows from the following derivation. Take a width-one scaling matrix $\sigma_{\mathfrak a}$. Let $\Gamma_{\mathfrak a}$ stabilize $\mathfrak a$ in $\Gamma=\Gamma_0(q)$, including $-I$. Then
$\sigma_{\mathfrak a}^{-1}\Gamma_{\mathfrak a}\sigma_{\mathfrak a}
=\{\pm T^k:k\in\mathbb Z\}$. Define, for $m>0$,
\[
 U_{\mathfrak a m}(z,s)=
 \sum_{\delta\in\Gamma_{\mathfrak a}\backslash\Gamma}
 (\Im\sigma_{\mathfrak a}^{-1}\delta z)^s
 e(m\sigma_{\mathfrak a}^{-1}\delta z),\qquad
 \langle F,G\rangle=\int_{\Gamma\backslash\mathbb H}
 F(z)\overline{G(z)}\,\frac{dx\,dy}{y^2}.
\]
Each coset enters once, with weight one. The series and
unfolding converge absolutely for $\Re s>1$.
Unfold the first series and decompose the second into cusp double
cosets of positive lower-left entry $\gamma$. Translation in the
real variable leaves the Kloosterman factor
$S_{\mathfrak a\mathfrak b}(m,-\nu;\gamma)$, with the
cusp convention of (2.7.T14). Explicitly, a matrix
$\sigma_{\mathfrak b}^{-1}\delta\sigma_{\mathfrak a}
=(\begin{smallmatrix}a&b\\\gamma&d\end{smallmatrix})$
gives phase $e((\nu a-md)/\gamma)$; its negative inverse
has entries $(-d,b;\gamma,-a)$, giving the same phase in
$S_{\mathfrak a\mathfrak b}(m,-\nu;\gamma)$.
The remaining integral is
\[
 I_\gamma=\gamma^{-2s_2}\int_{\mathbb H}
 y^{s_1+s_2-2}|z|^{-2s_2}
 \exp\!\left(2\pi i\{mz-\nu/(\gamma^2z)\}\right)dx\,dy.
\]
The zero-lower-left coset, if present for equivalent cusps,
vanishes by the integral of $e((m+\nu)x)$ over $[0,1]$.
Put $s_-=s_1-s_2$, $s_+=s_1+s_2$,
$z=i\varrho e^{i\alpha}$, $-\pi/2<\alpha<\pi/2$, and
$\varrho=(\sqrt{\nu/m}/\gamma)e^u$. Then
\[
 I_\gamma=\gamma^{-s_+}(\nu/m)^{s_-/2}
 \int_{-\pi/2}^{\pi/2}\cos^{s_+-2}\alpha
 \int_{\mathbb R}e^{s_-u}
 e^{-(4\pi\sqrt{m\nu}/\gamma)\cosh(u+i\alpha)}du\,d\alpha.
\]
For each interior $\alpha$, shift $u+i\alpha$ to the real
line. The vertical sides vanish exponentially, and the inner integral
is $2e^{-is_-\alpha}K_{s_-}(4\pi\sqrt{m\nu}/\gamma)$.
Let $\Re s_i>1$. At either angular endpoint the radial growth is of order $\cos^{- |\Re s_-|}\alpha$. For $\Re s_-=0$, a logarithm may also occur. The remaining exponent is $\Re s_+-2-|\Re s_-|=2\min\Re s_i-2>0$. Hence the original absolute integral is finite.
Hence Fubini and the endpoint limits are legitimate. Applying (2.7.T20) gives the following identity for $\Re s_1,\Re s_2>1$:
\[
\begin{aligned}
\langle U_{\mathfrak a m}(s_1),
             \overline{U_{\mathfrak b\nu}(s_2)}\rangle
={}&\pi2^{3-s_1-s_2}
 \frac{\Gamma(s_1+s_2-1)}{\Gamma(s_1)\Gamma(s_2)}
 (\nu/m)^{(s_1-s_2)/2}\\
&\times\sum_\gamma\gamma^{-s_1-s_2}
 S_{\mathfrak a\mathfrak b}(m,-\nu;\gamma)
 K_{s_1-s_2}(4\pi\sqrt{m\nu}/\gamma).
\end{aligned}
\tag{2.7.T19}\label{p1:eq:2.7.T19}
\]
The polar substitution and contour shift above give
the factor $2K_{s_1-s_2}$. The remaining angular integral is
\[
\int_{-\pi/2}^{\pi/2}\cos^{s_1+s_2-2}\alpha\,
 e^{-i(s_1-s_2)\alpha}\,d\alpha
 =\pi2^{2-s_1-s_2}
   \frac{\Gamma(s_1+s_2-1)}{\Gamma(s_1)\Gamma(s_2)}.
\tag{2.7.T20}\label{p1:eq:2.7.T20}
\]
For an independent normalization, unfold once and use (2.2). This gives
\[
\langle U_{\mathfrak a m}(s),u_j\rangle
 =\overline{\rho_{j\mathfrak a}(m)}
 \frac{\sqrt\pi(4\pi m)^{1/2-s}}{\Gamma(s)}
 \Gamma(s-\tfrac12-i\kappa_j)\Gamma(s-\tfrac12+i\kappa_j).
\]
This is the integral of $y^{s-3/2}e^{-2\pi my}K_{i\kappa_j}(2\pi my)$. Use the standard integral representation of $K$ and a beta integral. The second factor
$\langle u_j,\overline{U_{\mathfrak b\nu}(s_2)}\rangle$
is the same expression with $\rho_{j\mathfrak b}(-\nu)$,
$\nu,s_2$ in place of $\overline{\rho_{j\mathfrak a}(m)}$,
$m,s$. There is no conjugation of $s_2$ in this bilinear
pairing. Multiplication gives exactly (2.7.T21). Applying the same
unfolding to the Eisenstein expansion (2.4), with its standard
Parseval measure, gives (2.7.T22). Thus Parseval gives the spectral side with
\[
\begin{aligned}
P&=\frac{\pi m^{1/2-s_1}\nu^{1/2-s_2}(4\pi)^{1-s_1-s_2}}
          {\Gamma(s_1)\Gamma(s_2)},\\
\Lambda(s_1,s_2;t)&=
 \prod_{j=1}^2\Gamma(s_j-\tfrac12-it)\Gamma(s_j-\tfrac12+it).
\end{aligned}
\tag{2.7.T21}\label{p1:eq:2.7.T21}
\]
Its Maass summand is
$P\overline{\rho_{j\mathfrak a}(m)}
\rho_{j\mathfrak b}(-\nu)\Lambda(s_1,s_2;\kappa_j)$.
Its continuous integral is
\[
\begin{split}
&P\sum_{\mathfrak c}\int
 (\nu/m)^{it}\overline{\varphi_{\mathfrak c\mathfrak a m}(1/2+it)}
 \varphi_{\mathfrak c\mathfrak b,-\nu}(1/2+it)\\
&\qquad\times\Lambda(s_1,s_2;t)|\Gamma(1/2+it)|^{-2}\,dt.
\end{split}
\tag{2.7.T22}\label{p1:eq:2.7.T22}
\]
The factor at the second cusp has index $-\nu$. It arises because
the second Eisenstein parameter in Parseval is conjugated:
$\overline{E(z,1/2+it)}=E(z,1/2-it)$. So the continuous
product pairs a conjugated coefficient at $1/2+it$ with an unconjugated one at $1/2+it$.

For the unscaled Eisenstein series the continuous Parseval measure is
$dt/(4\pi)$. Its coefficient of
$\sqrt yK_{it}(2\pi|n|y)e(nx)$ is
\[
b_{(\mathfrak c,\mathfrak a)}(n,t)=
2\pi^{1/2+it}|n|^{it}\
\frac{\varphi_{\mathfrak c,\mathfrak a,n}(1/2+it)}
     {\Gamma(1/2+it)}.
\]
Hence $(4\pi)^{-1}\overline{b_{(\mathfrak c,\mathfrak a)}(m,t)}
b_{(\mathfrak c,\mathfrak b)}(-\nu,t)$ equals
\[
(\nu/m)^{it}
\overline{\varphi_{\mathfrak c,\mathfrak a,m}(1/2+it)}
\varphi_{\mathfrak c,\mathfrak b,-\nu}(1/2+it)
|\Gamma(1/2+it)|^{-2}.
\]
The coefficient product contributes $4\pi$, canceling the
Parseval denominator. Thus (2.7.T22) has plain $dt$ and no hyperbolic-cosine denominator. The second $\varphi$ is unconjugated at the same $t$.

Write $\sigma_i=\Re s_i>1$,
$\delta=|\sigma_1-\sigma_2|$,
$\sigma_*=\min(\sigma_1,\sigma_2)$, and
$x=4\pi\sqrt{m\nu}/\gamma$, with fixed $m,\nu\ge1$. Replace both Poincaré summands by their moduli. For one double coset with positive lower-left entry, the absolute integral is
\[
\ll_{\sigma_1,\sigma_2,m,\nu}
\gamma^{-\sigma_1-\sigma_2}
\bigl(1+\gamma^\delta+\log(2+\gamma)\bigr).
\tag{2.7.T23}\label{p1:eq:2.7.T23}
\]
The radial integral is a fixed positive factor times
$K_\delta(x\cos\alpha)$. For $\delta>0$, its small-argument
bound is $O_\delta((x\cos\alpha)^{-\delta})$; for $\delta=0$,
it is $O(1+|\log(x\cos\alpha)|)$. The angular weight is
$\cos^{\sigma_1+\sigma_2-2}\alpha$. Thus the power endpoint has
exponent $2\sigma_*-2>0$, while the logarithmic endpoint is
integrable and contributes $O(1+\log(2+\gamma))$ after angular
integration. The implied constant depends on the fixed spectral parameters and
on $m,\nu$, and is uniform in $\gamma$.

Consequently the controlling absolute-sum majorant is
\[
\sum_\gamma N_{\mathfrak a\mathfrak b}(\gamma)
\gamma^{-2\sigma_*}\bigl(1+\log(2+\gamma)\bigr)<\infty.
\tag{2.7.T24}\label{p1:eq:2.7.T24}
\]
For the cusp pair $(\infty,1/s)$, (2.7.T14) gives
$\gamma=cs\sqrt r$, $(c,r)=1$, and
$N(\gamma)=\varphi(sc)\le sc$. The sum is bounded by a constant
times
$\sum_{c\ge1}c^{1-2\sigma_*}(1+\log(2+c))$, which converges
because $\sigma_*>1$. For a general fixed cusp pair, the cited
cumulative count
$\sum_{\gamma\le T}N_{\mathfrak a\mathfrak b}(\gamma)
\ll_{\mathfrak a,\mathfrak b}T^2$ and dyadic decomposition bound
the tail by
$\sum_{k\ge0}(1+k)2^{-2(\sigma_*-1)k}<\infty$. Any bounded
initial range is finite.

The logarithm in (2.7.T23) is attained: for $\sigma_1=\sigma_2=2$, $m=\nu=1$ and the admissible moduli $\gamma=\sqrt2 c$, $c$ odd, the absolute cell integral is $\pi\gamma^{-4}\log\gamma+O(\gamma^{-4})$. The summable majorant (2.7.T24) supplies the uniform step beyond the fixed-$\gamma$ angular endpoint estimate.

Both sides continue to the connected domain
$\Re s_i>3/4,\ |\Re(s_1-s_2)|<1/2$. For the geometric side,
Weil's bound and $K_\mu(x)\ll x^{-|\Re\mu|}(1+|\log x|)$
give locally uniform absolute convergence. For the spectral side, apply Stirling and DI's positive large sieves. These give locally uniform absolute convergence. Since $i\kappa_j\le7/64$, the exceptional gamma factors have no poles here. Equality therefore
holds at $s_1=1+iu,s_2=1-iu$. With
\[
\mathcal H(t,u)=\frac{\cosh\pi t}
                 {\cosh\pi(t-u)\cosh\pi(t+u)},
\tag{2.7.T25}\label{p1:eq:2.7.T25}
\]
this gives
\[
\begin{aligned}
&\sum_\gamma\frac{4\pi\sqrt{m\nu}}{\gamma^2}
 S(m,-\nu;\gamma)K_{2iu}(4\pi\sqrt{m\nu}/\gamma)\\
&\qquad=\frac{\pi^2}2\sum_j\frac{\bar\rho_j(m)\rho_j(-\nu)\mathcal H(\kappa_j,u)}
                      {\cosh\pi\kappa_j}
 +\frac\pi2\sum_{\mathfrak c}\int E_{\mathfrak c}(t)\mathcal H(t,u)\,dt,
\end{aligned}
\tag{2.7.T26}\label{p1:eq:2.7.T26}
\]
where
\[
 E_{\mathfrak c}(t)
 :=(\nu/m)^{it}
   \overline{\varphi_{\mathfrak c\mathfrak a m}(1/2+it)}
   \varphi_{\mathfrak c\mathfrak b,-\nu}(1/2+it).
\]
The factor $|\Gamma(1/2+it)|^{-2}$ of (2.7.T22) is part of the
spectral kernel; $E_{\mathfrak c}$ excludes it.
The constant eigenfunction contributes $0$ to the sum over $j$: its
Fourier coefficients vanish at the nonzero frequencies $m$ and $-\nu$.

To integrate (2.7.T26) against a test function, set
\[
\mathcal R(u)=\int_0^\infty K_{2iu}(x)\phi(x)x^{-2}\,dx,\qquad
k(u)=\frac4{\pi^2}u\sinh(2\pi u)\mathcal R(u).
\tag{2.7.T27}\label{p1:eq:2.7.T27}
\]
Kontorovich--Lebedev inversion gives
$\int_\mathbb R k(u)K_{2iu}(x)\,du=\phi(x)/x$;
see \cite[\S10.43(v), (10.43.30)--(10.43.31)]{DLMF}.
Use the following kernel identity.
\[
\int_\mathbb R u\sinh(2\pi u)\mathcal H(t,u)K_{2iu}(x)\,du
       =\cosh(\pi t)xK_{2it}(x).
\tag{2.7.T28}\label{p1:eq:2.7.T28}
\]
For $|\Re\mu|+|\Re\nu|<1$, the absolutely convergent product integral is
\[
\int_0^\infty K_\mu(x)K_\nu(x)\,dx
 =\frac{\pi^2}
 {4\cos(\pi(\mu+\nu)/2)\cos(\pi(\mu-\nu)/2)}.
\tag{2.7.T29}\label{p1:eq:2.7.T29}
\]
Use the Bessel product representation
$K_\mu K_\nu(x)=2\int_0^\infty K_{\mu+\nu}(2x\cosh v)
\cosh((\mu-\nu)v)\,dv$. Integrate in $x$. For $|\Re b|<1$, use
$\int_0^\infty\cosh(bv)/\cosh v\,dv=(\pi/2)\sec(\pi b/2)$.
The latter follows by expanding $1/\cosh v$ and the partial
fractions of secant. In the subsequent KL application,
$(\mu,\nu)=(2it,2iu)$ with $u\in\mathbb R$; here $t\in\mathbb R$ on
the regular and continuous spectrum, while an exceptional parameter
satisfies $|\Im t|\le 7/64$. Thus $|\Re\mu|+|\Re\nu|\le 7/32<1$ and
$|\Re(\mu-\nu)|<1$. Apply KL inversion to
$xK_{2it}(x)$, whose KL transform is, by (2.7.T29),
$(\pi^2/4)\mathcal H(t,u)/\cosh\pi t$. This proves (2.7.T28).

All interchanges here are justified as follows. Split the integral representation of $K_{iv}$ at $\sinh U=\max(1,1/x)$. Integrate by parts on its tail. This gives
\[
|K_{iv}(x)|\ll e^{-\pi|v|/2}
                    (1+|v|+\log(2+1/x)).
\tag{2.7.T30}\label{p1:eq:2.7.T30}
\]
Apply the self-adjoint operator $(x\partial_x)^2-x^2$ repeatedly to $\phi(x)/x$. Then
$\mathcal R(u)\ll_{\phi,j}e^{-\pi|u|}(1+|u|)^{1-2j}$.
On real $t$, $|\mathcal H(t,u)|\le
4\cosh(\pi t)e^{-2\pi\max\{|t|,|u|\}}$; exceptional parameters
give an exponentially decaying function of $u$ as well.
DI's positive quadratic estimates imply convergence after the
remaining $(1+|t|)^{-4}$ bound. The geometric side is absolutely
summable by Weil. These bounds justify Fubini on both sides of
(2.7.T26). Equations (2.7.T27)--(2.7.T28) now give exactly
(2.7.T16)--(2.7.T17) for sign $-$. For sign $+$, let $r>R\ge1$. Insert (2.7.T15) in the classical same-sign formula. This gives the minus signs in (2.7.T16)--(2.7.T17) and the coefficient (2.7.T18). These are the
normalizations of Section 2, (2.14).

Let $\phi$ be supported in $[Z/21,9Z]$. With $\mathfrak B=\max_{j\le5}\sup x^j|\phi^{(j)}(x)|$, the transform bounds are
\[
|\widehat\phi(t)|+|\check\phi(t)|
 \ll\mathfrak B(1+|\log Z|)(1+Z)^{-1}
       \min\{1,((1+Z)/|t|)^3\}\quad(t\in\mathbb R),
\tag{2.7.T31}\label{p1:eq:2.7.T31}
\]
\[
|\widehat\phi(i\nu)|+|\check\phi(i\nu)|
 \ll\mathfrak B(1+|\log Z|)(1+Z^{-{2\nu}})/(1+Z),
\quad0<\nu\le1/4,
\tag{2.7.T32}\label{p1:eq:2.7.T32}
\]
and the same bound as (2.7.T31) for $\widetilde\phi(\ell)$ with
$|t|$ replaced by $\ell$. To prove them, put
$\Phi(y)=\int\cos(xy)\phi(x)\,dx/x$ and
$\Phi_c(y)=\int e^{-ixy}\phi(x)\,dx/x$.
Five integrations by parts give
$|\Phi^{(j)}(y)|+|\Phi_c^{(j)}(y)|
\ll\mathfrak B Z^j\min(1,(Z|y|)^{-5})$, $j\le3$.
The exact Bessel representations are
\[
\begin{aligned}
\check\phi(t)&=\frac4\pi\int_0^\infty
  \cos(2tu)\Phi(\sinh u)\,du,\\
\widehat\phi(t)&=-2\int_0^\infty
  \cos(2tu)\Phi(\cosh u)\,du,\\
\widetilde\phi(\ell)&=\frac1{2\pi}\int_{-\pi}^{\pi}
  e^{i\ell u}\Phi_c(\sin u)\,du.
\end{aligned}
\]
Their absolute bounds follow by splitting at $\sinh u=1/Z$;
three integrations by parts give the cubic decay. The periodic
holomorphic integral has no boundary terms. Replacing cosine by
$\cosh(2\nu u)$ and splitting at the same point gives (2.7.T32).
Truncate the oscillatory representations first. For real $t$, the
tails are $O((1+|t|)/(x\sinh U))$. For $t=i\nu$, they are
$O((1+\nu)e^{2\nu U}/(x\sinh U))$. Since $2\nu\le1/2$, these
tails tend to zero and justify the interchanges.

\subsubsection{Cusp transport and all Hecke factors}\label{p1:sec:t-cusp-transport}

Choose integers $w,y$ with $rw-sy=1$, and let
\[
W=r^{-1/2}\begin{pmatrix}r&y\\rs&rw\end{pmatrix}.
\tag{2.7.T33}\label{p1:eq:2.7.T33}
\]
Then
$\sigma_{1/s}^{-1}W=\begin{pmatrix}1&y/r\\0&1\end{pmatrix}$,
with $y/r\equiv-\bar s/r(1)$. Also $W\Gamma_0(rs)W^{-1}
=\Gamma_0(rs)$. Here is an integer computation of the latter fact.
For
$A=r^{-1/2}\begin{pmatrix}r\alpha&\beta\\rs\gamma&r\delta\end{pmatrix}$,
$r\alpha\delta-s\beta\gamma=1$, conjugating
$\begin{pmatrix}a&b\\rsc&d\end{pmatrix}$ gives entries
\[
\begin{pmatrix}
 a\alpha\delta r-\alpha b\gamma rs+\beta c\delta rs-\beta d\gamma s&
 -a\alpha\beta+\alpha^2br+\alpha\beta d-\beta^2cs\\
 rs(a\delta\gamma-b\gamma^2s+c\delta^2r-d\delta\gamma)&
 -a\beta\gamma s+\alpha b\gamma rs+\alpha d\delta r-\beta c\delta rs
\end{pmatrix}.
\tag{2.7.T34}\label{p1:eq:2.7.T34}
\]
Both $W=A(1,y,1,w)$ and $W^{-1}=A(w,-y,-1,1)$ have this
form, proving both inclusions. Thus composition by $W$ is unitary,
commutes with the Laplacian, and preserves each eigenspace. The
weight-$\ell$ slash action has the same properties.

Comparison of Fourier series after the translation $y/r$ gives
\[
\rho_{F,1/s}(\nu)=e(\nu\bar s/r)c_{F\circ W}(\nu),\quad
\psi_F(1/s,\nu)=e(\nu\bar s/r)\psi_{F|_\ell W}(\infty,\nu).
\tag{2.7.T35}\label{p1:eq:2.7.T35}
\]
For Eisenstein series, substitute $W\gamma W^{-1}$ in the
defining coset series for $\Re z>1$. It gives
$E_{\mathfrak c}(Wz,u)=E_{W^{-1}\mathfrak c}(z,u)$, using
$W^{-1}\sigma_{\mathfrak c}$ as a width-one scaling matrix.
Changing the source-cusp representative or translating its scaling
matrix does not change the Fourier coefficient at the fixed
observation cusp $\infty$.
Meromorphic continuation therefore gives
\[
\varphi_{\mathfrak c,1/s,\nu}(u)
 =e(\nu\bar s/r)\varphi_{W^{-1}\mathfrak c,\infty,\nu}(u).
\tag{2.7.T36}\label{p1:eq:2.7.T36}
\]
In (2.7.T12), (2.7.T14) contributes
$e(-\varsigma A_an\bar s/r)$. Equations (2.7.T35)--(2.7.T36)
cancel this phase exactly for both signs, in every spectrum.

For $(k,q)=1$, define the normalized Hecke operator below.
\[
T_kF(z)=k^{-1/2}\sum_{ad=k}\sum_{b\bmod d}F((az+b)/d).
\tag{2.7.T37}\label{p1:eq:2.7.T37}
\]
Sum each Fourier term, of index $\nu'$, over $b$. This forces
$d\mid\nu'$, and the term moves to index $\nu=a\nu'/d$. The
coefficient at $\nu\ne0$ is
\[
c_{T_kF}(\nu)=\sum_{a\mid(k,\nu)}c_F(k\nu/a^2).
\tag{2.7.T38}\label{p1:eq:2.7.T38}
\]
The prefactor is 1 because, in the normalization (2.2) of
\cite[(1.15)]{DI82}, the factor $\sqrt y$ contributes $\sqrt{a/d}$,
and $k^{-1/2}d\sqrt{a/d}=1$.
For weight $\ell$, use
$T_k=k^{(\ell-1)/2}\sum_{ad=k}d^{-\ell}\sum_b(\cdot)$;
the normalized coefficients
$a_F(\nu)=\nu^{-(\ell-1)/2}\psi_F(\infty,\nu)$
satisfy the same identity at coprime $(k,\nu)$.

Split $n=n'n''$, $n''\mid a_T^\infty$, $(n',a_T)=1$,
and put $k=A_an''$. Then $(k,q)=(k,n')=1$.
The number of pieces is
\[
\mathcal D\le(2N_1)^\epsilon
                  \prod_{p\mid a_T}(1-p^{-\epsilon})^{-1}
 \ll_\epsilon\mathcal X^{2\epsilon},
\quad \tau(k)^2\ll_\epsilon\mathcal X^\epsilon.
\tag{2.7.T39}\label{p1:eq:2.7.T39}
\]
The product bound is uniform in $a_T$. Choose $p_\epsilon$ so that $(1-p^{-\epsilon})^{-1}\leq p^\epsilon$ for $p>p_\epsilon$. The smaller primes contribute a constant depending only on $\epsilon$. The remaining product is at most $|a_T|^\epsilon$.
On a Maass eigenspace $V$, the transported $n''$-functional is
$Y\circ T_k\circ W$, where
$Y(v)=\sum_{n'}\beta_{n'n''}c_v(\varsigma n')$.
By Riesz representation and unitarity,
\[
\sum_{u_j\in V}|\ell_{2,n''}(u_j)|^2
 \le\tau(k)^2k^{\theta_{\mathrm{KS}}}\sum_{u_j\in V}|Y(u_j)|^2.
\tag{2.7.T40}\label{p1:eq:2.7.T40}
\]
The self-adjoint unramified $T_k$ diagonalizes on the full
level space. By Kim--Sarnak its norm is at most
$\tau(k)k^{\theta_{\mathrm{KS}}/2}$. For holomorphic forms replace $k^{\theta_{\mathrm{KS}}}$
by 1, by Deligne. For negative index, $c_v(-n')=\overline{c_{\bar v}(n')}$. Conjugation permutes an orthonormal basis. Hence the positive-index quadratic bound holds with conjugated coefficients. The argument allows complex
coefficients and any orthonormal basis of each eigenspace.

The continuous spectrum requires its own proof. Represent a cusp by
$u/w$, $w\mid q$, put $w'=(w,q/w)$. The orbit conditions in
\cite[Lemma~3.6, (3.9)--(3.11)]{DI82}
$(c_1,d_1)=1,(c_1,q)=w,c_1d_1\equiv uw\pmod{ww'}$, with $c_1=w\gamma$,
give exactly
\[
\begin{aligned}
\varphi_{u/w,\infty,n}(z)&=
 (w'/(wq))^z
 \sum_{\gamma\ge1,(\gamma,q/w)=1}\gamma^{-2z}\\
&\qquad
 \sum_{\substack{d_1\bmod\gamma w\\(d_1,\gamma w)=1,\ \gamma d_1\equiv u(w')}}
                         e(-nd_1/(\gamma w)).
\end{aligned}
\tag{2.7.T41}\label{p1:eq:2.7.T41}
\]
For a character $\chi\bmod w'$, let
$\mathcal V_\chi(n;z)=\sum_{u\bmod w'}^*\chi(u)\varphi_{u/w,\infty,n}(z)$.
For $(k,q)=(k,n')=1$, either sign of $n'$, the identity is
\[
\mathcal V_\chi(kn';z)=\bar\chi(k)
          \left(\sum_{d\mid k}\chi^2(d)d^{1-2z}\right)\mathcal V_\chi(n';z).
\tag{2.7.T42}\label{p1:eq:2.7.T42}
\]
To prove it, sum (2.7.T41) over $u$. Write $c_1=c_1'c_1''$, where
$c_1''$ is supported on primes of $q$ and $(c_1',q)=1$.
CRT factors the inner sum as
$\chi(c_1')G_\chi(-n;c_1'')c_{c_1'}(n)$.
The outer character supplies another $\chi(c_1')$.
Replacing the Gauss-sum variable by $\bar kd_1$ gives
$\bar\chi(k)$ on $c_1''$. The other factor is
\[
L^{(q)}(2z,\chi^2)^{-1}
     \sum_{d\mid|n|,(d,q)=1}\chi^2(d)d^{1-2z},
\]
by $c_{c_1'}(n)=\sum_{d\mid(c_1',n)}d\mu(c_1'/d)$; this divisor sum
factors at $kn'$. All manipulations converge absolutely for
$\Re z>1$. Both sides are meromorphic Fourier coefficients,
so (2.7.T42) holds meromorphically, hence at $z=1/2+it$ outside
a discrete set. There are finitely many cusp characters.
Orthogonality in $u\bmod w'$ gives
\[
\begin{split}
&\sum_{\mathfrak c}\left|\sum_{n'}v_{n'}|kn'|^{it}
            \varphi_{\mathfrak c,\infty,kn'}(1/2+it)\right|^2\\
&\qquad\le\tau(k)^2\sum_{\mathfrak c}\left|\sum_{n'}v_{n'}|n'|^{it}
            \varphi_{\mathfrak c,\infty,n'}(1/2+it)\right|^2.
\end{split}
\tag{2.7.T43}\label{p1:eq:2.7.T43}
\]
Here the character eigenvalue is
$\bar\chi(k)\sum_{d\mid k}\chi^2(d)(k/d^2)^{it}$,
of modulus at most $\tau(k)$. Transport (2.7.T36) permutes the
complete cusp set. For negative indices,
$\varphi_{\mathfrak c\infty,-n'}(1/2+it)
=\overline{\varphi_{\mathfrak c\infty,n'}(1/2-it)}$.
Changing $t$ to $-t$ and conjugating coefficients converts
each symmetric integration window to DI's positive-index
continuous quadratic form. The discrete exceptional set is null
for the $dt$-integral.

\subsubsection{Regular, continuous, and holomorphic estimates}

Put $Q_T=RS$, $L_{n''}=N'/n''\ge1/2$. Include $|t|<1$ among the spectral windows $K$. By (2.7.T31), the undivided transform satisfies the following bound.
\[
\mathfrak h(K)\ll\mathfrak B(\mathbf t)\mathcal X^{o(1)}
 (1+Z)^{-1}\min\{1,((1+Z)/K)^3\}.
\tag{2.7.T44}\label{p1:eq:2.7.T44}
\]
For Maass forms use the positive DI weights $1/\cosh\pi\kappa_j$. For holomorphic forms use the positive weights $(\ell-1)!/(4\pi)^{\ell-1}$. These weights are split equally
between the two Cauchy--Schwarz factors. For the continuous spectrum use plain $dt$. Take the supremum of $e_\varsigma(t)$ outside the positive quadratic form.
By (2.7.T17), $|e_\varsigma|\le|\Phi_\varsigma|/2$, with
$\Phi_+=\widehat\phi,\Phi_-=\check\phi$. Thus every norm uses a
positive measure.

\cite[Theorem~2]{DI82} at $\infty$ gives the $m$-side
\[
\sum_{q\asymp Q_T}\tau(q)(K^2+M^{1+\epsilon}/q)M
 \ll\mathcal X^{3\epsilon}(Q_TK^2+M)M.
\tag{2.7.T45}\label{p1:eq:2.7.T45}
\]
For the $n$-side, Cauchy--Schwarz over $n''$ costs $\mathcal D$. Apply (2.7.T40) or (2.7.T43). Then use \cite[Theorem~2]{DI82} at scale $L_{n''}$. The bound is
\[
\ll\mathcal D\sum_{n''}\tau(k)^2k^{\theta_{\mathrm{KS}}}
             (K^2+L_{n''}^{1+\epsilon}/Q_T)\|b^\flat(n'')\|_2^2.
\tag{2.7.T46}\label{p1:eq:2.7.T46}
\]
On continuous and holomorphic parts (2.7.T46) holds with $k^{\theta_{\mathrm{KS}}}$
replaced by 1. \cite[Theorem~2]{DI82} allows the endpoint $M=1/2$
in (2.7.T45).

Use $K^2+U\le((1+Z)^2+U)\max(1,K/(1+Z))^2$.
The sum over dyadic $K$ of the latter factor times the cubic
decay in (2.7.T44) is $O(1+\log(2+Z))$. The cubic decay also absorbs a factor $K^\epsilon$: with it the sum is $O((1+Z)^\epsilon\log(2+Z))$. Thus the bound (2.6) in the form $(K^2+N/q)(KN)^\epsilon$ gives the same estimate. Cauchy--Schwarz and
(2.7.T39) therefore yield
\[
\begin{aligned}
|\mathcal L_{\mathbf t}^{\rm reg}|
\ll{}&\mathcal X^{4\epsilon}\mathfrak B(\mathbf t)(1+Z)^{-1}
 [(Q_T(1+Z)^2+M)M]^{1/2}\\
&\times\left[\sum_{n''}(A_an'')^{\theta_{\mathrm{KS}}}
 ((1+Z)^2+L_{n''}/Q_T)\|b^\flat(n'')\|_2^2\right]^{1/2}.
\end{aligned}
\tag{2.7.T47}\label{p1:eq:2.7.T47}
\]
The same calculation with plain $dt$ proves (2.7.T47) for
the continuous part. It uses the symmetry shown after
(2.7.T43). The same cubic summation proves absolute convergence
of its integral, justifying the window partition. Thus the
continuous part uses the measure $dt$ throughout.

\subsubsection{Exceptional spectrum and the one-integer block}\label{p1:sec:t-exceptional}

Write $L=L_{n''}$, $\Xi=1/Z$, $\Omega=1+Q_T/L$.
Since $0<\vartheta_j\le\theta_{\mathrm{KS}}$,
\[
1+\Xi^{\vartheta_j}
 \le2(1+\Xi/\sqrt{\Omega})^{\vartheta_j}\Omega^{\vartheta_j/2}.
\tag{2.7.T48}\label{p1:eq:2.7.T48}
\]
Fix $n''$. Apply Cauchy--Schwarz, assigning
$Y=(1+\Xi/\sqrt{\Omega})^2$ to the $m$-side and
$\Omega$ to the $n$-side. With $\ell_{2,n''}$ as in (2.7.T40), the two exceptional quadratic forms are
\[
\begin{aligned}
\mathcal M_{n''}&=\sum_{r,s}\sum_{\rm exc}^{(rs)}Y^{\vartheta_j}
 \Bigl|\sum_m\alpha_m\rho_{j\infty}(m)\Bigr|^2,\\
\mathcal N_{n''}&=\sum_{r,s}\sum_{\rm exc}^{(rs)}\Omega^{\vartheta_j}
 |\ell_{2,n''}(u_j)|^2,
\end{aligned}
\]
with $r,s$ over the support of $b$. Let $L\ge1/2$. Apply (2.7.T40) and Theorem 2.2 at level $rs>Q_T$, length $L$ and $Y=\Omega$. Since $L\Omega/(rs)\le1+L/Q_T$, the factor in (2.8) is $\ll\max(1,L)^\epsilon(1+L/Q_T)$. Summing over $r,s$ gives
\[
\mathcal N_{n''}\ll\mathcal X^{2\epsilon}(A_an'')^{\theta_{\mathrm{KS}}}
                 (1+L/Q_T)\|b^\flat(n'')\|_2^2.
\tag{2.7.T49}\label{p1:eq:2.7.T49}
\]
Put $Q'=4Q_T$, $Y_0(M)=(Q'+M)^2/M$, and
\[
Z_{n''}=\frac{C^2S^2R}{(Q_T+M)^2(L+Q_T)}.
\tag{2.7.T50}\label{p1:eq:2.7.T50}
\]
The identity
$\Xi^2/\Omega=(16\pi^2A_an'')^{-1}
C^2S^2R/[M(L+Q_T)]$ gives
\[
\max(1,Y/Y_0(M))^{\theta_{\mathrm{KS}}}
 \le1+(A_an'')^{-\theta_{\mathrm{KS}}}Z_{n''}^{\theta_{\mathrm{KS}}}.
\tag{2.7.T51}\label{p1:eq:2.7.T51}
\]
Let $M\ge1$. Then (2.9) applies to every integer interval
inside $(M,2M]$. Use the positive Hilbert space with weights $Y^{\vartheta_j}$. Apply Abel summation to $m^{it_1}$. The norm is at most $1+|t_1|\log2$ times the maximum interval norm. Take the maximum after the norm. It remains outside the spectral sum. Apply (2.9) at $Y_0(M)$. Then use monotonicity
\[
Y^{\vartheta_j}\le\max(1,Y/Y_0)^{\theta_{\mathrm{KS}}} Y_0^{\vartheta_j},
\tag{2.7.T52}\label{p1:eq:2.7.T52}
\]
For fixed $q$, the coefficient $\alpha_m$ is independent of the
factorization $q=rs$. Put
$w(q)=\#\{(r,s)\text{ in the current support}:rs=q\}$ and
$\Delta=\max_{2\le q\le Q'}w(q)$, and let $\omega_0(q)$
count the distinct prime factors of $q$. Since $(r,s)=1$,
$w(q)\le2^{\omega_0(q)}\le\tau(q)\ll_\eta q^\eta$.
In the Hilbert space over all levels $2\le q\le Q'$ and their exceptional
forms, vector-valued Abel summation gives
$\mathcal M_{n''}^{1/2}\le\Delta^{1/2}
(1+|t_1|\log2)\sup_{I\subset(M,2M]}
F_Y(I)^{1/2}$, where
$F_Y(I)=\sum_{2\le q\le Q'}\sum_{\rm exc}^{(q)}Y^{\vartheta_j}
|\sum_{m\in I}\rho_{j\infty}(m)|^2$ and $I$ ranges over integer
intervals. Thus the supremum is outside the whole quadratic form.
Apply (2.9) with exponent $\eta=\epsilon/2$ at $Y_0(M)$.
By (2.7.T4), $M\le H\le4\mathcal X^2$ when $\epsilon_1\le1$;
also $Q'=4Q_T\le4\mathcal X$. Consequently
$\Delta(Q'M)^\eta\ll_\epsilon\mathcal X^{2\epsilon}$.
Together with (2.7.T51)--(2.7.T52), this is absorbed in
$\mathcal X^{3\epsilon}$ and gives
\[
\mathcal M_{n''}\ll\mathcal X^{3\epsilon}(1+t_1^2)
 [1+(A_an'')^{-\theta_{\mathrm{KS}}}Z_{n''}^{\theta_{\mathrm{KS}}}](Q_T+M)M.
\tag{2.7.T53}\label{p1:eq:2.7.T53}
\]

\cite[(8.17)]{DI82} excludes $M=1/2$. We prove that endpoint here.
Let $p(q)$ be the least prime not dividing $q$. Every smaller
prime divides $q$. The classical Chebyshev lower bound
$\Theta(x):=\sum_{\ell\le x}\log\ell\gg x$ for $x\ge2$
(see \cite{IK04}) gives
\[
p(q)\ll1+\log q.
\tag{2.7.T54}\label{p1:eq:2.7.T54}
\]
Here $\Theta(p(q)-1)\le\log q$, because the product of all
smaller primes divides $q$; the case $p(q)=2$ is immediate.
The implied constant in (2.7.T54) is absolute and independent of $a_T$.

On an exceptional eigenspace $V$, let $v$ represent
$u\mapsto c_u(1)$. Self-adjointness and (2.7.T38) show that
$T_pv,T_{p^2}v$ represent $c_u(p),c_u(p^2)$.
The Hecke identity $T_p^2=T_{p^2}+I$, with
$\|T_p\|\le2p^{\theta_{\mathrm{KS}}/2}$, gives
\[
\sum_{u_j\in V}|c_{u_j}(1)|^2
 \le8p^{\theta_{\mathrm{KS}}}\sum_{u_j\in V}|c_{u_j}(p)|^2
       +2\sum_{u_j\in V}|c_{u_j}(p^2)|^2.
\tag{2.7.T55}\label{p1:eq:2.7.T55}
\]
Multiply by $Y^{\vartheta_j}$, which is constant on $V$. Group
levels by $p(q)\le P:=C_0(1+\log(4Q_T))$, where $C_0$ is an
absolute constant supplied by (2.7.T54). Positivity permits
enlarging each group to all $2\le q\le Q'$. For $k=p,p^2$, the singleton $\{k\}$ is an integer interval in $(k/2,k]$, with $k/2\ge1$, so (2.9) applies at $N=k/2$. In DI's terms, subtract the initial intervals $(k/2,k]$ and $(k/2,k-1]$. The resulting factor 4 is included in $c_\epsilon$.

For the endpoint, the same factorization count multiplies the
level-$q$ quadratic form by at most $\tau(q)$. The following bound
holds for any positive exponent. In its use in (2.7.T57), take
$\eta=\min(\epsilon/2,1/4)$ and use
$\tau(q)\ll_\epsilon Q_T^{\epsilon/2}$ for
$Q_T<q\le4Q_T$. This keeps their combined loss within the
displayed $\epsilon$-power after adjustment of its constant. At
$Y_k=(Q'+k/2)^2/(k/2)$ it gives
\[
\sum_{2\le q\le Q'}\sum_{\rm exc}^{(q)}
 Y^{\vartheta_j}|c_{u_j}(k)|^2
 \ll_\epsilon (Q'k)^\epsilon(Q'+k/2)k
                          \max(1,Y/Y_k)^{\theta_{\mathrm{KS}}}.
\tag{2.7.T56}\label{p1:eq:2.7.T56}
\]
Since $Y_0(1/2)/Y_k=k(Q'+1/2)^2/(Q'+k/2)^2\le k\le9k/4$, insert (2.7.T56), with its
exponent taken as $\eta$, in (2.7.T55).
Sum over the at most $P$ possible primes:
\[
\begin{aligned}
&\sum_{Q_T<q\le4Q_T}\tau(q)\sum_{\rm exc}^{(q)}
 Y^{\vartheta_j}|c_{u_j}(1)|^2\\
&\qquad\ll_\epsilon Q_T^\epsilon P^6
 (Q_T+\tfrac12)\tfrac12\,
 \max(1,Y/Y_0(1/2))^{\theta_{\mathrm{KS}}},
\end{aligned}
\tag{2.7.T57}\label{p1:eq:2.7.T57}
\]
for $0<\epsilon\le1/4$, with the implied constant including
the divisor bound and the change from $\epsilon$ to $\eta$.
Before simplification, the sum is at most a constant times
\[
c_\epsilon(9/4)^{\theta_{\mathrm{KS}}} Q'^{1+\epsilon}
 \sum_{p\le P}(8p^{2+2\theta_{\mathrm{KS}}+\epsilon}
                         +2p^{4+2\theta_{\mathrm{KS}}+2\epsilon});
\]
this is $\ll_\epsilon Q_T^{1+\epsilon}P^6$, because
$5+2\theta_{\mathrm{KS}}+2\epsilon\le6$.
The loss $P^6$ is logarithmic and uniform in $a_T$.
In particular $P^6\ll_\epsilon\mathcal X^\epsilon$; together
with the factorization loss already included in (2.7.T57), this
fits the $\mathcal X^{3\epsilon}$ allowance of (2.7.T53).
Together with (2.7.T51), this proves (2.7.T53) also for $M=1/2$.
Since $\alpha_1=1$, Abel summation costs nothing here.

Apply the triangle inequality over $n''$ before Cauchy--Schwarz.
Then (2.7.T49), (2.7.T53), and (2.7.T32) give
\[
\begin{aligned}
|\mathcal L_{\mathbf t}^{\rm exc}|
\ll{}&\mathcal X^{6\epsilon}\mathfrak B(\mathbf t)(1+|t_1|)
 (1+Z)^{-1}\sqrt{M(Q_T+M)}\\
&\times\sum_{n''}\big[(A_an'')^{\theta_{\mathrm{KS}}/2}
                 +Z_{n''}^{\theta_{\mathrm{KS}}/2}\big]
        ((L+Q_T)/Q_T)^{1/2}\|b^\flat(n'')\|_2.
\end{aligned}
\tag{2.7.T58}\label{p1:eq:2.7.T58}
\]
There is no exceptional continuous or holomorphic spectrum.

\subsubsection{The four terms and completion of the proof}

Put
\[
\mathfrak T=A_aMN'/(C^2S),\quad M_0=CS/D,\quad
\varpi=4\mathcal X^{\epsilon_2}.
\tag{2.7.T59}\label{p1:eq:2.7.T59}
\]
The quantity $\mathfrak T$ is unrelated to the cell endpoint $T$ of
Section~3. Then $M\le\varpi M_0$,
$Z^2Q_T=16\pi^2\mathfrak T$, so
$(1+Z)^{-2}Q_T^{-1}\asymp(Q_T+\mathfrak T)^{-1}$.
Square (2.7.T47) and multiply by $D^2R$. Apart from the stated auxiliary powers, sum the following over $n''$.
\[
(A_an'')^{\theta_{\mathrm{KS}}} D^2R
 \frac{M(M+Q_T+\mathfrak T)(L+Q_T+\mathfrak T)}
                       {Q_T+\mathfrak T}\|b^\flat(n'')\|_2^2.
\tag{2.7.T60}\label{p1:eq:2.7.T60}
\]
The first, noncancelling branch of (2.7.T58) has the same bound. Termwise, use
\[
(Q_T+M)(L+Q_T)/Q_T
 \le[Q_T(1+Z)^2+M][(1+Z)^2+L/Q_T].
\tag{2.7.T61}\label{p1:eq:2.7.T61}
\]
To obtain (2.42), use the following identity.
\[
\frac{Q_T(M+Q_T+\mathfrak T)(L+Q_T+\mathfrak T)}
                            {Q_T+\mathfrak T}
=(L+Q_T)(M+Q_T)+Q_T\mathfrak T
                         -\frac{LM\mathfrak T}{Q_T+\mathfrak T}.
\tag{2.7.T62}\label{p1:eq:2.7.T62}
\]
Multiply (2.7.T62) by $M/Q_T$, then drop the last (nonpositive) term:
\[
\frac{M(M+Q_T+\mathfrak T)(L+Q_T+\mathfrak T)}{Q_T+\mathfrak T}
\le\frac{M(L+Q_T)(M+Q_T)}{Q_T}+M\mathfrak T.
\]
Since
$D^2M/S\le\varpi CD$ and
$M+Q_T\le\varpi S(C+DR)/D$, it follows that
\[
\begin{aligned}
&D^2R\frac{M(M+Q_T+\mathfrak T)(L+Q_T+\mathfrak T)}
                            {Q_T+\mathfrak T}\\
&\qquad\le\varpi^2
 \{CS(L+RS)(C+DR)+A_aN'RS\}.
\end{aligned}
\tag{2.7.T63}\label{p1:eq:2.7.T63}
\]
The last term follows from
$D^2RM\mathfrak T=A_aN'RS(DM/(CS))^2$.
Using $L\le N_1/n''$, $N'\le N_1$, $A_a\le A$ proves
the first two terms of (2.42), with
$A^{\theta_{\mathrm{KS}}}(n'')^{\theta_{\mathrm{KS}}}$ kept explicit.

For the cancelling branch of (2.7.T58), use
\[
Z_{n''}^{\theta_{\mathrm{KS}}}(L+Q_T)
 \le(C^2S^2R)^{\theta_{\mathrm{KS}}}(Q_T+M)^{-2\theta_{\mathrm{KS}}}(N_1+Q_T)^{1-\theta_{\mathrm{KS}}},
\tag{2.7.T64}\label{p1:eq:2.7.T64}
\]
and $\sum_{n''}\|b^\flat(n'')\|_2
\le\mathcal D^{1/2}\|b\|_2$. Since
$M\mapsto M(Q_T+M)^{1-2\theta_{\mathrm{KS}}}$ is increasing,
substitute $M_0$, paying at most $\varpi^{2-2\theta_{\mathrm{KS}}}$.
The resulting identity is
\[
\begin{aligned}
&\frac{D^2RM_0(Q_T+M_0)^{1-2\theta_{\mathrm{KS}}}
                  (C^2S^2R)^{\theta_{\mathrm{KS}}}(N_1+Q_T)^{1-\theta_{\mathrm{KS}}}}{Q_T}\\
&\qquad=CS(CD\sqrt R)^{2\theta_{\mathrm{KS}}}
             (N_1+RS)^{1-\theta_{\mathrm{KS}}}(C+DR)^{1-2\theta_{\mathrm{KS}}}.
\end{aligned}
\tag{2.7.T65}\label{p1:eq:2.7.T65}
\]
This is the third term of (2.42). The fourth was (2.7.T3).
The truncated tail is smaller than it.

To account for all losses, reconstruct the block sum as
$D\sqrt R(2\pi)^{-4}\int\mathcal L_{\mathbf t}\,d\mathbf t$.
Equation (2.7.T11) gives
$\int\mathfrak B(\mathbf t)(1+|t_1|)\,d\mathbf t
\ll\mathcal X^{17\epsilon_2}$.
The amplitude's $\varpi$ loss in (2.7.T63) is at most
$\mathcal X^{\epsilon_2}$ after taking its square root.
The log partitions and divisor factors cost $\mathcal X^{O(\epsilon_1)}$. This includes the DI small $\epsilon$ losses, $\mathcal D$, and the one-integer $P^6$ loss. The bound is uniform in $a_T$.
Consequently the final loss is
\[
\mathcal X^{O(\epsilon_1)+18\epsilon_0}.
\tag{2.7.T66}\label{p1:eq:2.7.T66}
\]
Choose the freely available $\epsilon_1$ small enough relative
to the requested $\epsilon$. Equations (2.7.T63), (2.7.T65), and (2.7.T3)
then prove (2.41)--(2.43), with loss
$\mathcal X^{\epsilon+18\epsilon_0}$, for both signs and
including the $M=1/2$ block. This proves Theorem 2.8 in the form used in Section~3, (3.14)--(3.15).


\section{Dispersion for abstract sources}\label{sec:pi-psdt}

This section proves a local implication for an abstract source.
It uses Linnik's dispersion method \cite{Linnik63}. The second-dispersion
setup has close precedents in Maynard \cite{MaynardII} and Lichtman
\cite[\S5]{Lichtman23}; the proof below tracks its own coefficient
masks, restoration terms, and moving residue.
The local hypotheses are (H1)--(H5). The source-wide estimate (H6)
and the box-family coefficient bound (H7) are used only in the
global summation in Section~\ref{sec:pi-global}.
Each application verifies the required conditions for its own measures.
The geometric and source partition results appear below.
Theorem 2.8 is stated in
Section 2, (2.38)--(2.43), and derived from the classical inputs in
\S\ref{p1:sec:theorem-2-8-proof}. We keep its full moving-frequency bracket.
The local estimate has three terms.
The restoration error $E_{\rm GR}$ is bounded apart from the spectral term.

\subsection{Notation and source hypotheses}

\emph{Local notation.} Put $L=\log X$; \S\ref{sec:pi-global} writes $\log N$ explicitly. $T$ is the lower endpoint of the selected Type-II cell, as in Lemma~\ref{lem:aa37-typeII-lengths} for the prime source. Here $\vartheta_{\rm sp}=7/32$. The zero-mode functional is $\mathcal Z(V)$ of (3.31); $Z$ is the scale of $z$. $B$ is a Heath--Brown box. $D$ and $C$ are the scales of $b$ and $c$ in (3.62), as in Theorem 2.8. $B_{z,y}$ is the count (3.70). $S$ is the scale of $\gamma_2$. $S_B^{\rm PSDT}$ denotes the signed box sum in (3.12). Type I-1 and I-2 denote boxes with one and two large variables (\S\ref{sec:pi-nonII}). Since $\log X=\log N+O(\log\log N)$, fixed powers of $\log X$ and $\log N$ are comparable.
$I_T$ is the selected integer interval of a Type-II cell; $I_N$ is reserved for the source interval of Theorem~\ref{thm:aa37-interior}.
Some letters have local meanings, fixed at each use.
The letter $C$ is a fixed logarithmic exponent in (3.13), (3.92) and similar bounds, and the scale of $c$ from (3.62) on.
Similarly, $\varepsilon$ is the source cut of Theorem~\ref{thm:aa37-interior}, $\epsilon$ is a small fixed exponent, and $\varepsilon_y\in\{\pm1\}$ is a sign.
The letter $Y$ is a summation length in (3.16)--(3.19) and (3.35).
It is the auxiliary scale $H_*SF$ in the spectral step near (3.70).
In Theorem~\ref{thm:aa37-interior} it is the smooth-prime cutoff $D^{\sigma^2}$.
$\mathcal T$ is the tuple set of \S\ref{sec:pi-psdt}; it is unrelated to $T$.

Let $N$ be a sufficiently large even integer and let $a$ be a positive
integer. Fix $0<K_{\rm scale}<\infty$ and

\[
0<\sigma\le10^{-10},\quad \tau=\sigma^9,\quad\delta=30\sigma,
\quad\vartheta_{\rm sp}=7/32,\quad\eta=\sigma/1000,
\]
\[
\eta_L=10^{-4},\qquad
\ell_0=\frac{53}{412}-\frac{\eta_L}{2}
       =\frac{264897}{2060000},\qquad
c_\sigma=\ell_0\sigma^2/4.
\tag{3.1}
\]

The local product scale is
$N/(\log N)^{K_{\rm scale}}\le X\le N$.
All implied constants in this section may depend on the fixed
$K_{\rm scale}$ and the fixed source bounds. They may also depend on
the parameters displayed at the point of use. Put
\[
\begin{gathered}
L=\log X,\quad z_0=X^{1/(\log\log X)^3},\quad
u_0=N^{\ell_0\sigma^2},\\
p_0=\min\{p\text{ prime}:p\nmid 2aN\}.
\end{gathered}
\tag{3.2}
\]
In particular, $N/u_0=N^{1-\ell_0\sigma^2}$.

Thus $p_0\ge3$ and $p_0\nmid aN$.
Every prime below $p_0$ divides $2aN$.
Chebyshev's bound $\sum_{p<p_0}\log p\gg p_0$ gives
$p_0\ll\log(2aN)$. By C4 below, $a\le X(\log N)^{K_{\rm scale}}$,
so $p_0\ll\log N<z_0\le u_0$ eventually. Write

\[
\mathcal S_t(q;a)=1_{t\equiv a\ (q)}-
                    \frac{1_{(t,q)=1}}{\varphi(q)}.
\tag{3.3}
\]

A \textbf{Class-II box} consists of the following data.
These are source hypotheses.

\textbf{C1. Exact unrestricted coefficient.} Fix a source term, a prefix cell,
a sign of the well-factorable remainder in Iwaniec's theorem \cite[Theorem~4]{AA37},
and a tail-label vector $b$. Its coefficient is exactly

\[
C_b(q)=\sum_{d,(p_j)}\lambda_d
                  1_{q=d\prod_{j\in J}p_j},\qquad |\lambda_d|\le1,
\tag{3.4}
\]

where each $p_j$ ranges independently over its fixed prime set and
$p_j\ge u_L$.
Here $u_L\ge u_0$ is the source's lower cutoff for labelled tail primes,
and $P(u_L;N)=\prod_{p<u_L,\,p\nmid2N}p$.
The variable $d\mid P(u_L;N)$ is squarefree with the
cap of that theorem. The full modulus is at most $X^{0.6}$.
Every supported full modulus is prime to $2N$.
Selectors are constant on the block.
Exact identities remove the equal-tail, prefix-diagonal, natural-comparator,
cut, and selector-transfer corrections.
Each is a separate contribution outside (3.4).

\textbf{C2. Exact factors.} There is a labelled partition
$J=J_1\sqcup J_2\sqcup J_3$, with $d$ attached to $J_1$, such that

\[
C_b=C_{\rm slot}(\gamma_1*\gamma_2*\gamma_3),\qquad
C_{\rm slot}=\prod_{i=1}^3|J_i|!=O_\sigma(1),\qquad |\gamma_i|\le1.
\tag{3.5}
\]

The last two factors are $u_0$-rough. If $J_s=\varnothing$ for $s=2,3$, then $\gamma_s(m)=\mathbf1_{m=1}$. The factor $\gamma_1$ always contains the smooth-divisor sum of Iwaniec's theorem. If $J_1=\varnothing$, then $\gamma_1(m)=\sum_d\lambda_d\mathbf1_{m=d}$.
Equality (3.5) holds on all supported moduli,
with repeated-prime multiplicities.
The source's labelled-tuple map proves it.
It covers moduli with repeated primes as well as squarefree ones.

\textbf{C3. Packing.} Localize dyadically the first factor $\gamma_1$
at $\asymp Q$, $\gamma_3$ at $\asymp R$, and $\gamma_2$
at $\asymp S$. Here $T$ is the lower endpoint of the selected Type-II
cell; $N_I$ is its integer length.
Enlarge the sufficiently-large threshold for $N$ to absorb fixed dyadic constants. Then

\[
Q<TX^{-29\sigma},\qquad QRS<M_I X^{-89\sigma},
\qquad N_I\ge QX^{\sigma/3},\quad M_I\ge X^{1/2}.
\tag{3.6}
\]

Here $Q$ is the scale of the first factor $\gamma_1$ (the outer
slot scale, written $Q_{\rm slot}$ in Appendix~\ref{app:proof-aa37}).
It is not the scale of the full modulus $q$ of (3.4), which is $QRS$.
Thus (3.6) gives the first inequality in (3.7).
We list the four spectral packing inequalities together.
\[
Q<TX^{-\sigma},\quad T^2R^2S<X^{1-8\sigma},
\quad A_0^{\vartheta_{\rm sp}} TQR^2S^5<X^{2-15\sigma},
\]
\[
(TR/S)^{2\vartheta_{\rm sp}}TQR^2S^5<X^{2-15\sigma},
\qquad A_0=X^{1+\sigma/20}.
\tag{3.7}
\]
Each application must verify the stated packing and interval
conditions for its own source. Appendix~\ref{app:proof-aa37} does so
for the prime source in Theorem~\ref{thm:aa37-interior}.

\textbf{Removed squareful and other restricted terms.}
A source may remove forbidden tails, forbidden prefixes or
retained allowed repeats.
Their contributions are bounded separately.
One admissible bound is $O_\sigma(N^{1-\ell_0\sigma^2}L^4)$,
after all source-label and representation multiplicities.
Theorem~\ref{thm:main} concerns the unrestricted coefficient (3.4).
The removed terms need their own estimates.

The prime source in Theorem~\ref{thm:aa37-interior} is
$\mathbf1_{\mathbb P}(n)\mathbf1_{N/(\log N)^B<n<(1-\varepsilon)N}$.
Appendix~\ref{app:proof-aa37} verifies its coefficient, interval,
cutoff, support, and boundary conditions. The hypotheses above also
state the local estimate for any other source that satisfies them.

Apply the finite HB expansion below.
Then split into disjoint short product boxes of
relative width $L^{-C_{\rm box}}$. Boundary boxes contribute separately.
On each interior Class-II box select an HB variable in the window.
Maximize its lower-endpoint exponent $e_i=\log_N V_i$.
Break ties by the fixed variable index, as in the classification in
Appendix~\ref{app:proof-aa37}.
Call it $n$, and the product of the other variables $m$.

In (3.8), allow a roughness cutoff $\zeta\ge z_0$,
where $z_0=z_0(X)$ is the baseline in (3.2).
All estimates below are uniform in $\zeta$.
Upper bounds use only $P^-(n)\ge z_0$.
The SW proof (3.25) for the selected factor uses the cutoff $\zeta$.
It remains uniform under stronger roughness.

\[
\alpha_n\in\left\{\mu(n),1,\frac{\log n}{L}\right\}
1_{P^-(n)\ge \zeta}1_{I_T}(n),\quad |\alpha_n|\le1,
\qquad |\beta_m|\ll\tau_{23}(m).
\tag{3.8}
\]

The fixed HB binomial coefficients and tags may be kept as exterior
scalars. $I_T\subset(T,2T]$ has length $N_I\asymp TL^{-C_{\rm box}}$.
The $m$-interval is majorized by

\[
\psi_I(m)=\phi_0((m-m_0)/M_I),\qquad
M_I\asymp (X/T)L^{-C_{\rm box}},
\tag{3.9}
\]

Fix $\phi_0\in C_c^\infty([-1,1])$, with $0\le\phi_0\le1$ and $\phi_0=1$ on $[-1/2,1/2]$.
It is independent of $N$, $X$, the source and the box.
It is also independent of all shell and tuple variables.
For every $j\ge0$, the seminorms $\|\phi_0^{(j)}\|_1$ and $\|\phi_0^{(j)}\|_\infty$ are fixed constants.
The support of $\psi_I(m)=\phi_0((m-m_0)/M_I)$ lies in the doubled interval.
The majorant depends on the box only through $m_0$ and $M_I$.
On that entire support the cut geometry gives

\[
1\le |mn-a|\le X^3.
\tag{3.10}
\]

\textbf{C4. Moving frequency.} All original moduli are prime to $a$,
and $0<a\le X(\log N)^{K_{\rm scale}}\le A_0$. This is the moving
frequency in the spectral theorem. It may vary with $N$.

\textbf{C5. Window.} The source selection gives

\[
N^{60\sigma}\le T\le N^{1/3+14\sigma},\qquad
X^{60\sigma}\le T\le X^{1/3+15\sigma}.
\tag{3.11}
\]

The second assertion follows from the first and
$\log(N/X)/\log N=o(1)$, with $N_0$ chosen after fixed
$\sigma$. \textbf{Claim P (nonresonance):} the doubled support satisfies (3.10).
The source assumptions for the non-II
extension are in \S\ref{sec:pi-nonII}.

\subsection{Statement and theorem inputs}

\textbf{Length hypothesis for the selected integer interval (LEN).}
Uniformly over the selected boxes, the interval has length
$N_I\asymp T L^{-C_{\rm box}}$ and satisfies
$N_I\ge(2T)^{1-2^{-126}}$.
The integer interval used in (3.17) has this same length, with the
endpoint convention specified by the source.
Lemma~\ref{lem:aa37-typeII-lengths}, (B.22)--(B.27), proves
this for the prime source.
Each application proves it for its own source.

For the coefficient of one Heath--Brown box define

\[
\Delta_B(q;a)=\sum_{\mathbf v\in B}w(\mathbf v)
                 \mathcal S_{\prod v_i}(q;a),\qquad
S_B^{\rm PSDT}=\sum_q C_B(q)\Delta_B(q;a).
\tag{3.12}
\]

\subsection{An abstract dispersion estimate}

The local argument is source-independent.
We state it before the verification for specific sources.
Let $\mathcal B$ be a finite
family of dyadic boxes at scale $X$, and let $a$ be a moving frequency.
For $B\in\mathcal B$ write $C_B(q)$ for its signed modulus coefficient.
Define $\Delta_B(q;a)$ by (3.12).

The following hypotheses define the class used in this paper.

\begin{description}[leftmargin=2.4em,style=nextline]
\item[(H1) Exact coefficient] On every box, $C_B(q)$ has the exact
form (3.4), with its fixed labelled prime sets, squarefree
smooth-divisor variable $d$, and repeated-prime multiplicities. Every
supported modulus is coprime to $2a$.

\item[(H2) Three-factor decomposition] The labelled variables split
into three groups satisfying (3.5). The last two factors are
$u_0$-rough, as in C2. The identity holds on all
supported moduli, including repeated-prime multiplicities.

\item[(H3) Geometric packing and box data]  After dyadic localization,
the scales $Q,R,S$ and the test scale $T$ satisfy (3.6)--(3.7).
On each box, the Heath--Brown coefficients factor as $\alpha_n\beta_m$
with (3.8).  The smooth majorant has the form (3.9).
Its center satisfies $|m_0|\le C_{\rm ctr}X/T$, where $C_{\rm ctr}$
is a fixed source-class constant, uniform over all boxes, shells,
tuple variables, and permitted $N,X,a$.
Its full doubled support satisfies (3.10).
The selected integer interval satisfies LEN, including its strong
length bound and integer endpoint convention.
For this abstract dispersion estimate, assume modulus support in
$q\le X^{\vartheta_0}$, with $\vartheta_0$ as in (3.12a).

\item[(H4) Moving frequency]  The frequency is the positive integer $a$,
with $N$ even, $0<a\le X(\log N)^{K_{\rm scale}}\le A_0$, and all original moduli are coprime
to $a$.

\item[(H5) Selected window]  The tested scale obeys (3.11), with all
parameters fixed after $\sigma>0$ is fixed.

\item[(H6) Source-wide remainders and boundary] Let a source supply
a logarithm-freezing estimate and boundary estimate for the full box family.
After summing over all source labels, shells, boxes, and representation
multiplicities, these estimates must give
\[
E_{\rm log}\ll N(\log N)^{c_{\log}-C_{\rm box}},\qquad
E_{\rm bdry}\ll N(\log N)^{c_R-C_{\rm box}/2},
\]
with fixed $c_{\log}$ and $c_R$. This condition is used in the
global summation, not in the local theorems below.

\item[(H7) Box-family coefficient bound] The finite box family
$\mathcal B=\{B_j\}$ carries outer coefficients $c_j$, coming from
the source labels and the Heath--Brown identity. At fixed $\sigma$
they satisfy (3.91):
$\sum_j|c_j|\le C_{\sigma,C_{\rm box},\mathcal C_{\rm src}}(\log N)^{K_0}$ with a fixed
$K_0\le24C_{\rm box}+26$. This condition is used only for the
aggregate bound (3.92), not for the per-box estimates.
\end{description}

Write $\mathcal C_{\rm src}$ for the fixed source constants in
(H1)--(H5), such as $C_{\rm slot}$, $C_{\rm ctr}$ and the implied
constants in (3.8)--(3.11).

The Siegel--Walfisz bound (3.25) for the selected factor follows in
\S\ref{sec:pi-sw} from the classical theorem \cite[Corollary~5.29]{IK04}. BFI \cite{BFI86} and
Theorem 2.8, proved from classical inputs in
\S\ref{p1:sec:theorem-2-8-proof}, are theorem inputs, not source hypotheses.
The source-wide boundary bound (3.24) is hypothesis (H6). It is used
in the source estimates, not in the per-box theorem (3.13).

Put
\[
 \Theta_0=\frac{19101}{32000},\qquad
 \vartheta_0=\Theta_0-\eta_L>\frac12.
\tag{3.12a}
\]
The number $\Theta_0$ is the small-atom empty-prefix packing level
of Proposition~\ref{prop:pack-ceiling} in \S\ref{sec:pack}.
The exponent $\vartheta_0$ is a support convention in the abstract
dispersion estimate. The actual support level comes from the packing inequalities
(3.6)--(3.7) after all Iwaniec factors are restored.
In Theorem~\ref{thm:aa37-interior}, (B.2)--(B.4) give
$q<N^{e_Q}$ with $e_Q<\Theta(u)\le\Theta(0)=153/256<0.6$; hence
$q\le X^{0.6}$ on every source shell for sufficiently large $N$.
The fixed exponent $\vartheta_0$ belongs to the abstract support
hypothesis of the introductory source-box estimate. The variable ceiling
$\Theta(u)=(5-7/32-3u)/8$ in Theorem~\ref{thm:aa37-interior}
controls the Iwaniec coefficient family as $u$ varies.
At $u=1/500$, $\Theta(u)=19101/32000=\Theta_0$,
whereas $\vartheta_0=\Theta_0-10^{-4}$.
The introductory estimate applies to its stated abstract support class;
the Iwaniec result at this endpoint uses the $X^{0.6}$ theorem below.
Theorem~\ref{thm:standalone} is the abstract dispersion estimate under
the support and packing hypotheses in (H3), including the stronger
support $q\le X^{\vartheta_0}$.
Theorem~\ref{thm:main} is the dispersion engine used for the actual
source support. Theorem~\ref{thm:aa37-interior} in the application
uses that engine together with the source estimates of
Appendix~\ref{app:proof-aa37}. At $u=1/500$ it gives the limiting
ceiling $\Theta(1/500)$ for this specified signed mean value.

\begin{theorem}[Dispersion estimate under the (H3) packing hypotheses]\label{thm:standalone}
\textup{(a)} Assume (H1)--(H5). Then, for every fixed $A'>0$ and
every box $B$,
\[
 \left|\sum_{q\le X^{\vartheta_0}}C_B(q)\Delta_B(q;a)\right|
 \ll_{\sigma,K_{\rm scale},C_{\rm box},A',\mathcal C_{\rm src}}
 XL^{-A'}+X^{1-c_\sigma}L^C,
\tag{3.12b}
\]
where $c_\sigma=\ell_0\sigma^2/4>0$ and
$C=C(\sigma,C_{\rm box},A',K_{\rm scale},\mathcal C_{\rm src})$ is the exponent of (3.13).
The implied constant and the threshold for $X$ are uniform in the box,
the even integer $N$, and all moving frequencies $a$ permitted by
(H1)--(H5). They may
depend on the fixed source bounds and on $\sigma$. They are
ineffective, since the proof uses the Siegel--Walfisz theorem in (3.25).
Under (H3) every supported modulus satisfies $q\le X^{\vartheta_0}$, so
the truncation in (3.12b) is vacuous and its left side is $|S_B^{\rm PSDT}|$.

\textup{(b)} Assume in addition (H7). Fix $A>0$ and take
$A'=A+24C_{\rm box}+31$ in \textup{(a)}. Then the weighted sum
$E_{\rm II}=\sum_j|c_j||S_{B_j}^{\rm PSDT}|$ over the box family
satisfies (3.92), that is,
$E_{\rm II}\ll_{\sigma,A,C_{\rm box},K_{\rm scale},\mathcal C_{\rm src}}
N(\log N)^{-A}$.
This bound is uniform in the scales $X$ of the boxes. For
$N/(\log N)^{K_{\rm scale}}\le X\le N$ one has
$\log X\ge\log N-K_{\rm scale}\log\log N\ge\tfrac12\log N$ for large $N$,
so $X(\log X)^{-A'}\le2^{A'}N(\log N)^{-A'}$ and
$X^{1-c_\sigma}(\log X)^C\le N^{1-c_\sigma}(\log N)^C$.
Part \textup{(b)} holds for any fixed $C_{\rm box}>0$. The condition
$C_{\rm box}\ge\lceil2A+32\rceil$ of Theorem~\ref{thm:aa37-interior}
serves only the source-wide errors of (H6).

Under the support hypothesis in (H3), the theorem gives arbitrary
logarithmic saving for these signed sums through
$q\le X^{\vartheta_0}=X^{0.59680625}$, where $\vartheta_0>1/2$.
\end{theorem}

\begin{proof}
The support bound assumed in (H3) makes the sum in (3.12b)
the full sum $S_B^{\rm PSDT}$ of (3.12).
Hypotheses (H1)--(H5) imply those of Theorem~\ref{thm:main}.
Hence Theorem~\ref{thm:main} applies to $B$ and gives (3.13), which is (3.12b).
This proves \textup{(a)}. Part \textup{(b)} follows from \textup{(a)}
and (H7) by the summation (3.92) in \S\ref{sec:pi-global}.
\end{proof}

\begin{theorem}[Dispersion theorem]\label{thm:main}
Assume (H1)--(H5), where in (H3) the support bound $q\le X^{\vartheta_0}$
may be relaxed to $q\le X^{0.6}$.
Thus each box satisfies (3.8)--(3.10) and LEN.
Then, for each fixed $A'>0$,

\[
|S_B^{\rm PSDT}|\ll_{\sigma,K_{\rm scale},C_{\rm box},A',\mathcal C_{\rm src}}
 XL^{-A'}+X^{1-c_\sigma}L^C.
\tag{3.13}
\]

Here $C=C(\sigma,C_{\rm box},A',K_{\rm scale},\mathcal C_{\rm src})$.
It multiplies only power-saving terms.
It enters only the final choice of $N_0$.
The exponent $C$ is fixed after the proof parameters.
It is uniform in the box, each class in the random partition, the residue,
and the moving shift.
Constants and $N_0$ may depend on the fixed $\sigma$.
All bounds are for fixed $\sigma>0$ as $N\to\infty$.
The implied constants and the threshold $N_0$ are ineffective, since
the proof uses the Siegel--Walfisz theorem in (3.25).
The estimate (3.13) is per box and does not use (H6) or (H7).
If (H7) also holds and $A'=A+24C_{\rm box}+31$, then the weighted sum
over the box family satisfies (3.92).
\end{theorem}

Hypotheses (H1)--(H5) give the coefficient, factor identity, packing, box,
frequency and window conditions, including LEN and the integer endpoint convention.
BFI \cite[Theorem~0(a)]{BFI86} supplies (3.48), Theorem 2.8 of
\S\ref{p1:sec:theorem-2-8-proof} supplies (3.14)--(3.15), and
(3.25) supplies SW.
The source-wide condition (H6) enters the source error terms, including
(B.29) in the proof of Theorem~\ref{thm:aa37-interior} and the auxiliary
bound (3.95). It is not used in the per-box Theorem~\ref{thm:main}.
The box-family bound (H7) enters only (3.92).
The spectral part gives a term $X^{1-1.725\sigma}L^C$ (\S\ref{sec:pi-spectral}).
It is absorbed in $X^{1-c_\sigma}L^C$, since $c_\sigma<\sigma/2<1.725\sigma$.

The dependencies have the following order:
\[
\begin{gathered}
\text{source identities, box geometry and source estimates}\ \Longrightarrow\ (\mathrm{H1})\text{--}(\mathrm{H7}),\\
\text{Section 2 and (3.25)}\ \Longrightarrow\ \text{BFI, Theorem 2.8 and SW},\\
\text{(H1)--(H5)}+\text{theorem inputs}\ \Longrightarrow\ (3.13),\\
(3.13)+(\mathrm{H7})\ \Longrightarrow\ (3.92).
\end{gathered}
\]
The hypotheses are stated independently of any particular source.
Appendix~\ref{app:proof-aa37} verifies them for
Theorem~\ref{thm:aa37-interior}.
Source comparison and the estimates for removed squareful terms are separate steps.
Aggregate logarithm-freezing estimates enter in \S\ref{sec:pi-global}.
Theorem~\ref{thm:main} uses (H1)--(H5) and the stated theorem inputs.

The inputs are the following theorems, printed in Section 2.
The Siegel--Walfisz theorem (SW) \cite[Corollary~5.29]{IK04} gives prime distribution in reduced progressions
to moduli at most a fixed power of the logarithm.
\cite[Theorem~0(a), (2.2)]{BFI86}, with \cite[(A2)]{BFI86},
gives the BDH quadratic estimate.
It applies to sequences satisfying its norm-normalized SW condition.
On the domain $(rd,sc)=1$, Theorem 2.8 gives

\[
\begin{aligned}
&\left|\sum_{(r,s)=1}\sum_{0<n\le N_T}b_{n,r,s}
 \sum_{c,d:(rd,sc)=1}g(c,d,n,r,s)
 e\left(\frac{a_Tn\overline{rd}}{sc}\right)\right|\\
&\qquad\ll (2+A_T C D N_T R_T S_T)^{\epsilon+18\epsilon_0}J_*(a_T).
\end{aligned}
\tag{3.14}
\]

Here $0<|a_T|\le A_T$, where $A_T$ is the frequency bound $A$ of Theorem 2.8, and $b=0$ unless $(a_T,rs)=1$. The positive
scales meet the endpoint conditions of Section 2. The bracket is

\[
\begin{aligned}
J_*^2={}&\sum_{n''\mid a_T^\infty}A_T^{\vartheta_{\rm sp}}(n'')^{\vartheta_{\rm sp}}
 \bigl[C S_T(N_T/n''+R_TS_T)(C+DR_T)\\
&\qquad+|a_T|N_TR_TS_T\bigr]\|b^\flat(n'')\|_2^2\\
&+\bigl[C S_T(CD\sqrt{R_T})^{2\vartheta_{\rm sp}}
 (N_T+R_TS_T)^{1-\vartheta_{\rm sp}}\\
&\qquad\times(C+DR_T)^{1-2\vartheta_{\rm sp}}
       +D^2N_TR_T\bigr]\|b\|_2^2.
\end{aligned}
\tag{3.15}
\]

Here $b^\flat(n'')$ is the dyadic-block slice of $\tilde b(n'')$
with $(n',a)=1$. The slice satisfies
$\|b^\flat(n'')\|_2\le\|\tilde b(n'')\|_2$.
Since $b$ is supported on the present dyadic block,
$b^\flat(n'')=\tilde b(n'')$ term by term.
Hence (3.15) is exactly (2.42)--(2.43).

The smooth derivatives in (3.14) obey the variable-based hypothesis
\eqref{p1:eq:2.39}, with the exponent $1-\epsilon_0$ displayed there:
each derivative is scaled by the corresponding positive argument
$c,d,n,r,s$ itself. In particular, the initial interval
$0<n\le N_T$ is not treated as a dyadic block of scale $N_T$.
Theorem 2.8 is used in the form of Section 2, (2.38)--(2.43).
It is proved in \S\ref{p1:sec:theorem-2-8-proof}.
Its spectral identity, including the continuous spectral measure, is exactly (2.14).
The I-2 estimate (Lemma~\ref{lem:typeI2}) also uses the composite-modulus Weil--Estermann bound \cite{Weil48,Estermann61} in \S\ref{sec:pi-nonII}.

The sum in (3.13) is a signed modulus sum; triangle inequalities over
source labels and boxes are applied outside it.

\subsection{Elementary preparation, moments, and SW}\label{sec:pi-sw}

\textbf{Divisor moments.} For integers $K,k\ge1$,

\[
\tau(n)^K\le\tau_{2^K}(n),\qquad
\sum_{n\le Y}\frac{\tau_k(n)}n\le(1+\log Y)^k.
\tag{3.16}
\]

At $p^v$, the first inequality is
$(v+1)^K\le\binom{v+2^K-1}{v}$.
The ratio of successive right/left quotients is at least 1.
This holds since $(1+t)^K\le1+(2^K-1)t$ for $0\le t\le1$, by convexity.
The second inequality follows by summing over the $k$ factors. Also
$\tau_j(n)\le\tau(n)^{j-1}$.

Suppose $Y^{1-1/k}\le H\le Y$. Select a largest factor in a
$k$-factorization. Then

\[
\sum_{Y<n\le Y+H}\tau_k(n)
 \le k\sum_{d\le(2Y)^{1-1/k}}\tau_{k-1}(d)(H/d+1)
 \le3kH(1+\log(2Y))^{k-1}.
\tag{3.17}
\]

Let $1\le u\le Y$ and $0<\varpi\le1$.
Split off prime factors larger than $Y^{\varpi/2}$.
Greedily group the rest into factors at most $Y^\varpi$.
There are at most $\lfloor2/\varpi\rfloor+1$ groups.
The large-prime part has at most $2/\varpi$ prime factors with multiplicity. Hence

\[
\tau(u)^K\le2^{2K/\varpi}
 \sum_{d\mid u,\ d\le Y^\varpi}
       \tau(d)^{K(\lfloor2/\varpi\rfloor+1)}.
\tag{3.18}
\]

For $u=|mn-a|\le X^3$, take $Y=X^3,\varpi=1/12$.
Put $h=(d,n)$. If $h\nmid a$, the congruence $d\mid mn-a$ has no solutions.
Otherwise it has at most $|J|h/d+1$ solutions on $J$.
Sum (3.18) and apply (3.16). For $|J|\le3M_I$ this gives

\[
\sum_{m\in J}\tau(|mn-a|)^K
 \le2^{24K}(1+\log X)^{2^{25K}}\{3M_I\tau(n)^{25K+1}+X^{1/4}\}.
\tag{3.19}
\]

Write $\widehat I_M$ for the integer support of $\psi_I$;
it is contained in an interval of length at most $3M_I$.
Sum (3.19) with $K=5$ over $n\in I_T$.
The $n$-sum of $\tau(n)^{126}\le\tau_{2^{126}}(n)$ is bounded by (3.17) and LEN.
Since $M_I\ge X^{1/2}$, the term $N_IX^{1/4}$ is absorbed by $M_IN_I$. Thus

\[
\sum_{n\in I_T}\sum_{m\in\widehat I_M}
                 \tau(|mn-a|)^5\ll M_IN_IL^{c_1},
\qquad c_1=2^{125}+2^{126}+2.
\tag{3.20}
\]
Here $c_1$ is an upper bound; the exponent $2^{125}+2^{126}-1$ already suffices.

For the prime source, Lemma~\ref{lem:aa37-typeII-lengths}
takes $T$ as the lower endpoint of the selected cell.
Under threshold (B.24), it proves (B.27), namely
$N_I\ge(2T)^{1-2^{-126}}$.
In (3.17) take $Y'=\lfloor T\rfloor$ and $H'=N_I$.
The integers in $I_T=(T,r_XT]$ of (B.22) are exactly
$(Y',Y'+N_I]$.
The stated bound implies $H'\ge(2Y')^{1-2^{-126}}$.
The mesh uses the lower endpoint of the dyadic shell, which is a different scale.
For the product source, LEN gives the lower length and the integer-interval convention.
Thus (3.17) applies to $\tau_{2^{126}}$, after increasing $N_0(\sigma)$ for each fixed
$\sigma>0$.
The same argument with $\varpi=\kappa/6$ treats
intervals of length at least $X^\kappa$.
It gives $\sum_m|\beta_m|^2\ll M_IL^c$ by elementary means; Shiu's estimate \cite{Shiu80} would also suffice.

\noindent\begin{minipage}{\linewidth}
\textbf{Finite Heath--Brown identity \cite{HB82}.} For $U\ge1$ and $n\le U^3$,

\[
\Lambda(n)=\sum_{j=1}^3(-1)^{j-1}\binom3j
 \sum_{n=m_1\cdots m_jn_1\cdots n_j,\ m_i\le U}
           \mu(m_1)\cdots\mu(m_j)\log n_1.
\tag{3.21}
\]
\end{minipage}

For the proof put $M(s)=\sum_{m\le U}\mu(m)m^{-s}$.
The series $1-\zeta M$ has no coefficients of index at most $U$. Hence
$-(\zeta'/\zeta)(1-\zeta M)^3$ has none of index at most $U^3$.
Expand the cube and use $-(\zeta'/\zeta)\zeta^j=(-\zeta')\zeta^{j-1}$.
Let $X$ be the dyadic shell of the tested coordinate, and $t\in(X,2X]$.
Take $U=2X^{1/3}$. Then (3.21) applies, since $t\le2X<U^3=8X$.
For the product source use $U_i=2X_i^{1/3}$ for the shell $X_i$ of each coordinate.
There are at most 24 variables.
Roughness factors exactly over the product. The identity

\[
1_{\mathbb P}(t)=\frac{\Lambda(t)}{\log t}
               -\sum_{k\ge2,\ t=p^k}\frac1k
\tag{3.22}
\]

\noindent has its proper-prime-power error bounded at source level in \S\ref{sec:pi-global}.

\textbf{Logarithms and boundaries.} On one short box replace each $1/\log p$
by its value at the box reference point. The pointwise error is
$O(L^{-C_{\rm box}-2})$ per coordinate. To sum the discrepancy
absolutely, use

\[
\sum_q|C_b(q)|\,|\mathcal S_t(q;a)|
 \ll_\sigma\tau(|t-a|)+L,\qquad t\ne a.
\tag{3.23}
\]

The first term counts divisors of $t-a$.
The second uses $\sum_{q\le X^{0.6}}1/\varphi(q)\ll L$.
The unrestricted labelled coefficient is $O_\sigma(1)$.
The number of prime labels is fixed,
and the small-prime part is disjoint from them.
For the prime source, the pointwise majorant in
Appendix~\ref{app:proof-aa37} gives (3.23) with
$F_a(n)=d(a-n)+2(1+L)$.
The source-wide logarithm-freezing estimate is (B.18):
$|E_{\rm LOG}|\ll_\sigma NL^{18-C_{\rm box}}$.
It is summed over tuples, each lying in one box, so its box-count
factor is 1. Here the logarithm is $L_N=\log N$, and its divisor-moment
input is (B.17). The box width is expressed in $L_X=\log X$.
Fixed powers of $L_X$ and $L_N$ are comparable.
For the product source, the corresponding aggregate bound is a
source condition stated in \S\ref{sec:pi-global}.

Write $L_N=\log N$. Consider the source cuts, the HB caps, and the monomial ordering and radial cuts of the product source.
On each dyadic source shell, each boundary region has relative width
$O(L_N^{-C_{\rm box}})$ in an active coordinate. Counting that
coordinate first and summing the geometric shell scales gives a total
of $O(NL_N^{5-C_{\rm box}}+N^{6/7}L_N^5)$ boundary indices, including
the fixed divisor multiplicity. An active coordinate can be chosen
among the positive-power factors, so the integer endpoint terms have a
fixed power saving. For the prime source, Cauchy--Schwarz with this
count and the source-wide second moment in (B.17) gives

\[
E_{\rm bdry}\ll N(\log N)^{c_R-C_{\rm box}/2}.
\tag{3.24}
\]

For the prime source in Theorem~\ref{thm:aa37-interior}, (3.24)
holds with $c_R=16$ by the global estimate (B.19) in
Appendix~\ref{app:proof-aa37}. This verifies (H6) for that source.

The prime-source estimate is invoked under
Theorem~\ref{thm:aa37-interior}, whose residue range supplies $a\ge N$.
This lower bound is compatible with the local upper bound:
$a\le N\log N/3\le(2/3)X(\log N)^{B+1}$ when
$X>N/(2(\log N)^B)$.
The prime-source cut and $a\ge N$ give
$|mn-a|\ge\varepsilon N/2$ on the doubled support.
Thus resonant boxes are empty and (3.10) holds.

\textbf{SW for the selected factor.} For every $A_{\rm SW}>0$, subinterval $I\subset I_T$,
integer $d\ge1$, and reduced class $l\bmod k$,

\[
\sum_{n\in I,\ n\equiv l(k),\ (n,d)=1}\alpha_n
 -\frac1{\varphi(k)}\sum_{n\in I,\ (n,dk)=1}\alpha_n
 \ll_{A_{\rm SW}} T\tau(d)L^{-A_{\rm SW}}.
\tag{3.25}
\]

Here and below the interval indicator is used only once.
Let $k>L^K$. After $K$ is chosen large, the trivial bound is
$O(T/k+1+T/\varphi(k))\ll TL^{-A_{\rm SW}}$.
The elementary bound $\varphi(k)\gg k/\log\log(k+2)$, or
$\varphi(k)\gg k^{1/2}$, suffices.
Let $k\le L^K$ and put $y_1=\exp(L^{2/3})$.
For sufficiently large $X$, $z_0>y_1$, so no selected integer
has largest prime below $y_1$.
Integers divisible by $p^2$, $p\ge y_1$, contribute $O(T/y_1)$.
Every other integer has a unique form $n=Pm$, with $P>P^+(m)$, $P\ge y_1$,
and $m\le2T/y_1$. For fixed $m$, its coefficient is
$\mu(m)\mu(P)=-\mu(m)$, $1$, or $\log(Pm)/L$.
The last has bounded variation on the allowed prime interval.
If $(m,dk)>1$, both relevant counts vanish,
or the coprimality with $d$ removes the term.
Otherwise $P\equiv l\bar m(k)$.
Apply prime SW and partial summation, uniformly at endpoints $P\ge y_1$.
The endpoint errors sum to $O_{A_{\rm SW}}(TL^{-A_{\rm SW}})$, since
$\log y_1=L^{2/3}$ and $\sum_{m\le2T/y_1}1/m\ll L$.
Excluding $P\mid dk$ costs at most
$O((\omega(d)+\omega(k))T/y_1)\ll_{A_{\rm SW}} T\tau(d)L^{-A_{\rm SW}}$.
Endpoint subtraction handles arbitrary $I$. This proves (3.25).
All errors are absolute, of size $O(TL^{-A_{\rm SW}})$.
Only the classical Siegel--Walfisz theorem is used for primes.
The proof is uniform for every $\zeta\ge z_0$.
A nonzero selected coefficient has all prime factors at least $\zeta$.
Raise the lower endpoint of the sum over the largest prime variable.
Its SW bound is unchanged. If $\zeta>2T$, the coefficient is zero.

\subsection{First dispersion with restoration before CRT}

Work first with the normalized convolution in (3.5).
Split all three factors dyadically.
Apply the triangle inequality before squaring.
There are $O(L^3)$ triples, including the outer factor. For one triple put

\[
\lambda_r=\sum_{r_3r_2=r}\gamma_3(r_3)\gamma_2(r_2),
\quad r_3\asymp R,\ r_2\asymp S.
\tag{3.26}
\]

Since $r$ is $u_0$-rough and at most $X^{0.6}$,
$|\lambda_r|\le\tau(r)=O_\sigma(1)$. The arbitrary small-prime factor
of C1 remains in the outer $q$, as the only smooth divisor.
Let $\psi_0$ be a smooth nonnegative bump, equal to 1 on $[1,2]$,
supported on $[1/2,5/2]$. Cauchy--Schwarz in $(q,m)$ gives

\[
\begin{aligned}
&|S|^2\ll QM_IL^c S',\quad\\
&S'=\sum_q\psi_0(q/Q)1_{(q,a)=1}\sum_m\psi_I(m)
 \left|\sum_r\lambda_r\sum_n\alpha_n\mathcal S_{mn}(qr;a)\right|^2.
\end{aligned}
\tag{3.27}
\]

Write $t=(q,r_1,r_2,n_1,n_2)$, $k_i=mn_i-a$, and

\[
c(t)=\psi_0(q/Q)1_{(q,a)=1}\lambda_{r_1}\bar\lambda_{r_2}
                       \alpha_{n_1}\bar\alpha_{n_2},
\quad \omega(t)=1_{(n_1,qr_1)=(n_2,qr_2)=1}.
\tag{3.28}
\]

Expanding the square gives $S'=S_1-2\Re S_2+S_3$, where

\[
\begin{aligned}
S_3&=\sum_t c(t)\sum_m\psi_I(m)1_{qr_1\mid k_1}1_{qr_2\mid k_2},\\
S_2&=\sum_t c(t)\sum_m\psi_I(m)
       1_{qr_1\mid k_1}1_{(mn_2,qr_2)=1}/\varphi(qr_2),\\
S_1&=\sum_t c(t)\sum_m\psi_I(m)
       1_{(mn_1,qr_1)=1}1_{(mn_2,qr_2)=1}/(\varphi(qr_1)\varphi(qr_2)).
\end{aligned}
\tag{3.29}
\]

Let $\mathcal T$ be the tuple set. Let $\mathfrak B$ be its union
of cross-copy gcd failures

\[
(n_1,n_2)>1,\ (r_1,r_2)>1,\ (n_1,r_2)>1,
\ (n_2,r_1)>1,\ (q,r_1r_2)>1,
\tag{3.30}
\]

and let $\mathfrak D=\{|n_1-n_2|<TL^{-C_M}\}$.
Set $\mathfrak G'=\mathcal T\setminus(\mathfrak B\cup\mathfrak D)$.
On every tuple subset $V$ define the formal sums

\[
\begin{aligned}
MT(V)&=\sum_{t\in V}c(t)\omega(t)
                \frac{M_I\widehat\phi_0(0)}{qr_1r_2\varphi(q)},\\
\mathcal Z(V)&=\sum_{t\in V}c(t)\omega(t)1_{n_1\equiv n_2(q)}
                \frac{M_I\widehat\phi_0(0)}{qr_1r_2}.
\end{aligned}
\tag{3.31}
\]

\textbf{Restoration lemma.} Choose $A_E,C_\Sigma,C_M$ in the order of \S\ref{sec:pi-global}. Then

\[
S'=M_IE_2^*+\mathcal Z(\mathcal T)-MT(\mathcal T)+O(E_{\rm GR}),
\tag{3.32}
\]
\[
\begin{aligned}
E_{\rm GR}\ll{}&L^{3C_\Sigma+c}M_IN_I\left(\frac{N_I}{z_0Q}+1\right)
 +C_\epsilon X^\epsilon M_IN_I\left(\frac{N_I}{u_0Q}+1\right)\\
&+\frac{M_IN_I^2}{QL^{A_E}}.
\end{aligned}
\tag{3.33}
\]

We first prove the positive estimates needed for CRT.
For $S_i$, let $\Upsilon_i(V)$ be the positive majorant.
It replaces $c(t)$ by $|c(t)|$ and deletes coprimalities.
Retain the two congruences for $\Upsilon_3$, one for $\Upsilon_2$.
Let $|MT|(V),|\mathcal Z|(V)$ be the positive versions of (3.31).
We write $\Upsilon_{\rm exc}$ for the positive contribution of
$(S_3)$ from the complement of $\Sigma^*$ defined in (3.37).

The total-variation progression estimate is

\[
\left|\sum_{m\equiv b(k)}f(m)-k^{-1}\int f\right|\le V(f).
\tag{3.34}
\]

Compare $f(b+jk)$ to its cell average on
$[b+jk,b+(j+1)k]$. Bound the difference by the cell variation,
and sum. Möbius expansion of the coprimality then gives
$\sum_{(m,k)=1}\psi_I=(\varphi(k)/k)M_I\widehat\phi_0(0)
+O(\tau(k))$.

From $n/\varphi(n)=\sum_{d\mid n}\mu^2(d)/\varphi(d)$ we get

\[
\sum_{r\le Y}\frac1{\varphi(r)}\ll\log(2Y),\quad
\sum_{r\le Y,p\mid r}\frac1{\varphi(r)}\ll\frac{\log(2Y)}{p-1},
\]
\[
\sum_{q\asymp Q}\frac1{\varphi(q)^2}\ll L^2/Q,\qquad
\sum_{q\asymp Q,p\mid q}\frac1{\varphi(q)^2}\ll L^2/(pQ).
\tag{3.35}
\]

For $\Upsilon_1,\Upsilon_2,|MT|$ the weight is bounded by
$O_\sigma(M_I/(\varphi(q)^2\varphi(r_1)\varphi(r_2)))$, since
$qr_1\ll QRS<M_I$. For $|\mathcal Z|$ it is
$O_\sigma(M_I/(qr_1r_2))1_{n_1\equiv n_2(q)}\omega(t)$.
These two shapes prove

\[
\Upsilon_1(\mathfrak B)+\Upsilon_2(\mathfrak B)+|MT|(\mathfrak B)+|\mathcal Z|(\mathfrak B)
 \ll L^cM_IN_I\left(\frac{N_I}{z_0Q}+1\right).
\tag{3.36}
\]

For the first failure of (3.30), fix $n_1$.
Choose one of its at most $L$ primes $p\ge z_0$.
Count $n_2$ in multiples of $p$. The first shape gives $N_I/z_0+1$.
For the weight of $|\mathcal Z|$, $(n_1,q)=1$.
So the extra congruence gives a class modulo $pq$,
with $N_I/(pq)+1$ choices.
For a common prime $p\ge u_0$ of $r_1,r_2$, fix $r_1$.
Sum the $p$-divisible $r_2$ weight using (3.35).
For $p\mid(n_1,r_2)$, fix $n_1$ and do the same weighted $r_2$ sum.
For $p\mid(n_2,r_1)$, fix $r_1$.
Then count $n_2$ in multiples of $p$, or modulo $pq$ for the weight of $|\mathcal Z|$.
For $p\mid(q,r_i)$, use the last estimate of (3.35).
This covers all five failures, and $u_0\ge z_0$ gives (3.36).

Split the unrestricted $(t,m)$-sum at

\[
\Sigma^*=\{\tau(|k_1|)\tau(|k_2|)<L^{C_\Sigma}\}.
\tag{3.37}
\]

The complement is estimated on all tuples, coprime or not.
Put $u=\tau(|k_1|)$ and $v=\tau(|k_2|)$.
The exceptional indicator is at most $uvL^{-C_\Sigma}$.
The two divisor choices $r_i\mid k_i$ cost at most $uv$.
There are also at most $u$ choices for $q\mid k_1$.
The resulting weight is $u^3v^2$, of degree five.
To keep the progression restriction while counting, split into $u\ge v$
and $v>u$.
On the first part $u^3v^2\le u^5$.
For each fixed $(m,n_1)$, sum over the at most $u$ divisors $q$ first.
The condition $q\mid k_1$ and $(q,a)=1$ implies $(m,q)=1$.
Thus $q\mid k_2$ forces $n_2\equiv n_1\pmod q$.
There are at most $4N_I/Q$ such $n_2$, by (3.6).
This bounds the first part by $L^{-C_\Sigma}(4N_I/Q)\sum u^5$.
Exchanging the two copies $(n_1,r_1)\leftrightarrow(n_2,r_2)$ gives the same bound for the second part.
Consequently

\[
\Upsilon_{\rm exc}\ll L^{-C_\Sigma}\frac{N_I}{Q}
 \sum_{n\in I_T,m\in\widehat I_M}\tau(|mn-a|)^5
 \ll \frac{M_IN_I^2}{Q}L^{c_1-C_\Sigma}.
\tag{3.38}
\]
The factor $4$ from $4N_I/Q$ and the factor $2$ from the two parts are
absorbed in the implied constant.

For a modulus-related failure, fix $(m,n_1)$. The number of
$(q,r_1)$ with $qr_1\mid k_1\ne0$ is
$\tau_3(|k_1|)\le\tau(|k_1|)^2\ll_\epsilon X^{\epsilon/3}$.
A common prime $p\ge u_0$ comes from $r_1$, $n_1$, or $q$.
It forces $n_2$ into one class modulo $pq$.
For $p\mid r_1,r_2$ use $p\nmid m$.
For $p\mid n_1,r_2$ use $mn_2\equiv a(p)$.
For $p\mid n_2,r_1$ use $n_2\equiv0(p)$.
If $p\mid q,r_2$, then $pq\mid qr_2\mid k_2$ gives the class directly.
If $p\mid q,r_1$ only, exchange the copies.
There are $O(L)$ candidate primes and $O(N_I/(u_0Q)+1)$ choices in each class.
There are at most $\tau(|k_2|)$ remaining choices of $r_2$. Hence

\[
\Upsilon_3(\mathfrak B_{\rm modulus})
 \ll_\epsilon X^\epsilon M_IN_I(N_I/(u_0Q)+1).
\tag{3.39}
\]

For $(n_1,n_2)>1$ on $\Sigma^*$, there are at most $L^{2C_\Sigma}$
pairs $(q,r_1)$. There are at most $L$ primes $p\mid n_1$, $p\ge z_0$.
There are at most $2N_I/(z_0Q)+1$ choices for $n_2$.
There are at most $L^{C_\Sigma}$ choices for $r_2$. Thus

\[
\Upsilon_3(\mathfrak B_n\cap\Sigma^*)
 \ll L^{3C_\Sigma+1}M_IN_I(N_I/(z_0Q)+1).
\tag{3.40}
\]

On the near-diagonal the same count replaces $N_I/z_0$ by
$TL^{-C_M}$. The density shapes and (3.35) give the same or smaller
bound. Therefore

\[
\begin{aligned}
&\Upsilon_3(\mathfrak D\cap\Sigma^*)+\Upsilon_1(\mathfrak D)+\Upsilon_2(\mathfrak D)
 +|MT|(\mathfrak D)+|\mathcal Z|(\mathfrak D)\\
&\qquad\ll L^{3C_\Sigma+c}M_IN_I(TL^{-C_M}/Q+1).
\end{aligned}
\tag{3.41}
\]

Use the common constants from \S\ref{sec:pi-global}:
$C_\Sigma=A_E+2c_*+3$ and
$C_M=5C_\Sigma+A_E+C_{\rm box}+c_*+3$.
They satisfy $C_\Sigma>A_E+c_1+2$ and
$C_M>3C_\Sigma+C_{\rm box}+A_E+c$.
To check this, divide by the last term of (3.33). After this division, the right side
of (3.41) is at most
\[
L^{3C_\Sigma+c+A_E}
\left(\frac{T L^{-C_M}}{N_I}+\frac{Q}{N_I}\right)
\ll L^{3C_\Sigma+c+A_E+C_{\rm box}-C_M}
   +L^{3C_\Sigma+c+A_E}X^{-\sigma/3}=o(1).
\]
Here LEN supplies $T/N_I\ll L^{C_{\rm box}}$, while (3.6)
supplies $Q/N_I\le X^{-\sigma/3}$. The chosen $C_M$ makes
the first logarithmic exponent negative; the second term tends to zero
for fixed $\sigma$. The same target absorbs (3.38), since
$C_\Sigma>A_E+c_1+2$.
This removes bad tuples and the near-diagonal from \textbf{all} of
$S_1,S_2,S_3,Z,MT$ before their main terms are compared.

Now use CRT on $\mathfrak G'$. Let
$Q_1=qr_1,Q_2=qr_2,g=(r_1,r_2)$.
The two congruences have one class modulo $[Q_1,Q_2]=qr_1r_2/g$ precisely when
$(n_1,Q_1)=(n_2,Q_2)=1$ and $n_1\equiv n_2(qg)$.
To see this, equate $a\bar n_1$ and $a\bar n_2$ modulo $(Q_1,Q_2)=qg$.
On $\mathfrak G'$, $g=1$ and $q,r_1,r_2$ are pairwise coprime.

\textbf{Fixed cutoff.} On the smoothed support put

\[
k(t)=qr_1r_2\le K_*=40QR^2S^2,\quad
H_2=QR^2S^2X^\eta/M_I,\quad H_*=40H_2.
\tag{3.42}
\]

These are fixed for the entire shell triple. Poisson gives exactly

\[
\begin{aligned}
E_2^*={}&\sum_{t\in\mathfrak G'}
 \frac{c(t)\omega(t)1_{n_1\equiv n_2(q)}}{k(t)}\\
&\times\sum_{1\le|h|\le H_*}\widehat\phi_0(hM_I/k(t))
                      e(h(b(t)-m_0)/k(t)),\\
S_3(\mathfrak G')={}&\mathcal Z(\mathfrak G')+M_IE_2^*+O(X^{-90}).
\end{aligned}
\tag{3.43}
\]

Put $\lambda=M_I/k(t)$. This is the Poisson prefactor of one
progression sum. In (3.43) it appears as the factor $M_I$ in front of
$E_2^*$ and the factor $1/k(t)$ inside $E_2^*$.
Integration by parts gives
$|\widehat\phi_0(u)|\le C_j|u|^{-j}$.

Here one may take $C_j=(2\pi)^{-j}\|\phi_0^{(j)}\|_1$.
It is independent of the box and the tuple.
The integer $j=\lceil104/\eta\rceil+1$ is fixed once $\sigma$ is fixed.
Choose $X_0$ after these seminorms and all fixed coefficient parameters.
So the tail bound is uniform over the whole box family.
If $H_*<1$, then $\lambda>X^\eta$, since $\lambda H_*\ge X^\eta$.

The tail of one progression
sum is at most $C'_j(\lambda H_*)^{1-j}$ if $H_*\ge1$, and at
most $C'_j\lambda^{1-j}$ otherwise. In either case it is
$O(X^{-\eta(j-1)})$, since $\lambda H_*\ge X^\eta$.
Take $j=\lceil104/\eta\rceil+1$. After increasing $X_0$, the
weighted tail per tuple is at most $X^{-103}$. There are at most
$X^5$ tuples. This proves (3.43).
The cap (3.42) is the same for all tuples.
So it survives the Mellin separation of the inner variable below.

On $\mathfrak G'$, (3.34) gives

\[
S_1(\mathfrak G')=MT(\mathfrak G')+O(N_I^2L^c/Q).
\tag{3.44}
\]

For $S_2$, retain $(m,r_2)=1$. Möbius expansion over $e\mid r_2$
gives one class modulo $qr_1e$, so

\[
\sum_m\psi_I(m)1_{mn_1\equiv a(qr_1)}1_{(m,r_2)=1}
 =\frac{M_I\widehat\phi_0(0)\varphi(r_2)}{qr_1r_2}
   +O(\tau(r_2)).
\tag{3.45}
\]

Divide by $\varphi(qr_2)=\varphi(q)\varphi(r_2)$.
The main term is again $MT$.
The summed error is $O(RSN_I^2L^c)$.
It is absorbed using $QRS<M_IX^{-89\sigma}$.
Exchanging the two copies shows that $MT(\mathfrak G')$ is real.
Complete $\mathcal Z$ and $MT$ formally to $\mathcal T$, with errors bounded by (3.36), (3.41).
This proves (3.32).
CRT and the zero-mode identity are used on $\mathfrak G'$ only.

\subsection{Main-term cancellation and the full GR error}

Put $\alpha^{(r)}_n=\alpha_n1_{(n,r)=1}$ and

\[
\Delta(\xi;q,b)=\sum_{n\equiv b(q)}\xi_n-
                  \varphi(q)^{-1}\sum_{(n,q)=1}\xi_n.
\tag{3.46}
\]

Expanding the product of these two centered counts gives exactly

\[
\begin{aligned}
\mathcal Z(\mathcal T)-MT(\mathcal T)
={}&\sum_{q,r_1,r_2}\frac{\psi_0(q/Q)1_{(q,a)=1}
               \lambda_{r_1}\bar\lambda_{r_2}M_I\widehat\phi_0(0)}{qr_1r_2}\\
&\qquad\sum_{b\bmod q}^{*}\Delta(\alpha^{(r_1)};q,b)
                  \overline{\Delta(\alpha^{(r_2)};q,b)}.
\end{aligned}
\tag{3.47}
\]

For the norm-normalized condition \cite[(A2)]{BFI86}, apply (3.25) to $dr$, with
$\tau(dr)\le\tau(d)O_\sigma(1)$. If
$\|\alpha^{(r)}\|_2^2\le N_IL^{-A_3-1}$, the elementary estimate

\[
\sum_{q\le3Q}\sum_b^*|\Delta(\alpha^{(r)};q,b)|^2
 \ll (T\log(3Q)+3Q)\|\alpha^{(r)}\|_2^2
 \ll N_ITL^{-A_3}
\tag{3.48}
\]

\noindent suffices, since $\log(3Q)\le L$ is absorbed by the factor $L^{-1}$.
If $\|\alpha^{(r)}\|_2^2>N_IL^{-A_3-1}$, then $\|\alpha^{(r)}\|_2(2T)^{1/2}\gg TL^{-(C_{\rm box}+A_3+1)/2}$.
Then (3.25) with exponent
$A_{\rm BFI}+(C_{\rm box}+A_3+1)/2$ implies \cite[(A2)]{BFI86}.
The condition $Q<TX^{-29\sigma}$ implies the condition
$3Q\le T(\log T)^{-B_1}$ of \cite[Theorem~0(a)]{BFI86} eventually,
which then gives (3.48) in this case too.
Apply $|uv|\le(|u|^2+|v|^2)/2$ in (3.47), and sum
the reciprocal $r_i$ weights. This gives

\[
|\mathcal Z(\mathcal T)-MT(\mathcal T)|
 \ll M_IN_I^2L^{C_{\rm box}+2-A_3}/Q.
\tag{3.49}
\]

Take $A_3=A_E+C_{\rm box}+2$.
Increase it by a fixed constant to absorb any endpoint logarithm.
Combine (3.27), (3.32), (3.33), (3.49), with $\epsilon=\ell_0\sigma^2/8$.
For a requested local precision $A'>0$, set $A_{\rm work}=A'+4$
before choosing $A_E,C_\Sigma,C_M$ as in \S\ref{sec:pi-global}.
The resulting restoration contributions to $|S|$ are at most

\[
XL^{-A_{\rm work}}+XL^Cz_0^{-1/2}
 +XL^C(Q/N_I)^{1/2}
 +XX^{\epsilon/2}L^C\{u_0^{-1/2}+(Q/N_I)^{1/2}\}.
\tag{3.50}
\]

Both terms in braces are kept.
Their power savings are at least $7\ell_0\sigma^2/16$ and
$\sigma/6-\ell_0\sigma^2/16$, respectively.
Each exceeds $c_\sigma$, and $z_0^{-1/2}$ beats every logarithm.
More explicitly, for fixed $A_{\rm work}$ and sufficiently large $X$,
$XL^Cz_0^{-1/2}\ll XL^{-A_{\rm work}}$, because
$z_0^{-1/2}=X^{-1/(2(\log\log X)^3)}$ beats every fixed logarithm.
Thus this term joins the first term in (3.51), rather than the
fixed-power-saving term.
The separate term $XL^C(Q/N_I)^{1/2}$ in (3.50) has saving
$\sigma/6>c_\sigma$, by $N_I\ge QX^{\sigma/3}$ in (3.6),
and is therefore absorbed in the displayed power-saving term.
Increase $X_0$ after the fixed exponents. This gives

\[
|S|\ll XL^{-A_{\rm work}}
      +M_IQ^{1/2}L^c|E_2^*|^{1/2}+X^{1-c_\sigma}L^C.
\tag{3.51}
\]

This completes the first dispersion.
A later saving for $E_2^*$ changes only the middle term.

\subsection{Reciprocity, the complete mask, and the second square}

Split the combined products again as $r_i\leftarrow r_is_i$, where
$r_i\asymp R,s_i\asymp S$. The remaining support is

\[
(q,r_1s_1r_2s_2)=(r_1s_1,r_2s_2)=(n_1,n_2)=1,
\]
\[
(n_1,r_2s_2)=(n_2,r_1s_1)=1,
\quad(n_i,qr_is_i)=1,
\quad |n_1-n_2|\ge TL^{-C_M}.
\tag{3.52}
\]

Solving the two CRT congruences and using
$\bar u/v+\bar v/u\equiv1/(uv)\pmod1$ for $(u,v)=1$ gives

\[
\frac{ah\overline{n_1r_2s_2}}{qr_1s_1}
 +\frac{ah\overline{n_2qr_1s_1}}{r_2s_2}
 \equiv\frac{ah(n_1-n_2)\overline{n_2qr_1s_1}}{n_1r_2s_2}
       +\frac{ah}{qr_1s_1r_2s_2n_1}\pmod1.
\tag{3.53}
\]

All inverses exist by (3.52). Put $f=(n_1-n_2)/q\ne0$.
Then the outer tuple $\omega=(f,n_1,n_2,r_1,r_2,s_2)$ determines $q=(n_1-n_2)/f$.
The modular phase becomes

\[
e\left(\frac{ahf\overline{n_2r_1s_1}}{n_1r_2s_2}\right).
\tag{3.54}
\]

Split $f\ne0$ by sign and dyadic intervals
$f\in\varsigma_f(F,2F]$, starting at $F=1/2$. Split $h\ne0$ by
sign and intervals $H_0<|h|\le\min(2H_0,H_*)$, starting at
$H_0=1/2$. These disjoint intervals include $f,h=\pm1$ exactly
once. Write $\varsigma_f,\varsigma_h\in\{\pm1\}$ for the signs of $f$ and $h$.
The cap on the last, shortened $h$-block depends on $h$ alone.
Always $F\ll T/Q$.

The nonmodular factor is exactly

\[
\widehat\phi_0\left(\frac{hM_I}{qr_1s_1r_2s_2}\right)
 e\left[\frac{h}{qr_1s_1r_2s_2}
                   \left(\frac a{n_1}-m_0\right)\right].
\tag{3.55}
\]

Its dependence on $(h,s_1)$ is only through $|h|/s_1$. For
$v=|h|S/(H_0s_1)\in[1/2,2]$, choose a fixed smooth cutoff.
It equals 1 there and is supported in $(1/4,4)$. Rewrite (3.55):

\[
\begin{gathered}
W_{Y(\omega)}(v)=w(v)\widehat\phi_0(Y_2v)e(Y_1v),\\
Y_1=\frac{\varsigma_hH_0(a/n_1-m_0)}{Sqr_1r_2s_2},\qquad
Y_2=\frac{\varsigma_hH_0M_I}{Sqr_1r_2s_2}.
\end{gathered}
\tag{3.56}
\]

Since $H_0\le H_*=40H_2$,

\[
|Y_1|+|Y_2|\ll (1+a/X)L^{C_{\rm box}}X^\eta
 \ll L^{C_{\rm box}+K_{\rm scale}}X^\eta=:L_{\rm osc}.
\tag{3.57}
\]

Mellin inversion separates the ratio. Put
$m(t)=\sup_\omega|\widetilde W_{Y(\omega)}(it)|$.
The zeroth and second logarithmic derivative estimates give

\[
m(t)\ll\min\{1,L_{\rm osc}^2/(1+|t|)^2\},\qquad
\int_\mathbb Rm(t)\,dt\ll L_{\rm osc}.
\tag{3.58}
\]

If $m(t)>0$, divide the outer coefficient by $m(t)$.
The norm bounds below are for this normalized coefficient.
The factor $m(t)$ is restored before the $t$-integration.
If $m(t)=0$, the integrand is zero.
So the integrable decay of $m$ is kept.

The full inner coprimality mask is

\[
(s_1,ac)=1,\qquad(s_1,J_\omega)=1,
\quad c=n_1r_2s_2,\quad J_\omega=qn_2.
\tag{3.59}
\]

The second condition combines $(s_1,q)=1$ and $(s_1,n_2)=1$.
To separate the second, split $\mathcal P_X=[u_0,X^3]\cap\mathbb P$ at random.
Each prime goes independently and uniformly into one of two classes $B_\chi,R_\chi$.
For an integer $m$, let
$\mathcal P_X(m)=\{p\in\mathcal P_X:p\mid m\}$ and
$\nu_X(m)=|\mathcal P_X(m)|$. If all prime factors of $s$ lie in
$\mathcal P_X$, then exactly

\[
1_{(s,J)=1}=\mathbb E_\chi
 [2^{\nu_X(s)+\nu_X(J)}
  1_{\mathcal P_X(s)\subseteq B_\chi}
  1_{\mathcal P_X(J)\subseteq R_\chi}].
\tag{3.60}
\]

A common prime makes the event impossible.
Otherwise the event has probability $2^{-\nu_X(s)-\nu_X(J)}$.
Every prime of a supported $s_1$ lies in $\mathcal P_X$.
So $(s_1,J_\omega)=1$ depends only on the primes of $J_\omega$ in $\mathcal P_X$.
Since $s_1\le X^{0.6}$ and $J_\omega\le X^{0.94}$,
one may take

\[
K_s=\lceil0.6/(\ell_0\sigma^2)\rceil,\qquad
K_J=\lceil1/(\ell_0\sigma^2)\rceil.
\tag{3.61}
\]

Replace the inner coefficient by
$\lambda_s^\chi=\lambda_s2^{\nu_X(s)-K_s}
1_{\mathcal P_X(s)\subseteq B_\chi}$
and each outer representation coefficient by
$A_\omega^\chi=A_\omega2^{\nu_X(J_\omega)-K_J}
1_{\mathcal P_X(J_\omega)\subseteq R_\chi}$.
Both multipliers have modulus at most 1.
The original sum is $2^{K_s+K_J}\mathbb E_\chi W_\chi$.
This factor is $O_\sigma(1)$, and $\mathbb E_\chi$ is an average.
Fix $\chi$ before aggregating representations.
The partition acts on $\lambda_s$ and $A_\omega$ only, so (3.25) for $\alpha_n$ holds for each $\chi$.

Set $b=n_2r_1,c=n_1r_2s_2$, and

\[
I(b,c,f)=\sum_{s\asymp S,(s,ac)=1}\sum_h
             \lambda'_{s,t}\eta_{h,t}
             e(ahf\overline{bs}/c),
\tag{3.62}
\]

where $\lambda'_{s,t}=(S/s)\lambda_s^\chi(s/S)^{it}$, and
$\eta_{h,t}=(|h|/H_0)^{-it}1_{|h|\le H_*}$ on the chosen signed
block. Both are 1-bounded. Dyadically localize $b\asymp D,c\asymp C$.
The outer support has $(b,c)=(b,p_0)=1$. So the nonnegative majorant

\[
\rho=\psi_0(b/D)\psi_0(c/C)1_{(b,c)=(b,p_0)=1}
                                   1_{f\in\varsigma_f(F,2F]}
\tag{3.63}
\]

\noindent equals 1 on it. Cauchy--Schwarz now gives

\[
|W_{\chi,t}|\le\left(\sum_{b,c,f}|A_{b,c,f}^\chi|^2\right)^{1/2}
               W_4^{1/2},\quad W_4=\sum_{b,c,f}\rho|I|^2.
\tag{3.64}
\]

Fix $(b,c,f)$. Factor $b=n_2r_1,c=n_1r_2s_2$; then $q$ is determined.
So there are at most $\tau(b)\tau_3(c)$ outer representations.
For fixed $n_1,n_2$ the permitted $q\mid n_1-n_2$ have divisor multiplicity.
So there are $O(T^2R^2S)X^{o(1)}$ representations in total.
A normalized representation coefficient
has modulus $O_\sigma((QR^2S^2)^{-1})$. Summing over representations with Cauchy--Schwarz,

\[
\left(\sum|A^\chi|^2\right)^{1/2}
 \ll X^{o(1)}\frac{T}{QRS^{3/2}}.
\tag{3.65}
\]

This bound holds for the coefficients $A^\chi$ of each class $\chi$.

Expanding $|I|^2$, put

\[
z=ss',\quad k=h_1s'-h_2s,\quad y=fk.
\tag{3.66}
\]

The terms with $k=0$ have at most $X^{o(1)}H_*S$ choices of
$(s,s',h_1,h_2)$. To see this, write $(s,s')=g(u,v)$, $(u,v)=1$, and
$(h_1,h_2)=j(u,v)$. Sum $(g,j)$ over their ranges. Thus

\[
|W_{k=0}|\ll X^{o(1)}H_*S DFC
 \ll (T^2R^2S^3)X^{-8\sigma+\eta+o(1)},
\tag{3.67}
\]

\noindent where the final inequality is $T^2R^2S<X^{1-8\sigma}$.
For $k\ne0$, split $y=\varepsilon_y m$, $m>0$, and $z$ into
$Z\in\{S^2,2S^2\}$. The coefficient independent of $(b,c)$ is

\[
\begin{aligned}
\beta^{\chi,\varepsilon_y}_{z,m}={}&1_{(z,ap_0)=1}
 \sum_{ss'=z}\lambda'_{s,t}\bar\lambda'_{s',t}
 \sum_{\substack{f\in\varsigma_f(F,2F],\ h_i\ {\rm in\ block}\\
                  f(h_1s'-h_2s)=\varepsilon_y m}}
                     \eta_{h_1,t}\bar\eta_{h_2,t}.
\end{aligned}
\tag{3.68}
\]

Put $Y:=H_*SF$. By (3.70), $|y|\le8FH_*S$.
So the coefficient is supported on $1\le m\le8Y$.
The off-diagonal sum is exactly

\[
\sum_{z,m}\beta_{z,m}^{\chi,\varepsilon_y}
 \sum_{b,c:(b,c)=(b,p_0)=(z,c)=1}
      \psi_0(b/D)\psi_0(c/C)e(\varepsilon_y a m\overline{zb}/c).
\tag{3.69}
\]

\subsection{The W-prime norm, with endpoints}

For $S\ge1$, $1\le H\le S$, $F\ge1/2$, and $a\ne0$, let

\[
\begin{aligned}
B_{z,y}=\#\{&s_i\in(S,2S],\ (s_1s_2,a)=1,\\
&1\le|h_i|\le H,\ f\in(F,2F]:\\
&s_1s_2=z,\ f(h_1s_2-h_2s_1)=y\ne0\}.
\end{aligned}
\tag{3.70}
\]

For every nonzero integer $v$, define its $a$-supported part by
$v_a=\prod_{p\mid a}p^{v_p(|v|)}$.
In particular $y_a=\prod_{p\mid a}p^{v_p(|y|)}$.
Set

\[
\begin{gathered}
\mathcal T_0=\max\tau(z)\tau(|y|),\qquad
P_0=\prod_{p\mid a}(1-p^{\vartheta_{\rm sp}-1})^{-1},\\
P_1=P_0^2,\qquad P_2=\prod_{p\mid a}(1-p^{\vartheta_{\rm sp}-1/2})^{-1},
\end{gathered}
\tag{3.71}
\]

where the maximum is over $z\le4S^2,0<|y|\le8FHS$.
Then

\[
\sum_{z,y\ne0}B_{z,y}^2(y_a)^{\vartheta_{\rm sp}}
 \le\mathcal T_0\,2FP_0H^2S^2
             [16P_1+64(8+\log(2S))P_0+32P_2].
\tag{3.72}
\]

\emph{Proof.} For fixed $s=(s_1,s_2)$, let
$N_s(k)=\#\{h:1\le|h_i|\le H,h_1s_2-h_2s_1=k\}$.
Each $(z,y)$ has at most $\tau(z)\tau(|y|)$ choices of $(s,f)$.
Therefore Cauchy--Schwarz and $(fk)_a=f_ak_a$ reduce the left side to

\[
\mathcal T_0\left(\sum_{f\le2F}f_a^{\vartheta_{\rm sp}}\right)Q_0,
\qquad Q_0=\sum_s\sum_{k\ne0}k_a^{\vartheta_{\rm sp}} N_s(k)^2,
\quad \sum_{f\le2F}f_a^{\vartheta_{\rm sp}}\le2FP_0.
\tag{3.73}
\]

For the last inequality, majorize $f_a^{\vartheta_{\rm sp}}$ by
$\sum_{d\mid a^\infty,d\mid f}d^{\vartheta_{\rm sp}}$.
Then sum $\lfloor2F/d\rfloor\le2F/d$. It holds for every $F>0$.

Put $g=(s_1,s_2)$, $s=g(u,v)$, $S'=S/g$.
A fibre of $h_1s_2-h_2s_1$ has at most $1+2H/S'$ points.
It has at most one when $S'\ge2H$. Expand

\[
k_a^{\vartheta_{\rm sp}}=\sum_{d\mid a^\infty,d\mid k}\varphi_{\vartheta_{\rm sp}}(d),
\quad\varphi_{\vartheta_{\rm sp}}(p^j)=p^{j\vartheta_{\rm sp}}(1-p^{-\vartheta_{\rm sp}}),
\quad0\le\varphi_{\vartheta_{\rm sp}}(d)\le d^{\vartheta_{\rm sp}}.
\tag{3.74}
\]
Here $\varphi_{\vartheta_{\rm sp}}$ is multiplicative with
$\varphi_{\vartheta_{\rm sp}}(1)=1$; the prime-power formula is for $j\ge1$.

If $A_d(s)$ counts $h\in[-H,H]^2$ with $k\ne0,d\mid k$, then

\[
Q_0\le\sum_{d\mid a^\infty,d\le4HS}d^{\vartheta_{\rm sp}}
 \left[\sum_s A_d(s)+\sum_{s:S'<2H}\frac{2H}{S'}A_d(s)\right].
\tag{3.75}
\]

For the first sum write $h=e w$, $w$ primitive,
$\|w\|_\infty=j\le H$. There are at most $8j$ choices of $w$,
and $e\le H/j$. The congruence is
$d_e\mid w_1s_2-w_2s_1$, $d_e=d/(d,e)$.
Since $(s_1s_2,d_e)=1$, a contributing primitive vector has
$(w_1,d_e)=1$. The case $w_1=0$ forces $d_e=1$. There are at
most $2S$ choices for $s_1$, then $S/d_e+1$ for $s_2$.
By $\sum_{e\le E}(d,e)\le\tau(d)E$, this proves

\[
\sum_sA_d(s)\le16\tau(d)H^2S^2/d+16H^2S.
\tag{3.76}
\]

For the second sum $(d,g)=1$.
At most $8HS'/d$ nonzero levels of $h_1v-h_2u$ are divisible by $d$.
Each has at most $4H/S'$ points. Thus $A_d(s)\le32H^2/d$.
For each $g$, sum the at most $(S/g+1)^2$ pairs.
Use the exact inequality

\[
\sum_{g\le2S}\frac{(S/g+1)^2}{S/g}
 =\sum_{g\le2S}(S/g+2+g/S)\le S(8+\log(2S)).
\tag{3.77}
\]

So before the $d$-sum, the second term of (3.75) is at most
$64(8+\log(2S))H^3S/d\le64(8+\log(2S))H^2S^2/d$, using $H\le S$. Finally

\[
\begin{aligned}
\sum_{d\mid a^\infty}\tau(d)d^{\vartheta_{\rm sp}-1}&=P_1,\\
\sum_{d\mid a^\infty}d^{\vartheta_{\rm sp}-1}&=P_0,\\
\sum_{d\le4HS,d\mid a^\infty}d^{\vartheta_{\rm sp}}&\le(4HS)^{1/2}P_2.
\end{aligned}
\tag{3.78}
\]

The last contribution of (3.76) is at most $32H^2S^2P_2$, since
$H\le S$. This proves (3.72), also for real nonintegral $(S,H)$.

Uniformly for $|a|\le X^2$, elementary prime ordering gives

\[
\begin{gathered}
P_0,P_1,P_2\le\exp(C(\log X)^{23/32})=X^{o(1)},\\
\mathcal T_0\le\exp(C\log X/\log\log X)=X^{o(1)}.
\end{gathered}
\tag{3.79}
\]

For $\alpha\in\{1-\vartheta_{\rm sp},1/2-\vartheta_{\rm sp}\}$ we have
$\sum_{p\mid a}p^{-\alpha}\le
\sum_{j\le\omega(a)}j^{-\alpha}\ll(\log X)^{1-\alpha}$.
The bound for the products then follows from
$-\log(1-v)\ll v$. The divisor bound is elementary.

To apply (3.72) to (3.68), use the cap
$H=\min(2H_0,H_*)$. Empty blocks have zero coefficient.
Increase $X_0$ in (3.6)--(3.7). Then

\[
H_2/S\ll X^{-9\sigma+\eta}L^{C_{\rm box}},\qquad H_*\le S.
\tag{3.80}
\]

The strict power margin in (3.6)--(3.7)
absorbs the fixed logarithmic factor and the factor 40.
For $S=1/2$ this gives
$H_*=40H_2\ll X^{-9\sigma+\eta}L^{C_{\rm box}}\to0$.
If the middle factor is supported at $s=1\in(1/2,1]$,
set $S=1/2$. Then $H_*<1$, every nonzero-frequency block is empty,
and $E_2^*=0$. Thus Theorem 2.8 is needed only for $S\ge1$;
in that case a nonempty block has $H\ge1$.
Each summand of (3.68) has modulus at most 1, so
$|\beta_{z,m}|\le B_{z,\varepsilon_y\varsigma_fm}$.
Since $p_0\nmid a$, $(p_0m)_a=m_a$ exactly. Thus

\[
\sum_{n''\mid a^\infty}(n'')^{\vartheta_{\rm sp}}
                 \|\beta^\flat(n'')\|_2^2
 \ll X^{o(1)}(H_*S)^2F,
\qquad\|\beta\|_2^2\ll X^{o(1)}(H_*S)^2F.
\tag{3.81}
\]

The second estimate is the same proof with $\vartheta_{\rm sp}=0$.
Here $S\ge1$, as (3.72) requires; the case $S=1/2$ was treated above.
If $\gamma_1$ is supported at $1$, take $q=1,Q=1$.
If $\gamma_3$ is supported at $1$, it is a fixed coordinate.
If the full modulus is 1, the discrepancy is identically zero.

\subsection{Application of Theorem 2.8 and the spectral budget}\label{sec:pi-spectral}

For nonzero (3.69), set

\[
\begin{aligned}
&r_T=z\asymp Z,\ s_T=p_0,\ R_T=Z,\ S_T=p_0/2,\quad\\
&d_T=b,\ c_T=c,\ n_T=p_0m,\ N_T=8p_0Y,
\quad a_T=\varepsilon_y a.
\end{aligned}
\tag{3.82}
\]

Let the Theorem 2.8 coefficient equal $\beta_{z,m}$ at these
indices and zero elsewhere. The support conditions give exactly

\[
(zb,p_0c)=1\iff
(b,c)=(z,c)=(z,p_0)=(b,p_0)=1,
\qquad(a_T,zp_0)=1.
\tag{3.83}
\]

Here $(z,c)=1$ comes from $(s,ac)=(s',ac)=1$.
The $p_0$ conditions on $(b,z)$ come from their roughness.
The smoothed variable $c$ may be divisible by $p_0$.
Reducing the inverse modulo $c$ gives the exact phase identity

\[
e\left(\frac{a_T(p_0m)\overline{zb}_{p_0c}}{p_0c}\right)
 =e\left(\frac{\varepsilon_y a m\overline{zb}_c}{c}\right).
\tag{3.84}
\]

We make the dyadic support of the smooth weight explicit.
Choose a fixed smooth step $\chi$. It is zero on $(-\infty,0]$ and one
on $[1,\infty)$. Split

\[
\begin{aligned}
\psi_{0,1}(v)&=\psi_0(v)[1-\chi((v-0.8)/0.1)],\\
\psi_{0,2}(v)&=\psi_0(v)\chi((v-0.8)/0.1)[1-\chi((v-1.4)/0.2)],\\
\psi_{0,3}(v)&=\psi_0(v)\chi((v-1.4)/0.2).
\end{aligned}
\tag{3.85}
\]

They sum to $\psi_0$.
They lie in the strict dyadic windows $(0.45,0.90]$, $(0.80,1.60]$, $(1.40,2.80]$.
Expanding the $(b,c)$ bumps gives nine weights.
Multiply by fixed smooth cutoffs equal to 1 on the supported $n_T,z,p_0$ indices.
At the lower $n_T$ endpoint use a cutoff zero below $1/2$ and one above 1.
Its derivatives are $O_\nu(n_T^{-\nu})$.
Each weight satisfies the Theorem 2.8 derivative conditions with
$\epsilon_0=\sigma/2000>0$. Take $\epsilon=\sigma/100$,
$A_T=A_0$. Thus (3.14), exactly as stated, applies to (3.69).

We bound the four terms of its bracket. Put

\[
\begin{gathered}
V=T^4R^4S^6,\quad D_*=TR,\quad C_*=TRS,\\
Z\asymp S^2,\quad Y\ll X^{o(1)+\eta}T^2R^2S^3/X.
\end{gathered}
\tag{3.86}
\]

Since $\vartheta_{\rm sp}<1/2$, all powers of $(D,C)$ in (3.15) are nonnegative.
First replace them monotonically by $D_*,C_*$.
Then $C_*\le D_*Z$, and $Y\ll ZX^{-8\sigma+\eta+o(1)}$.
Apply (3.81). Suppress only uniform $X^{o(1)}$ and $p_0$-powers. This gives

\[
\begin{array}{c|c}
\text{bracket term}/V&\text{upper bound}\\ \hline
1&K_A=A_0^{\vartheta_{\rm sp}} TQR^2S^5/X^2\\
2&(A_0/X)K_A\\
3&K_R=(TR/S)^{2\vartheta_{\rm sp}}TQR^2S^5/X^2\\
4&(TQR^2S^5/X^2)(T^2R^2/X).
\end{array}
\tag{3.87}
\]

We check the regular third term directly.
Up to $p_0$-powers, its bracket before the coefficient norm is

\[
\begin{aligned}
C(CD\sqrt Z)^{2\vartheta_{\rm sp}}Z^{1-\vartheta_{\rm sp}}
  (DZ)^{1-2\vartheta_{\rm sp}}
 &=CDZ^2(C/Z)^{2\vartheta_{\rm sp}}\\
 &\le T^2R^2S^5(TR/S)^{2\vartheta_{\rm sp}}.
\end{aligned}
\tag{3.88}
\]

The factor here is $C/Z$; the factor $C/(SZ)$ would lose $S^{2\vartheta_{\rm sp}}$.
The bound $C\le DZ$ is applied after passing to the envelopes.
The first term is at most $A_0^{\vartheta_{\rm sp}} CDZ^2$ times the weighted norm.
The second uses $A_0YZ$. The fourth uses $D^2YZ$.
Substitute $(H_*S)^2F\ll X^{o(1)+2\eta}QT^3R^4S^6/X^2$.
This proves the other lines of (3.87).

The four packing inequalities (3.7) bound the exponents in (3.87) by
$-15\sigma,-(15-1/20)\sigma,-15\sigma,-23\sigma$.
The cutoff contributes at most $3\eta$ to $J_*^2/V$.
The product in (3.14) is eventually at most $X^{10.01}$.
This follows from $A_0<X^{1.01},T<X^{0.34},R,S<X^{0.6}$ and
the sharper $Y$ bound. The bounds on $R$ and $S$ use only
$RS\le X^{0.6}$. This comes from the full-modulus support $q\le X^{0.6}$
in (H3), since the full modulus has scale $QRS$; dyadic constants are
absorbed after enlarging $X_0$.
So the smoothness prefactor costs at most $0.19019\sigma$ in the exponent.

Before the cutoff loss, the exponents relative to
$V^{1/2}$ (from the square-root row on, relative to $V^{1/4}$), in
units of $\sigma$, are

\[
\begin{array}{c|r}
J_*&-299/40=-7.475\\
\text{Theorem 2.8 prefactor}&+0.19019\\
W_{k\ne0}&-7.28481\\
W_{k=0}&-8\\
\text{square root in (3.64)}&-3.642405\\
\text{Mellin integral}&+0.075\\
\text{combined fixed subpolynomial factors}&+0.11\\
\text{cutoff }X^\eta\text{ after both roots}&+0.00075.
\end{array}
\tag{3.89}
\]

The Mellin allowance follows from (3.57).
Any fixed logarithmic power is eventually at most $X^{0.074\sigma}$.
The remaining subpolynomial factors are the finitely many $(h,f,D,C,Z)$
localizations, the divisor representation counts in (3.65), (3.71), the
prime products (3.79), and powers of $p_0\ll L$.
Their product is uniformly $X^{o(1)}$.
So it is at most $X^{0.11\sigma}$ after choosing $X_0$.
The multiplier from the random partition is a
fixed $\sigma$-dependent constant. By
(3.58), (3.64), (3.65), (3.89),

\[
|E_2^*|\ll L^C\frac{T^2}{Q}X^{-3.45\sigma},
\tag{3.90}
\]

since $-3.642405+0.075+0.11+0.00075=-3.456655<-3.45$.
For $S=1/2$, $E_2^*=0$ (see above).
Inserting (3.90) in (3.51) gives a spectral contribution
$M_ITX^{-1.725\sigma}L^C\ll X^{1-1.725\sigma}L^C$, since $M_IT\ll X$.
As $c_\sigma<\sigma/2<1.725\sigma$, it is absorbed in $X^{1-c_\sigma}L^C$.

For a target $A'$, use $A_{\rm work}=A'+4$ on every shell
triple, uniformly before summing.
The triangle inequality over $O(L^3)$ triples costs
the logarithm of the outer modulus and the
two logarithms of the other factors.
Finally restore the factor $C_{\rm slot}$ of (3.5). It is $O_\sigma(1)$.
Equations (3.51), (3.90) prove Theorem~\ref{thm:main}.

\subsection{Global summation, prime powers, and normalization}\label{sec:pi-global}

At fixed $\sigma$, assume (H7): the box family has the
source and Heath--Brown coefficient bound

\[
\sum_j|c_j|\le C_{\sigma,C_{\rm box},\mathcal C_{\rm src}}(\log N)^{K_0},\qquad
K_0\le24C_{\rm box}+26.
\tag{3.91}
\]

\textbf{Source-wide error bounds.}
Assume the source's aggregate logarithm-freezing bound and the
boundary bound (3.24).
With the exponents $c_{\log}$ and $c_R$ of (H6), a sufficient pair of
raw bounds is
\[
 E_{\rm log}\ll_\sigma N(\log N)^{c_{\log}-C_{\rm box}},\qquad
 E_{\rm bdry}\ll_\sigma N(\log N)^{c_R-C_{\rm box}/2}.
\]
They include all source labels and representation multiplicities.
These bounds are separate from the interior-box hypotheses (H1)--(H5).
The source-wide estimates enter through (B.29) for the prime source and
through the auxiliary bound (3.95). They are not hypotheses of
Theorem~\ref{thm:main}.

The source labels and prefix cells are $O_\sigma(1)$. There are
at most 24 HB variables. Hence there are $O(L^{24})$ dyadic tuples, and
each dyadic tuple contains $O(L^{24C_{\rm box}})$ short boxes.
The product shells number $O(\log\log N)\le O(L)$.
The total count is $O_\sigma(L^{24C_{\rm box}+25})$. Since the
coefficients are $O_\sigma(1)$, (3.91) holds with
$K_0=24C_{\rm box}+26$; the count exponent is smaller.
The three factor localizations are inside (3.13).
The count includes every partition.

Choose parameters in the following order.
Fix $\sigma$ and the source combinatorics.
For Theorem~\ref{thm:aa37-interior}, through (3.92) and (B.29), use
the requested exponent $A_{\rm agg}=A$ and choose
$C_{\rm box}\ge\lceil2A+32\rceil$; then set
$A'=A+24C_{\rm box}+31$ and $A_{\rm work}=A'+4$.
This is the parameter convention used in Theorem~\ref{thm:aa37-interior}
and Theorem~\ref{thm:aa37-fixed-u}.

For the stronger source-wide bookkeeping in (3.95), which is outside
Theorem~\ref{thm:main}, fix a separate final exponent $A>0$ and put
$A_{\rm agg}=A+3$.
Choose the box width $C_{\rm box}=C_{\rm box}(A_{\rm agg})$ with
\[
 C_{\rm box}(A_{\rm agg})\ge
 \max\{2(A_{\rm agg}+c_R),\ A_{\rm agg}+c_{\log}\},
\]
so that the boundary and logarithm-freezing errors are
$O(N(\log N)^{-A_{\rm agg}})$. For the Iwaniec source,
$c_R=16$ by (B.19) and the logarithm-freezing exponent is
$c_{\log}=18$ by (B.18); since
$A_{\rm agg}+18<2A_{\rm agg}+32$, the condition is the integer choice
$C_{\rm box}(A_{\rm agg})\ge\lceil2A_{\rm agg}+32\rceil
=\lceil2A+38\rceil$. Set
$A'=A_{\rm agg}+24C_{\rm box}+31$ and $A_{\rm work}=A'+4$.
The value of $A'$ is $K_0+A_{\rm agg}+5$ at the maximal count
$K_0=24C_{\rm box}+26$ in (3.91).
Let $c_*$ dominate $c_1$ and every preliminary logarithmic loss denoted by $c$.
This includes the fixed loss in (3.41) and the endpoint losses before BFI.
It is chosen before $A_E,C_\Sigma,C_M$.
Use the same $c_*$ for every source and every application below.
Choose
\[
\begin{aligned}
 A_E&\ge2A_{\rm work}+2C_{\rm box}+c_*,\\
 C_\Sigma&=A_E+2c_*+3,\\
 C_M&=5C_\Sigma+A_E+C_{\rm box}+c_*+3.
\end{aligned}
\]
Then choose $A_3$ and the required BFI/SW exponents.
Later power-saving logarithmic losses may depend on $C_\Sigma,C_M$.
The constant $c_*$ is defined before them.
Choose $N_0$ last. It is ineffective, as are the implied constants,
because the Siegel--Walfisz theorem is ineffective.

The source may choose a larger width. Its fixed width is selected
before the analytic precisions, with the same displayed formulas.

By (3.13), (H7) in the form (3.91), and $N/(\log N)^{K_{\rm scale}}\le X\le N$,

\[
\begin{aligned}
E_{\rm II}:=\sum_j|c_j||S_j^{\rm PSDT}|
&\ll N(\log N)^{K_0-A'}
 +N^{1-c_\sigma}(\log N)^{K_0+C}\\
&\ll_{\sigma,A,C_{\rm box},K_{\rm scale},\mathcal C_{\rm src}}
N(\log N)^{-A}.
\end{aligned}
\tag{3.92}
\]

Every positive fixed power saving beats every fixed logarithm once $N_0$ is chosen.
The spectral term $X^{1-1.725\sigma}L^C$ is already absorbed in (3.13), since $c_\sigma<\sigma/2$.

\textbf{Source prime powers.} For the direct source the correction in (3.22) has
$O(N^{1/2}(\log N)^C)$ integers before discrepancy weights.
For the product source assume at most four tested coordinates,
each at least $N^\alpha$, $\alpha=100/1149$, and product $e<N$.
A proper power coordinate has reciprocal mass

\[
\sum_{N^\alpha\le p^k\le N,\ k\ge2}p^{-k}
 \ll N^{-\alpha/2}+(\log N)N^{-2\alpha/3}
 \ll N^{-\alpha/2}.
\tag{3.93}
\]

For $k=2$ compare with the tail of $\sum n^{-2}$.
For $k\ge3$ sum the integer $k$-power tail.
There are $O(\log N)$ possible exponents.
Fix one coordinate at $t$.
For $1\le r\le4$, the remaining positive coordinates
have at most $(N/t)(1+\log N)^{\max(r-2,0)}$ possibilities.
For $r=1$ no coordinate remains, and the
single admissible completion is bounded by $N/t$.
For $r=2$ the remaining coordinate is counted directly.
For $r\ge3$ the reciprocal-sum argument gives the displayed logarithmic power.
Take a union bound over the at most four possible power coordinates.
For the product discrepancy write $e$ for the tested product.
Theorem~\ref{thm:aa37-interior} supplies $N\le a\le N\log N/3$.
On the product support, $e<N\le a\ll N\log N$, so
$1\le |a-e|\ll N\log N$.
Use (3.23) and $\tau(|a-e|)\ll_\epsilon N^\epsilon$ with
$\epsilon=\alpha/4$. Then the whole source-level correction obeys

\[
E_{R0}\ll N^{1-\alpha/4}(\log N)^C
       =N^{1-25/1149}(\log N)^C
       \ll_A N(\log N)^{-A}.
\tag{3.94}
\]

This bound precedes the Heath--Brown subdivision, so its box multiplicity is 1.
For a product, the bound uses the coordinate floor $N^\alpha$ assumed
immediately before (3.93).
Thus (3.94) applies to every such product source.

The bounds (H6), the logarithm-freezing estimate, (3.92) (which uses
(H7)), and (3.94), with the fixed choices above, give

\[
\begin{aligned}
&E_{\rm II}+E_{\rm bdry}+E_{\rm log}+E_{R0}
       \ll_{\sigma,A,C_{\rm box},K_{\rm scale},\mathcal C_{\rm src}}
       N(\log N)^{-A-3}=o(U_C),\\
&U_C=C(N)N/(\log N)^2,\quad C(N)\ge c>0.
\end{aligned}
\tag{3.95}
\]

The ratio is
$O_{\sigma,A,C_{\rm box},K_{\rm scale},\mathcal C_{\rm src}}
((\log N)^{-A-1})$.

For each source family, its exact source identities give the
finite sum to which these bounds apply.
The restoration, BFI, partition, Poisson and Theorem 2.8 errors
enter once, inside $E_{\rm II}$.

For any normalization $W_{\rm norm}\gg N/(\log N)^2$, (3.92)
reads explicitly

\[
\begin{aligned}
\frac{\sum_j|c_j||S_j^{\rm PSDT}|}{W_{\rm norm}}
 &\ll(\log N)^{K_0+2-A'}\\
 &\quad+N^{-c_\sigma}(\log N)^{K_0+C+2}=o(1).
\end{aligned}
\tag{3.96}
\]

Consider the original direct prime source, before HB. Its mass is not yet replaced by a logarithmic-integral (Li) mass.
Its exact prime mass $X_A=\sum_pA_D(p)$ satisfies

\[
\sum_{(p,q)=1}A_D(p)=X_A
\tag{3.97}
\]

because $p>N/(\log N)^B>q$ for $q<N^{e_Q}$ and sufficiently large $N$.
Thus the formal and natural comparators coincide on that original source.
This holds for the source itself, before the split into signed HB terms.
If a Li mass $X_{\rm Li}$ is substituted, the exact difference is
$(X_A-X_{\rm Li})/\varphi(q)$.
Its weighted sum enters once, in the main-term mass contribution,
outside the natural-comparator correction.
For the composite product source, (3.97) may fail.

\subsection{The complete non-II cover}\label{sec:pi-nonII}

These results are separate from Theorem~\ref{thm:main}.
The source-to-box identity comes from the source itself.
For a direct shell $[X,2X)$, assume the following source conditions (SH1)--(SH5).
\begin{itemize}[leftmargin=3.2em]
\item[(SH1)] The coefficient $\gamma(q)$ is independent of tested variables, $(q,2N)=1$,
$|\gamma(q)|\le C_\sigma$, $q\le N^{e_Q}$, $e_Q<153/256$,
and $(a,q)=1$ on its support.
\item[(SH2)] All tested restrictions are constant on each nonboundary box.
Smooth factor weights have $\sup|w|+V(w)\ll1$.
\item[(SH3)] All HB factors are rough at the threshold $z^*=z_0(X)$ or $z_0(N)$.
\item[(SH4)] The discrepancy is exactly (3.3).
\item[(SH5)] The assignment priority is proper prime powers,
boundary boxes, Type II, Type I-1, then Type I-2.
\end{itemize}

Put $\omega=\log(N/X)/\log N$ and
$\kappa=30\log2/\log N$. Increase $N_0$ until
$\omega,\kappa\le\sigma/202$. Let $V_i$ be the lower HB scales, and
$e_i=\log V_i/\log N$. Then (3.21) implies

\[
1-\omega-\kappa<\sum_i e_i<1+\kappa,\qquad
e_i\le1/3+\kappa\quad\text{on Möbius variables}.
\tag{3.98}
\]

There are at most six variables.
Call a variable small below $60\sigma$, and large above $1/3+14\sigma$.
Suppose no variable is in the II window.
Since $360\sigma<1-\omega-\kappa$, some variable is large.
Since $1+42\sigma>1+\kappa$, at most two are large.
Möbius variables are never large.
Thus there are exactly one or two large smooth variables.
For a window variable,

\[
\log_X(2V)\le(1/3+14\sigma+\kappa)/(1-\omega)
       \le1/3+15\sigma,
\tag{3.99}
\]

which proves the shell-local window transfer. For one large variable,

\[
e_*-e_Q-(1-\omega)/10
 \ge0.9-153/256-0.9\omega-\kappa-300\sigma>0.3023;
\tag{3.100}
\]

hence its scale $V\ge qX^{1/10}$. With two large variables,
$N_1\ge N_2\ge X^{1/3}$. The at most four other factors are small.
Their product $m$ is split dyadically, $m\asymp M$, with

\[
M\le X^{2.401\cdot10^{-8}},\qquad N_1N_2M\le2^{30}X,\qquad q<X^{0.5977}.
\tag{3.101}
\]

These follow from
$(240\sigma+\kappa)/(1-\omega)<2.401\cdot10^{-8}$,
which uses $\sigma\le10^{-10}$,
and $e_Q/(1-\omega)<0.5977$.
Consider a product source with three or four coordinates at least $N^{100/1149}$.
Its smallest coordinate has exponent between $100/1149$ and $(1+\kappa)/3$.
It has at most six HB factors.
One of them has exponent exceeding $(100/1149-\kappa)/6>60\sigma$.
Each is below $1/3+14\sigma$.
Thus, after proper prime powers and boundary boxes are removed,
all its boxes are of type II.

\textbf{Elementary sieve sandwich.} For a finite prime set $\mathcal P$
below $z$, let $P=\prod_{p\in\mathcal P}p$,
$E=\sum_{p\in\mathcal P}1/p$, and

\[
\lambda_d^{(J)}=\mu(d)1_{d\mid P,\ \omega(d)\le J}.
\tag{3.102}
\]

Let $r=\omega((n,P))$. For $r\ge1$, Pascal's identity gives
$\sum_{j\le J}(-1)^j\binom rj=(-1)^J\binom{r-1}J$. The sum is 1 for $r=0$. Hence

\[
\begin{gathered}
\Psi^-(n):=\sum_{d\mid n}\lambda_d^{(2k-1)}
 \le1_{(n,P)=1}\le \Psi^+(n):=\sum_{d\mid n}\lambda_d^{(2k)},\\
\Psi^+\ge0,\qquad \Psi^+-\Psi^-=\binom r{2k}\ge0.
\end{gathered}
\tag{3.103}
\]

The weights are 1-bounded and supported on $d\le z^{2k}$.
Also
$\sum_{d\mid P,\omega(d)=j}1/d\le E^j/j!$.
Let $s=\log D/\log z\gg(\log\log X)^3$, and choose
$2k=2\lfloor s/2\rfloor$. Mertens' bound
$E\le\log\log z+O(1)$, with $s\gg(\log\log X)^3$, gives
$2k\ge s-2\ge e^2E$ for large $X$. With Stirling,

\[
E^{2k}/(2k)!\le(eE/(2k))^{2k}\le e^{2-s}
                   \ll_A L^{-A}.
\tag{3.104}
\]

It is the sieve approximation used in I-1 and I-2.

\setcounter{lemma}{2}
\begin{lemma}[Type I, one large variable (I-1)]\label{lem:typeI1}
Assume (SH1)--(SH5).
Consider a box with one large variable.
Let the complementary coefficient satisfy $|\alpha_m|\ll\tau(m)^{B_0}$.
Uniformly for every integer residue $a$ coprime to each supported $q$,

\[
\sum_q|\gamma(q)|\left|\sum_m\alpha_m
 \sum_{v\in I,\ P^-(v)\ge z^*}w(v)\mathcal S_{mv}(q;a)\right|
 \ll_{\sigma,A,B_0}XL^{-A}.
\tag{3.105}
\]
\end{lemma}

\begin{proof}
If $(m,q)>1$, the summand vanishes. Otherwise put
$b=a\bar m(q)$.
The progression and the reduced comparator contain only integers prime to $q$.
Apply (3.102)--(3.104) to the $q$-dependent set $\mathcal P_q=\{p<z^*:p\nmid q\}$, with
$D=X^{1/20}$. For every $d\mid P_q$, $(d,q)=1$, and

\[
\#\{v\le t:v\equiv b(q),d\mid v\}=t/(dq)+O(1),
\]
\[
\#\{v\le t:(v,q)=1,d\mid v\}
 =(\varphi(q)/q)(t/d)+O(\tau(q)).
\tag{3.106}
\]

After comparator normalization, the upper and lower sandwiches have the same main term.
Their gap gives discrepancy
$\ll(t/q)e^{2-s}+2D\tau(q)$, where
$s\ge(\log\log X)^3/40$.
Endpoint subtraction and partial summation cost $O(\sup|w|+V(w))$.
By (3.100), $D\tau(q)\ll(V/q)X^{-1/20+o(1)}$.
Thus the inner discrepancy is $O_A(V/(qL^A))$.
Sum with $\sum_m|\alpha_m|\ll ML^C$, $MV\asymp X$, and
$\sum_q1/q\ll L$. Then request a larger $A$. This proves
(3.105). The sieve weights use the prime set $\mathcal P_q$.
\end{proof}

\begin{lemma}[Completed hyperbola]\label{lem:hyperbola}
Let $c\bmod q$ be a unit. For integer
intervals $J_1,J_2$, put

\[
\mathcal D=\#\{(u,v)\in J_1\times J_2:uv\equiv c(q)\}
       -\frac{\#J_1^*\#J_2^*}{\varphi(q)}.
\tag{3.107}
\]

The full-modulus Kloosterman-sum bound
\cite[Corollary~11.12]{IK04},

\[
|S(h,k;q)|\le\tau(q)q^{1/2}(h,k,q)^{1/2},
\tag{3.108}
\]
implies, uniformly in interval lengths and positions,

\[
|\mathcal D|\ll\tau(q)q^{1/2}\log^2(2q)
        +\tau(q)^2\log^2(2q)/\varphi(q).
\tag{3.109}
\]
\end{lemma}

\begin{proof}
Put
\[
F_i(h)=\sum_{u\in J_i}e(hu/q).
\]
Additive completion then gives
\[
\mathcal D=q^{-2}\sum_{h,k}F_1(-h)F_2(-k)
\left(S(h,kc;q)-\frac{c_q(h)c_q(k)}{\varphi(q)}\right).
\]

The bracket vanishes if
$h=0$ or $k=0$, so complete periods cancel. For $0<h<q$,
$|F_i(h)|\le q/(2\min(h,q-h))$, and

\[
\sum_{0<h<q,d\mid h}|F_i(h)|\ll(q/d)\log(2q).
\tag{3.110}
\]

Use $(h,k,q)^{1/2}\le\sum_{d\mid(h,k,q)}d^{1/2}$ in (3.108)
and sum $\sum d^{-3/2}$. For the principal part use
$|c_q(h)|\le(h,q)=\sum_{d\mid(h,q)}\varphi(d)$, so its
one-variable weighted sum is $O(q\tau(q)\log(2q))$. This proves
(3.109). Modulo 1 the discrepancy is zero.
\end{proof}

\begin{lemma}[Type I, two large variables (I-2)]\label{lem:typeI2}
Assume (SH1)--(SH5) and (3.101).
Consider two large smooth variables with nonnegative weights $w_1,w_2$.
Let $|\alpha_m|\ll\tau(m)^{B_0}$.
Uniformly for every integer residue $a$ coprime to each supported $q$,

\[
\begin{aligned}
&\sum_{q\asymp Q}|\gamma(q)|\Bigl|\sum_m\alpha_m
 \sum_{n_i\in I_i,\ P^-(n_i)\ge z^*}
        w_1(n_1)w_2(n_2)\mathcal S_{mn_1n_2}(q;a)\Bigr|\\
&\qquad\ll_A XL^{-A}.
\end{aligned}
\tag{3.111}
\]
\end{lemma}

The smooth factor weights in the prime-source application of
Appendix~\ref{app:proof-aa37} are nonnegative.
For a general real weight, write $w=(w+c)-c$,
with $c=\sup|w|$, and apply the lemma to the nonnegative pieces.

\begin{proof}[Proof of Lemma~\ref{lem:typeI2}]
Again restrict to $(m,q)=1$. Let $C(FG)$ and $M_q(FG)$
be the positive congruence and comparator functionals with the
positive weights $w_1,w_2$. Put $\Delta=C-M_q$.
Apply (3.103) to all primes below $z^*$, with
$D_f=X^{\epsilon_I/2}$, $\epsilon_I=\delta/4=7.5\sigma$.
Include the Brun--Bonferroni threshold for this smaller level in $N_0$.
For every fixed logarithmic exponent $A_{\rm sv}$, put
\[
s_f=\frac{\log D_f}{\log z^*},\qquad
E=\sum_{p<z^*}\frac1p,\qquad J=2\lfloor s_f/2\rfloor,
\]
and require $J\ge\max\{2,e^2E,A_{\rm sv}\log L\}$.
At fixed $\sigma>0$ this holds eventually,
since $s_f\asymp\sigma(\log\log N)^3$ and $E=O(\log\log N)$.
Hence $(z^*)^J\le D_f$, $J-1\ge1$, and
\[
\frac{E^J}{J!}\le\left(\frac{eE}{J}\right)^J\le L^{-A_{\rm sv}}.
\]
Choose $A_{\rm sv}$ after the fixed logarithmic losses in (3.115).
So the same final choice of $N_0$ covers the use of (3.103)--(3.104) in I-2.
Write $R_i=1_{P^-(n_i)\ge z^*}$, $G_i=\Psi_i^+-\Psi_i^-\ge0$.
Since $0\le R_i\le \Psi_i^+$,

\[
\Psi_1^+\Psi_2^+-G_1\Psi_2^+-\Psi_1^+G_2\le R_1R_2\le \Psi_1^+\Psi_2^+.
\tag{3.112}
\]

Since $C,M_q$ are positive,

\[
\begin{aligned}
|\Delta(R_1R_2)|\le{}&|\Delta(\Psi_1^+\Psi_2^+)|+|\Delta(G_1\Psi_2^+)|
 +|\Delta(\Psi_1^+G_2)|\\
&+M_q(G_1\Psi_2^+)+M_q(\Psi_1^+G_2).
\end{aligned}
\tag{3.113}
\]

The lower weights $\Psi_i^-$ may be negative.
The bound (3.113) contains only the nonnegative $\Psi_i^+,G_i$.
Expand the first three terms over $d_i\le D_f$.
Their coefficients are bounded by 1. Terms with $(md_1d_2,q)>1$ vanish.
Otherwise put $c\equiv a\overline{md_1d_2}\pmod q$.
This is a unit, and $u_i=n_i/d_i$ lies on the modular hyperbola
$u_1u_2\equiv c\pmod q$.
Two-variable partial summation and (3.109) give a total
$\ll_\epsilon D_f^2q^{1/2}X^\epsilon$.

For the two main gaps, positivity lets us drop the restriction
$(n_1n_2,q)=1$. Full-integer progression counts and (3.104) give

\[
\sum_{n\in I_i}G_i(n)w_i(n)\ll_{A_{\rm sv}} N_iL^{-A_{\rm sv}}+D_f,
\quad\sum_{n\in I_i}\Psi_i^+(n)w_i(n)\ll N_i+D_f.
\tag{3.114}
\]

The progression counts range over all integers.
Summing over $(m,q)$ now gives

\[
\ll M D_f^2Q^{3/2}X^\epsilon L^C+XL^{C-A_{\rm sv}}
       +M D_f(N_1+N_2)L^C+M D_f^2L^C.
\tag{3.115}
\]

By (3.101), the first exponent is less than
$2.401\cdot10^{-8}+7.5\cdot10^{-10}+1.5(0.5977)+\epsilon<0.897$
for a sufficiently small fixed $\epsilon>0$. Also
$M(N_1+N_2)\le2^{31}X/N_2\ll X^{2/3}$.
So the exponents of the last gap terms are less than 0.668.
Choose $A_{\rm sv}$ after the fixed log losses.
This proves (3.111). Squareful moduli are covered by (3.108).
\end{proof}

Finally there are $O_\sigma(L^K)$ non-II direct boxes and shell
pieces for a fixed $K$. Request (3.105), (3.111) with exponents
larger than $K+A+5$. Their sum is then $O(NL^{-A})$; in particular,
divided by $N/(\log N)^2$ it tends to zero. For the prime source these
are the exponents $A_i=K_I+A+5$ of Appendix~\ref{app:proof-aa37}.
The I-1 and I-2 boxes are disjoint parts of the non-II contribution.
Under its coordinate hypotheses, the product source has none.
Their contributions lie outside $E_{\rm II}$.

For the prime source, Section~4 supplies the packing construction and
Proposition~\ref{prop:aa37-prime-source} in Appendix~\ref{app:proof-aa37}
gives the prime-source estimate.


\section{Three-slot packing for Iwaniec's remainder coefficients}\label{sec:pack}

This section proves a general packing theorem. It then verifies the
hypotheses for the remainder coefficients of Iwaniec's linear sieve
\cite[Theorem~4]{AA37}. The packing and coefficient identities provide
the modulus factorization used in the prime-discrepancy estimates below.
The functions below specialize Lichtman's optimization
\cite[Proposition~6.1, (6.2), (6.4)]{Lichtman23} to spectral exponent
$7/32$ and moving-residue exponent $1$. We prove the labelled-atom
statement here; its two-bin step is related to the corrected argument
in Pascadi's Lemma~5.6 \cite{Pascadi25}.

\subsection{A combinatorial packing theorem}

Put
\[
 \overline H=\frac{153}{256},\qquad \mu=\frac{25}{128},
 \qquad \vartheta=\frac7{32}.
\]
Here $\vartheta$ is the fixed spectral exponent $2\theta_{\rm RP}$.
Later in this section $\theta$ denotes the sieve level.
For $x\ge0$, define
\[
 v(x)=\min\left\{\frac{1+x}{2},\frac{32-34x}{39},
                         \frac{57}{64}-\frac32x\right\}.
\]
For real $x$ and $y\ge0$, put $\psi(x,y)=x\mathbf 1_{x\ge y}$.
Also put
\[
 w(a,b,c)=\psi\left(\min\left\{
       \frac{132-68c}{214},\frac{153-96c}{256},1-2b\right\},
       \frac{1+a}{2}\right).
\]
For an ordered triple $a\ge b\ge c\ge0$, set
\[
\begin{split}
 H(a,b,c)=\max\{&v(a),v(b),v(a+b),v(a+b+c),w(a,b,c),w(a,c,b),\\
                 &\psi(v(a+c),a+2b+c),
                   \psi(v(b+c),2a+b+c)\},
\end{split}
\]
Define $\widetilde H(a,b,c)=H(a,b,c)$ if $a\le\mu$.
If $a>\mu$, put $\widetilde H(a,b,c)=v(a)$. For a finite non-increasing
list $e_1\ge\cdots\ge e_m\ge0$, let $\bar e_i=e_i$ for $i\le m$
and $\bar e_i=0$ for $i>m$. Repeated values are distinct labelled
atoms. The empty list is allowed.

A list $f_1\ge\cdots\ge f_m\ge0$ satisfies $\mathrm{WF}(y)$ if
\[
 f_1+\cdots+f_{j-1}+2f_j\le y\qquad(1\le j\le m).
 \tag{4.1}\label{eq:pack-wf}
\]
Write
\[
 b_-=2h-1,\qquad b_A=\frac{32-39h}{34},\qquad
 b_B=\frac{2-\vartheta-2h}{3},\qquad b_+=\min(b_A,b_B).
\]

\refstepcounter{theorem}\textbf{Theorem~\thetheorem\ (three-slot packing).}\label{thm:pack-a}
Let $0<\delta\le3\cdot10^{-9}$,
$1/2\le h\le\overline H$, $d_*=h-2\delta$, and
$\delta\le g\le1/3-\delta/2$. Suppose the labelled list
$e_1\ge\cdots\ge e_m\ge0$ satisfies
\[
 \mathrm{WF}(d_*)\quad\hbox{and}\quad
 h\le\widetilde H(\bar e_1,\bar e_2,\bar e_3).
 \tag{4.2}\label{eq:pack-hypotheses}
\]
Then the deterministic rule $\mathcal R$ below partitions the labels
into $J_1\sqcup J_2\sqcup J_3$. If
$\Sigma_s=\sum_{i\in J_s}e_i$, then
\[
\begin{aligned}
 \Sigma_1&\le g,\\
 2g+\Sigma_2+2\Sigma_3&\le1-4\delta,\\
 \vartheta+g+\Sigma_1+5\Sigma_2+2\Sigma_3&\le2-4\delta,\\
 (1+2\vartheta)g+\Sigma_1+(5-2\vartheta)\Sigma_2
       +(2+2\vartheta)\Sigma_3&\le2-4(1+\vartheta)\delta.
\end{aligned}
 \tag{4.3}\label{eq:pack-loads}
\]
The output depends only on the labelled list and on $h,\delta,g$.

\refstepcounter{theorem}\label{lem:pack-1}\textbf{Lemma~\thetheorem.} For $x\ge0$,
$v(x)=\min\{(1+x)/2,57/64-3x/2\}$. Thus $v$ is concave,
$v\le\overline H$, and $\widetilde H\le\overline H$.

\emph{Proof.} Consider the middle affine function in the definition
of $v$. Its difference with $(1+x)/2$ is $(25-107x)/78$, which is
nonnegative for $x\le25/107$. Its difference with $57/64-3x/2$ is
$(1568x-175)/2496$, nonnegative for $x\ge175/1568$.
As $175/1568<25/107$, it is never the strict minimum of all three.
The minimum of the remaining two affine functions is concave and has
maximum $(1+\mu)/2=\overline H$. Each $w$ is zero or at most
$(153-96c)/256\le\overline H$. Every $\psi$ branch is zero or a value
of $v$. This proves both upper bounds. \hfill$\square$

\refstepcounter{theorem}\label{lem:pack-2}\textbf{Lemma~\thetheorem.} For $1/2\le h\le\overline H$ and
$0<\delta\le3\cdot10^{-9}$:
\begin{enumerate}
\item $h\le v(b)$ for $b\ge0$ exactly when $b_-\le b\le b_+$.
\item $b_+=b_B$, and $b_A>b_B$.
\item $b_+-b_-=(153-256h)/96\ge0$ and
      $b_+-\mu=(153/128-2h)/3\ge0$.
\item $d_*-b_+-2b_-\ge3/256-2\delta>0$ and
      $d_*-3b_-\ge3/256-2\delta>0$.
\item $\vartheta+2(d_*+2\delta)+3b_B=2$ and
      $(2+2\vartheta)(d_*+2\delta)+(3-4\vartheta)b_A=2$.
\end{enumerate}

\emph{Proof.} The three affine branches of $v$ give
$b\ge2h-1$, $b\le b_A$, and $b\le b_B$.
These are closed inequalities. The difference
$b_A-b_B=(567/16-49h)/102$ decreases with $h$ and at
$h=\overline H$ is $525/8704>0$, proving (2). The formulas in
(3) and the first bound in (4) follow by expansion. Their minima occur
at $h=\overline H$. For the second bound in (4), use
$d_*-b_-=1-h-2\delta$ and $3-5\overline H=3/256$.
For (5), use $2+2\vartheta=39/16$,
$3-4\vartheta=17/8$, and
$(3-4\vartheta)b_A=2-39h/16$. Finally
$3/256>6\cdot10^{-9}\ge2\delta$. \hfill$\square$

\refstepcounter{theorem}\label{lem:pack-3}\textbf{Lemma~\thetheorem.} If $h>0$ and $w(x,y,z)\ge h$, then
$z\le b_+-b_-$.

\emph{Proof.} Since $h>0$, the activation in $\psi$ is on. Hence
$(132-68z)/214\ge h$ and $(153-96z)/256\ge h$. These give
$z\le b_A-b_-$ and $z\le b_B-b_-$, respectively. \hfill$\square$

\refstepcounter{theorem}\label{lem:pack-4}\textbf{Lemma~\thetheorem\ (two bins).} Let $c_1,c_3\ge0$ and initial loads
$u_0\le c_1$, $w_0\le c_3$. Suppose a finite list
$f_1,\ldots,f_k\ge0$ satisfies
\[
 u_0+w_0+\sum_{i<j}f_i+2f_j\le c_1+c_3\qquad(1\le j\le k).
 \tag{4.4}\label{eq:pack-two-bin}
\]
Process the atoms in order. Place each in bin 1 if it fits, and
otherwise in bin 3 if it fits. Every atom is placed, and the final loads
are at most $c_1,c_3$.

\emph{Proof.} Suppose $f_j$ is the first atom fitting neither bin.
Then the current loads $U,W$ satisfy
$U+W=u_0+w_0+\sum_{i<j}f_i$, $U+f_j>c_1$, and
$W+f_j>c_3$. Adding the last two inequalities contradicts
\eqref{eq:pack-two-bin}. \hfill$\square$

The rule permits zero atoms and repeated values. Its three cases are:
\begin{enumerate}
\item If $g+b_+\ge d_*$, apply Lemma~\ref{lem:pack-4} to the whole list.
      Use bin caps $g,d_*-g$ for slots 1 and 2, and leave slot 3 empty.
\item If $g+b_+<d_*$, test in order
      $K=\{1\},\{1,2\},\{1,2,3\},\{1,3\},\{2,3\}$. Put
      $K'=K\cap\{1,\ldots,m\}$ and $b=\sum_{i\in K'}e_i$.
      Call $K$ legal if $b_-\le b\le b_+$ and the remaining atoms,
      in their original order, satisfy $\mathrm{WF}(d_*-b)$. For
      the first legal $K$, place $K'$ in slot 2. Apply Lemma~\ref{lem:pack-4}
      to the rest in slots 1 and 3, with caps $g,d_*-g-b$.
\item If no $K$ is legal, use caps $g,b_-,d_*-g-b_-$ for slots
      1, 2, 3. Put atom 1 in slot 2 if it exists. Process subsequent
      atoms in order, placing each in the first slot in which it fits.
      Put the first atom fitting no slot in slot 2. Call the
      resulting slot-2 load $b$. Apply Lemma~\ref{lem:pack-4} to the remaining
      atoms in slots 1 and 3, with caps $g,d_*-g-b$.
\end{enumerate}

\textbf{Proof of Theorem~\ref{thm:pack-a}.} Write $\mathcal R$ for the rule just
defined. Since $h\ge1/2$, $b_-\ge0$; also
$d_*=h-2\delta\ge1/2-6\cdot10^{-9}>1/3\ge g$.
It is enough to prove one of the following envelopes:
\[
\begin{split}
 (E_{\rm glob}):&\quad g\ge d_*-b_+,\quad
 \Sigma_1\le g,\quad\Sigma_2\le d_*-g,\quad\Sigma_3=0;\\
 (E_{\rm crit}(b)):&\quad b_-\le b\le b_+,\quad
 \Sigma_1\le g,\quad\Sigma_2\le b,\quad
 \Sigma_3\le d_*-g-b.
\end{split}
\]
The second, third, and fourth left sides in \eqref{eq:pack-loads}
have nonnegative load coefficients. Under $E_{\rm crit}(b)$,
Lemma~\ref{lem:pack-2}(5) gives
\[
\begin{aligned}
2g+\Sigma_2+2\Sigma_3&\le2d_*-b\le2d_*-b_-=1-4\delta,\\
\vartheta+g+\Sigma_1+5\Sigma_2+2\Sigma_3
 &\le\vartheta+2d_*+3b
 \le\vartheta+2d_*+3b_B=2-4\delta.
\end{aligned}
\]
For the fourth load, similarly,
\[
\begin{aligned}
 &(1+2\vartheta)g+\Sigma_1+(5-2\vartheta)\Sigma_2
       +(2+2\vartheta)\Sigma_3\\
 &\qquad\le(2+2\vartheta)d_*+(3-4\vartheta)b\\
 &\qquad\le(2+2\vartheta)d_*+(3-4\vartheta)b_A\\
 &\qquad=2-4(1+\vartheta)\delta.
\end{aligned}
\]
Here $3-4\vartheta=17/8>0$. Under $E_{\rm glob}$, use
$g\ge d_*-b_+$ in the two expressions whose coefficient of $g$
is negative. This gives
\[
\begin{aligned}
2g+\Sigma_2&\le g+d_*
 \le\tfrac13-\tfrac\delta2+\overline H-2\delta
 =\tfrac{715}{768}-\tfrac52\delta\le1-4\delta,\\
\vartheta+2g+5\Sigma_2&\le\vartheta+5d_*-3g
 \le\vartheta+2d_*+3b_+=2-4\delta,\\
(2+2\vartheta)g+(5-2\vartheta)\Sigma_2
 &\le(5-2\vartheta)d_*-(3-4\vartheta)g\\
 &\le(2+2\vartheta)d_*+(3-4\vartheta)b_+
 \le2-4(1+\vartheta)\delta.
\end{aligned}
\]
The first line uses $1-715/768=53/768>3\delta/2$.
In both envelopes $\Sigma_1\le g$, proving the first inequality
of \eqref{eq:pack-loads}. Equality at the case boundary belongs to
case 1; all displayed endpoint inequalities are weak.

Consider case 1. The clause $\mathrm{WF}(d_*)$ of
\eqref{eq:pack-hypotheses} is exactly the hypothesis
\eqref{eq:pack-two-bin} for the whole list. Here $u_0=w_0=0$ and the
total cap is $d_*$.
Lemma~\ref{lem:pack-4} gives $E_{\rm glob}$. Now suppose
$g+b_+<d_*$. If a tested set $K$ is legal, the remaining list
satisfies \eqref{eq:pack-two-bin} with total cap $d_*-b$. Lemma~\ref{lem:pack-4}
then gives $E_{\rm crit}(b)$.

We verify the rule when no tested set is legal. Two deletion facts
will be used. First let $K'$ be a leading set $\{1,\ldots,k\}$. Subtract
$b=\sum_{i\le k}e_i$ from every later inequality in
$\mathrm{WF}(d_*)$. This proves $\mathrm{WF}(d_*-b)$ for the remaining
list. Next let $K'$ be $\{1,3\}$, $\{2,3\}$, or, when $m=2$, the
truncation $\{2\}$ of $K=\{2,3\}$. Let $b_-\le b\le b_+$.
The only unselected atom among the first three has index $i_0=2$ or $1$.
If it is below $b_-$, it is first in the remaining list, and
\[
 2e_{i_0}+b<2b_-+b_+\le d_*
\]
by Lemma~\ref{lem:pack-2}(4). When $m=2$ this is the whole remaining list.
Now take a remaining rank at least four. Its original
prefix sum contains both selected atoms. Subtracting $b$ proves the
reduced mask. Thus the nonconsecutive set is legal whenever these
conditions hold.

Assume now that no tested set is legal. First, $\bar e_1\le\mu$. Otherwise the definition of $\widetilde H$ and
\eqref{eq:pack-hypotheses} give $h\le v(e_1)$. Then
$b_-\le e_1\le b_+$ by Lemma~\ref{lem:pack-2}(1), and $\{1\}$ is legal.
The same lemma and $b_+\ge\mu$ show $\bar e_1<b_-$; otherwise
$\{1\}$ is legal. Hence every padded coordinate is below $b_-$.
Since $\bar e_1\le\mu$, the branch condition gives
$\widetilde H=H$. A branch in its maximum is at least $h$.
Neither $v(\bar e_1)$ nor $v(\bar e_2)$ can be that branch.
Each would put the corresponding coordinate at least $b_-$. A branch
$v(\bar e_1+\bar e_2)$ or
$v(\bar e_1+\bar e_2+\bar e_3)$ would make a leading set legal.
If
\[
 \psi(v(\bar e_1+\bar e_3),\bar e_1+2\bar e_2+\bar e_3)\ge h,
\]
then $v(\bar e_1+\bar e_3)\ge h$, and
$b=\bar e_1+\bar e_3\in[b_-,b_+]$. The only unselected atom among
the first three is $e_2<b_-$. So the deletion fact makes $\{1,3\}$
legal, a contradiction. The other $\psi$ branch similarly makes
$\{2,3\}$ legal. Therefore either
$w(\bar e_1,\bar e_2,\bar e_3)\ge h$ or
$w(\bar e_1,\bar e_3,\bar e_2)\ge h$. Lemma~\ref{lem:pack-3} gives, respectively,
\[
 \bar e_3\le b_+-b_-,\qquad\hbox{or}\qquad
 \bar e_2\le b_+-b_-.
 \tag{4.5}\label{eq:pack-w-tail}
\]

In case 3, the slot-3 cap is
\[
d_*-g-b_-=1-h-g-2\delta\ge
1-\overline H-\frac13-\frac32\delta
=\frac{53}{768}-\frac32\delta>0.
\]
Thus all three caps are nonnegative. Atom 1 fits slot 2 because
$e_1<b_-$. In the first $w$ case, atom 2 cannot
be a stopping atom. Slots 1 and 3 are empty. If $e_2$ fit neither,
then $2e_2>d_*-b_-\ge2b_-$ by Lemma~\ref{lem:pack-2}(4), contrary to
$e_2<b_-$. Thus any stopping atom has index $j\ge3$ and
$e_j\le e_3\le b_+-b_-$. In the second $w$ case, a stopping atom
has $j\ge2$ and $e_j\le e_2\le b_+-b_-$. If there is no stop,
the final loads are bounded by $(g,b_-,d_*-g-b_-)$. This is
$E_{\rm crit}(b_-)$, also when $m=0$.

If the stop is at $j$, let $(u_0,v_0,w_0)$ be the loads before
$e_j$. After placing $e_j$ in slot 2, its load $b=v_0+e_j$ satisfies
\[
 b_-<b\le b_-+(b_+-b_-)=b_+.
\]
The lower bound holds because $e_j$ did not fit slot 2. The upper
bound follows from \eqref{eq:pack-w-tail}. The $j$-th inequality in
$\mathrm{WF}(d_*)$ gives
$u_0+w_0+e_j\le d_*-b$. Since $e_j$ did not fit slot 1,
$u_0+e_j>g$, so $w_0<d_*-g-b$. For every later rank, the original
mask is exactly the two-bin condition with total cap $d_*-b$.
Lemma~\ref{lem:pack-4} completes the assignment and gives $E_{\rm crit}(b)$.
The rule processes every atom once and has fixed tie priorities, so it
is well defined. All cases imply \eqref{eq:pack-loads}. \hfill$\square$

\subsection{A local ceiling and a small-atom criterion}

In this subsection write $\theta_{\rm RP}$ for the exponent in the
unramified Ramanujan bound and $\vartheta_{\rm RP}=2\theta_{\rm RP}$. The
Kim--Sarnak bound \cite{KS03} gives $\theta_{\rm RP}=7/64$ and
$\vartheta_{\rm RP}=7/32=\vartheta$.

\refstepcounter{theorem}\textbf{Proposition~\thetheorem\ (local ceiling and small-atom packing).}\label{prop:pack-ceiling}
Suppose
$0\le\vartheta_{\rm RP}\le7/32$, that is $0\le\theta_{\rm RP}\le7/64$,
and $0\le t\le1/3$. Let
$x,y,z_3\ge0$ satisfy the zero-loss local slot constraints
\[
\begin{aligned}
 x-t&\le0, &2t+y+2z_3&\le1,\\
 \vartheta_{\rm RP}+t+x+5y+2z_3&\le2,
 &(1+2\vartheta_{\rm RP})t+x+(5-2\vartheta_{\rm RP})y+(2+2\vartheta_{\rm RP})z_3&\le2,
\end{aligned}
\]
and $x+y+z_3\le1-t$. Then their maximum total load is
\[
 H_{\vartheta_{\rm RP}}=\frac{5-\vartheta_{\rm RP}}{8}=\frac{5-2\theta_{\rm RP}}8.
\]
For fixed $t$, equality holds precisely at
$x=t$, $y=(1-\vartheta_{\rm RP})/4$, $z_3=(3+\vartheta_{\rm RP})/8-t$.

Suppose instead that $0\le e_j\le u$ with $u\ge0$. Let the
non-increasing labelled atoms obey $\mathrm{WF}(h-2\delta)$, with
\[
 \frac12\le h\le H_{\rm small}(\vartheta_{\rm RP},u):=\frac{5-\vartheta_{\rm RP}-3u}{8},
 \qquad 0<\delta\le\frac1{100},\qquad 0\le g\le\frac13.
\]
Then they admit a three-slot partition with the four load
inequalities \eqref{eq:pack-loads}. Here $\vartheta_{\rm RP}$ replaces $\vartheta$.
In particular, $H_{\rm small}(7/32,1/500)=19101/32000$ and
$H_{7/32}=153/256$.

\emph{Proof.} Seven times $x-t\le0$, three times
$2t+y+2z_3\le1$, and once $\vartheta_{\rm RP}+t+x+5y+2z_3\le2$ give
$8(x+y+z_3)+\vartheta_{\rm RP}\le5$. Equality forces equality in these three
constraints, giving the stated point. Conversely, that point has
$y\ge25/128$, $z_3\ge1/24$, and total load plus $t$ at most
$23/24$. Substitution in the fourth constraint gives
$2-3\vartheta_{\rm RP}(1-\vartheta_{\rm RP})/4\le2$, proving the maximum and its equality cases.

For the small-atom assertion put
\[
\begin{aligned}
 d&=h-2\delta,& b_-&=2h-1,&
 b_B&=\frac{2-\vartheta_{\rm RP}-2h}{3},\\
 b_A&=\frac{2-(2+2\vartheta_{\rm RP})h}{3-4\vartheta_{\rm RP}},&
 b_+&=\min(b_A,b_B).
\end{aligned}
\]
The denominators are positive, and direct subtraction gives
\[
 b_B-b_-=\frac{5-\vartheta_{\rm RP}-8h}{3},\qquad
 b_A-b_-=\frac{5-4\vartheta_{\rm RP}-(8-6\vartheta_{\rm RP})h}{3-4\vartheta_{\rm RP}}.
\]
Hence $u\le b_+-b_-$ is equivalent to
\[
 h\le\frac{5-\vartheta_{\rm RP}-3u}{8},\qquad
 h\le\frac{5-4\vartheta_{\rm RP}-(3-4\vartheta_{\rm RP})u}{8-6\vartheta_{\rm RP}}.
\]
The second upper bound exceeds the first by
$[6\vartheta_{\rm RP}(1-\vartheta_{\rm RP})+14\vartheta_{\rm RP} u]/[8(8-6\vartheta_{\rm RP})]\ge0$.
So the assumed bound on $h$ gives $u\le b_+-b_-$. Start with capacities
$(g,b_-,d-g-b_-)$. They are nonnegative, since the last is at least
$1-5/8-1/3-2\delta\ge1/24-1/50>0$. Place atoms in order in the
first slot in which they fit. If all fit, take $b=b_-$. Otherwise, put the first failed atom $e_j$ in slot 2. Let $b$ be the
new slot-2 load. Its failure there and $e_j\le u$ give
$b_-<b\le b_-+u\le b_+$. Let $(U,V,W)$ be the previous loads.
Then $\mathrm{WF}(d)$ gives $U+W+e_j\le d-b$. Failure in slot 1 gives
$W<d-g-b$. For each later atom, subtract the fixed slot-2 load $b$
from its original $\mathrm{WF}(d)$ inequality. Lemma~\ref{lem:pack-4} places the
remaining atoms in slots 1 and 3 with capacities $g,d-g-b$.
Both cases reach
\[
 \Sigma_1\le g,\qquad\Sigma_2\le b,\qquad
 \Sigma_3\le d-g-b,\qquad b_-\le b\le b_+.
\]
The last three left sides of \eqref{eq:pack-loads}, with $\vartheta_{\rm RP}$ in
place of $\vartheta$, are at most
\[
 \begin{gathered}
 2d-b_-=1-4\delta,\qquad
 \vartheta_{\rm RP}+2d+3b_B=2-4\delta,\\
 (2+2\vartheta_{\rm RP})d+(3-4\vartheta_{\rm RP})b_A=2-4(1+\vartheta_{\rm RP})\delta,
 \end{gathered}
\]
respectively. This proves the partition claim. \hfill$\square$

These are a relaxed local spectral ceiling and a sufficient uniform
small-atom packing criterion. The actual local inequalities have
positive losses, so the relaxed ceiling is a supremum. For example,
take the losses $29\sigma$, $8\sigma$, and $15\sigma$ in the first
three slot constraints. They give
\[
 x+y+z_3\le H_{\vartheta_{\rm RP}}-\frac{242\sigma}{8}
                    -\frac{\vartheta_{\rm RP}\sigma}{160}
\]
when the moving-frequency factor is $X^{\vartheta_{\rm RP}(1+\sigma/20)}$.
The small-atom bound also needs strict room in an actual source.
If $h=\rho \theta+2\delta\le H_{\rm small}(\vartheta_{\rm RP},u)$ and
$\delta>0$, then $\theta<H_{\rm small}(\vartheta_{\rm RP},u)$.
The next subsection applies this criterion to the structured
coefficient family of Iwaniec's construction.

\subsection{Endpoint geometry for Iwaniec's coefficients}

Fix the parameters in Theorem~\ref{thm:aa37-interior} and put
\[
\begin{gathered}
 \tau=\sigma^9,\qquad \rho=1+\tau,\qquad
 \Theta(u)=H_{\rm small}(\vartheta,u)=\frac{5-\vartheta-3u}{8},\\
 \eta_\Theta=\Theta(u)-\theta,\\
 \delta=30\sigma,\qquad d_*=\rho \theta,\qquad h=d_*+2\delta.
\end{gathered}
 \tag{4.6}\label{eq:pack-aa37-parameters}
\]
Proposition~\ref{prop:pack-ceiling} gives the sufficient packing threshold $\Theta(u)$.
The relaxed local ceiling is $H_\vartheta=153/256$. The fixed
support exponent for these coefficients is $\Theta(u)$, and
$\Theta(u)<H_\vartheta$ because $u>0$. The margin
$\eta_\Theta=\Theta(u)-\theta$ leaves room for the smooth divisor and tail-prime
ranges. Since the theorem also imposes $D^{\sigma^2}\le z$, every
nonempty tail has the required roughness floor. Lemma~\ref{lem:aa37-pack} verifies
$\rho \theta+60\sigma<\Theta(u)$ and the hypotheses of
Theorem~\ref{thm:pack-a}.
Let $D=N^{\theta}$, $z=N^u$, and $Y=D^{\sigma^2}$. Write
$\mathcal P_N=\{p\text{ prime}:p\nmid2N\}$ and
$P_N(Y)=\prod_{p<Y,\ p\in\mathcal P_N}p$. Fix an Iwaniec
sign. Write $\mathcal D^\nu(D)$ for its finite family of nonempty
label vectors $\mathbf D=(D_1,\ldots,D_k)$. They have
$D_k\le\cdots\le D_1<D^{1/2}$, grid coordinates in
$\{D^{\sigma^2\rho^j}:j\ge0\}$, and the cubic conditions
\[
\begin{array}{ll}
 \nu=+:&D_1\cdots D_{2j}D_{2j+1}^3<D
                    \quad(2j+1\le k),\\
 \nu=-:&D_1\cdots D_{2j-1}D_{2j}^3<D
                    \quad(2j\le k).
\end{array}
 \tag{4.7}\label{eq:pack-aa37-cubic}
\]
The prime labels in a vector range over
$D_j\le p_j<\min(D_j^\rho,z)$, with $p_j\nmid2N$. A vector with
an empty coordinate interval contributes zero and may be omitted.

\setcounter{lemma}{\value{theorem}}
\refstepcounter{lemma}\textbf{Lemma~\thelemma\ (Iwaniec coefficient packing).}\label{lem:aa37-pack}
For every contributing vector and
every $\xi\in[60\sigma,1/3+14\sigma]$, let
$e_j=\log_N\min(D_j^\rho,z)$, sorted in non-increasing order.
Then $\mathrm{WF}(d_*)$ holds and
\[
 h\le\widetilde H(\bar e_1,\bar e_2,\bar e_3).
 \tag{4.8}\label{eq:pack-aa37-hyp}
\]
Consequently Theorem~\ref{thm:pack-a} partitions the tail labels into three slots.
Assign the Iwaniec smooth divisor $d_{AA}<D^\sigma$ to slot 1.
Let $Q,S,R$ be the resulting lower dyadic slot scales. Then
\[
\begin{aligned}
 \log_N Q&\le \xi-\left(1+\frac{7271}{256}\right)\sigma,\\
 \log_N(T^2SR^2)&\le1-60\sigma,\\
 \log_N(A_{\max}^{\vartheta}TQS^5R^2)
   &\le2-\left(15+\frac{95221}{1280}\right)\sigma,\\
 \log_N\left((TR/S)^{2\vartheta}TQS^5R^2\right)
   &\le2-\left(15+\frac{22407}{256}\right)\sigma,
\end{aligned}
\tag{4.9}\label{eq:pack-scales}
\]
where $T=N^\xi$ and $A_{\max}=X^{1+\sigma/20}$ for $X\le N$.

\emph{Proof.} Put $s_j=\log_ND_j$. The grid lower bound gives
$s_j\ge\sigma^2\theta$. The cubic conditions in
\eqref{eq:pack-aa37-cubic} imply $\sum_{i<j}s_i+2s_j<\theta$
for every rank $j$. At a cubic rank this follows by dropping one
nonnegative copy of $s_j$. At an intermediate rank,
\[
 \sum_{i<j}s_i+2s_j
 =\sum_{i<j-1}s_i+s_{j-1}+2s_j
 \le\sum_{i<j-1}s_i+3s_{j-1}<\theta,
\]
using $s_{j-1}\ge s_j$ and the cubic condition at rank $j-1$.
For the first rank in the minus case use $2s_1<\theta$, which follows
from $D_1<D^{1/2}$. Thus $\sum_{i<j}e_i+2e_j\le
\rho(\sum_{i<j}s_i+2s_j)<\rho \theta=d_*$.
Also $e_j\le u$, so $\bar e_1\le u<\mu$.
For this range, the branch $w(\bar e_1,\bar e_2,\bar e_3)$ is active.
Its activation threshold is at most $(1+u)/2\le0.501$, while
\[
 \min\left\{\frac{132-68u}{214},
              \frac{153-96u}{256},1-2u\right\}
 =\frac{153-96u}{256}
 =H_{\rm small}(\vartheta,u)=\Theta(u).
\]
Each branch in the minimum decreases with its argument, and
$\bar e_2,\bar e_3\le u$. Take $0\le u\le1/500$. Each other affine branch minus
$(153-96u)/256$ is nonnegative at both endpoints. Hence it is
nonnegative on the whole interval. In particular, the
minimum defining $w(\bar e_1,\bar e_2,\bar e_3)$ is at least
$\Theta(u)>0.501$, so its activation condition holds. Thus
$\widetilde H\ge\Theta(u)$. Next,
$\tau \theta+60\sigma<\eta_\Theta$, because
$60\sigma\le0.3\eta_\Theta$ and
$\tau \theta<0.6\sigma^9\le0.6\sigma<0.7\eta_\Theta$.
Therefore
\[
 \frac12<h=\rho \theta+60\sigma
 =\Theta(u)-\eta_\Theta+\tau \theta+60\sigma
 <\Theta(u)<\overline H.
\]
Also $0<\delta\le3\cdot10^{-9}$. The range for $\xi$ gives
$g=\xi-\delta\in[\delta,1/3-\delta/2]$.
Theorem~\ref{thm:pack-a} applies.

The smooth-only term is the empty-list case. Here
$\widetilde H(0,0,0)=H(0,0,0)=\overline H$, since the $w$ branch
equals $\min\{132/214,153/256,1\}=\overline H$ and its activation
threshold is $1/2$. Thus Theorem~\ref{thm:pack-a} applies to the empty list as well.
The smooth divisor remains in slot 1, and slots 2 and 3 are units.

For the scale bounds, the slot-1 smooth divisor contributes at most
$\sigma \theta<\sigma\overline H$ to the logarithmic load. The other
slot loads are bounded by $\Sigma_2,\Sigma_3$. Thus
\[
\begin{aligned}
 \log_NQ&\le\Sigma_1+\sigma\overline H,&
 \log_NS&\le\Sigma_2,&
 \log_NR&\le\Sigma_3,\\
 \log_NT&=\xi=g+\delta,&
 \log_NA_{\max}&\le1+\sigma/20.
\end{aligned}
\]
The first line of \eqref{eq:pack-loads} gives
\[
 \log_NQ-\xi\le-(30-\overline H)\sigma
 =-\left(1+\frac{7271}{256}\right)\sigma.
\]
The second gives
\[
 \log_N(T^2SR^2)\le1-2\delta=1-60\sigma.
\]
For the third scale, \eqref{eq:pack-loads} gives
\[
 \log_N(A_{\max}^{\vartheta}TQS^5R^2)
 \le2-3\delta+\left(\frac{\vartheta}{20}+\overline H\right)\sigma
 =2-\left(15+\frac{95221}{1280}\right)\sigma.
\]
For the fourth, it gives
\[
 \log_N((TR/S)^{2\vartheta}TQS^5R^2)
 \le2-(3+2\vartheta)\delta+\overline H\sigma
 =2-\left(15+\frac{22407}{256}\right)\sigma.
\]
The exact constants in \eqref{eq:pack-scales} hold for the lower scales.
We round each of the three aggregate slot scales $Q,R,S$ once, to a
dyadic representative within a factor two. Each monomial below has
total degree at most $8$ in $Q,R,S$. So rounding multiplies it by at
most $2^8$. Individual prime labels are not rounded.
This adds at most
$8\log 2/\log N$ to its base-$N$ exponent. Fix $\sigma>0$. For sufficiently large $N$, each positive margin
absorbs this cost. The tightest margin, for $Q<TX^{-29\sigma}$, is
$103\sigma/256-8\log 2/\log N>0$. \hfill$\square$

For fixed $B>0$ and $N/(2(\log N)^B)\le X\le N$, put
$\omega=\log_N(N/X)$. Then
$\omega\le(B\log\log N+\log2)/\log N$, so eventually
$\log_XN=(1-\omega)^{-1}<1+0.01\sigma$. The first exponent in
\eqref{eq:pack-scales}, after dividing by $T$, is negative. So
conversion from base $N$ to base $X$ improves that relative inequality.
Each other exponent is at most $2$;
conversion costs at most $0.02\sigma$. The displayed margins absorb
this cost.

For every contributing vector, the last-rank mask gives
$\sum_js_j<\theta$. Since each $s_j\ge\sigma^2\theta$, it also gives
$k<\sigma^{-2}$. Thus $\prod_jD_j<D$, and $p_j<D_j^\rho$ gives
\[
 d\prod_jp_j<D^\sigma D^\rho=D^{1+\sigma+\tau}
 =N^{(1+\sigma+\tau)\theta}<N^{\Theta(u)},
\]
where $(\sigma+\tau)\theta<\eta_\Theta$.
All prime factors avoid $2N$. Thus the full coefficient is supported
on $q<N^{\Theta(u)}$ with $(q,2N)=1$. Its coefficients are
$O_\sigma(1)$. The $Y$-smooth part of a supported modulus determines
$d$. The remaining labelled prime assignments number at most $k!$,
and $k<\sigma^{-2}$. For each sign the Iwaniec family has fewer than
$e^{8\sigma^{-3}}$ vectors. So the sum over vectors keeps a bound
depending only on $\sigma$. The smooth-only term satisfies the same
support and coefficient bounds.

The cubic-to-quadratic implication just proved is uniform in the vector
length. We record the two bounds obtained above:
\[
 k<\sigma^{-2},\qquad \sum_{j=1}^k s_j<\theta.
 \tag{4.10}\label{eq:pack-aa37-length}
\]
These bounds also control the coefficient regrouping below.

\subsection{Three-slot coefficient identity}

Fix a sign $\nu$ and a vector $\mathbf D$. The set of
equal-label pairs in the Iwaniec grid is
$I(\mathbf D)=\{(i,j):1\le i<j\le k,\ D_i=D_j\}$. Distinct grid labels have
disjoint half-open prime intervals
$[D_i,\min(D_i^\rho,z))$. Adjacent grid values differ by
the power $\rho$. Thus repeated primes can occur only in
$I(\mathbf D)$. Define
\[
 K_{\mathbf D}(\mathbf p)=
   \prod_{(i,j)\in I(\mathbf D)}(1-\mathbf1_{p_i=p_j}).
\]
For the Iwaniec coefficients $\lambda_d^\nu(\mathbf D)$, define the
unrestricted coefficient
\[
 C_{\mathbf D}^{\rm full}(q)=
 \sum_{\substack{d<D^\sigma\\d\mid P_N(Y)}}
 \lambda_d^\nu(\mathbf D)
 \sum_{\substack{D_j\le p_j<\min(D_j^\rho,z)\\p_j\nmid2N}}
 \mathbf1_{q=d\prod_{j=1}^k p_j}.
 \tag{4.11}\label{eq:pack-aa37-vector}
\]
The masked coefficient $C_{\mathbf D}^{\rm mask}$ is defined by
inserting $K_{\mathbf D}(\mathbf p)$ in the inner sum. The coefficient
of the smooth-only term, the sole empty-tail term in the Iwaniec formula,
is
$C_{\emptyset}(q)=\sum_{d<D^\sigma,d\mid P_N(Y)}
\phi_d^\nu\mathbf1_{q=d}$. Thus the full signed coefficient is
\[
 w_N^\nu(q)=C_{\emptyset}(q)+
       \sum_{\mathbf D\in\mathcal D^\nu(D)}
             C_{\mathbf D}^{\rm mask}(q).
 \tag{4.12}\label{eq:pack-aa37-weight}
\]

\refstepcounter{lemma}\textbf{Lemma~\thelemma\ (slot factorization and collision support).}\label{lem:aa37-slot}
Fix $\xi$ in
the range of Lemma~\ref{lem:aa37-pack} and let $J_1\sqcup J_2\sqcup J_3$ be its
packing partition. All $d$ and $p_j$ sums below use the ranges in
\eqref{eq:pack-aa37-vector}. Define
\[
\begin{aligned}
 \gamma_1(m)&=\frac1{|J_1|!}
    \sum_{d,(p_j)_{j\in J_1}}
       \lambda_d^\nu(\mathbf D)\mathbf1_{m=d\prod_{j\in J_1}p_j},\\
 \gamma_s(m)&=\frac1{|J_s|!}
    \sum_{(p_j)_{j\in J_s}}
       \mathbf1_{m=\prod_{j\in J_s}p_j},\qquad s=2,3.
\end{aligned}
\]
For $s=2,3$ set $\gamma_s=\mathbf1_{m=1}$ when $J_s$ is empty;
$\gamma_1$ always retains the smooth-divisor sum. Then
\[
 C_{\mathbf D}^{\rm full}(q)=C_{\rm slot}
       (\gamma_1*\gamma_2*\gamma_3)(q),\qquad
 C_{\rm slot}=|J_1|!|J_2|!|J_3|!,
 \tag{4.13}\label{eq:pack-aa37-slot-id}
\]
and $|\gamma_s(m)|\le1$. Consider the full Iwaniec family. The difference between the masked
and unrestricted coefficients has pointwise size $O_\sigma(1)$. It is
supported on moduli divisible by $p^2$ for some prime $p\ge Y$.

\emph{Proof.} A labelled tuple $(d,(p_j)_{j\le k})$ splits uniquely
into its three slot tuples. Multiplying the three normalized finite
sum definitions by $C_{\rm slot}$ therefore gives
\eqref{eq:pack-aa37-slot-id}. This regroups a finite sum and
allows repeated primes. Fix a product $m$. Its prime factors
below $Y$ determine $d$ uniquely: $d\mid P_N(Y)$, and every tail
prime is at least $D_j\ge Y$. Once $d$ is fixed, the labelled prime
assignments within a slot number at most $|J_s|!$. The Iwaniec bounds
$|\lambda_d^\nu(\mathbf D)|\le1$ give $|\gamma_s(m)|\le1$;
the same argument applies to $\gamma_1$. The smooth-only term has
the corresponding bound from $|\phi_d^\nu|\le1$. Since
$k<\sigma^{-2}$, the scalar
$C_{\rm slot}=|J_1|!|J_2|!|J_3|!$ is $O_\sigma(1)$.

Suppose the distinctness mask removes a tuple. Then some selected pair
has $p_i=p_j\ge Y$, and its modulus is divisible by $p_i^2$. For fixed
$\sigma$, \eqref{eq:pack-aa37-length} bounds the number of labels.
The Iwaniec family has fewer than $e^{8\sigma^{-3}}$ vectors for each
sign. The number of labelled
assignments to a fixed modulus is at most $k!$. Hence the aggregate
coefficient difference is $O_\sigma(1)$.

The prime-source cost of this collision support is estimated in
(B.5)--(B.6). \hfill$\square$

The smooth-only term $C_{\emptyset}$ is already a one-slot coefficient.
Take $J_1=J_2=J_3=\emptyset$, put
$\gamma_1(m)=\sum_{d<D^\sigma,d\mid P_N(Y)}
\phi_d^\nu\mathbf1_{m=d}$, and
$\gamma_2=\gamma_3=\mathbf1_{m=1}$. Its coefficient identity and
pointwise bound follow directly. The full coefficient $w_N^\nu$ in
\eqref{eq:pack-aa37-weight} is the sum of this term and the vector
terms above. Every partition is fixed
before the prime realizations are summed.

This section proves the packing inequalities and the slot identity for
the unrestricted coefficient. It also gives the collision support
needed in Appendix~B.
The prime-source estimate is proved there.

\section{Application to Iwaniec's remainder weights}
\label{sec:linear-sieve-application}

The source is the prime interval
$I_N=(N/(\log N)^B,(1-\varepsilon)N)$ from
Theorem~\ref{thm:aa37-interior}. The residue satisfies
$N\le a\le N\log N/3$ and is coprime to every supported modulus.
The saving is $(\log N)^{-A}$ for each fixed $A>0$.

\subsection{Iwaniec's remainder coefficients}

All parameters are those in Theorem~\ref{thm:aa37-interior}; in
particular $s=\log D/\log z=\theta/u$.
Let $N>2$ be even. Let $\mathcal P_N$ be the set of primes not
dividing $2N$, and write
\[
 P_N(t)=\prod_{\substack{p<t\\p\in\mathcal P_N}}p.
\]
Use the grid $\mathcal G(D)$ in \eqref{p1:eq:2.28}. The admissible
vector sets $\mathcal D^\nu(D)$ are those in \eqref{p1:eq:2.29}.
The parameter substitution is \eqref{p1:eq:2.33a}--\eqref{p1:eq:2.33b},
with $\epsilon_s=\sigma$ and $\delta_s=\tau$. The distinctness set is
$I(\mathbf D)=\{(i,j):1\le i<j\le k,\ D_i=D_j\}$, the equal-label
tail-position pairs. For a vector $\mathbf p=(p_1,\ldots,p_k)$ define
\[
 K_{\mathbf D}(\mathbf p)
   =\prod_{(i,j)\in I(\mathbf D)}
       (1-\mathbf 1_{p_i=p_j}).
\]
Let $\phi_d^\nu$ and $\lambda_d^\nu(\mathbf D)$ be the coefficients of
Iwaniec \cite[Theorem~4]{AA37}. Take $R=1$ and parameters
$(\mathcal P_N,D,z,\sigma)$ at the cutoff $z=N^u$. The proof uses only
their bound $|\cdot|\le1$ and their support. It never compares
coefficient values at different cutoffs.
The labelled-tail cutoff of C1 in Section~3, used in (H1), is
$u_L=Y=D^{\sigma^2}$.
Since $\theta>1/2>\ell_0$, this satisfies
$u_L=N^{\theta\sigma^2}\ge N^{\ell_0\sigma^2}=u_0$.
The smooth divisor has $d<D^\sigma$ and $d\mid P_N(Y)$. Every
tail interval starts at $D_i\ge Y$.
The exact structured remainder coefficients are
\begin{equation*}\tag*{(W)}
\begin{split}
 w_N^\nu(q)
 ={}&\sum_{\substack{d<D^\sigma\\d\mid P_N(Y)}}
          \phi_d^\nu\mathbf1_{q=d}\\
 &+\sum_{\mathbf D\in\mathcal D^\nu(D)}
   \sum_{\substack{d<D^\sigma\\d\mid P_N(Y)}}
       \lambda_d^\nu(\mathbf D)
   \sum_{\substack{p_i\in\mathcal P_N\\
             D_i\leq p_i<\min(D_i^{1+\tau},z)}}
       K_{\mathbf D}(\mathbf p)\,
       \mathbf1_{q=dp_1\cdots p_k}.
\end{split}
\end{equation*}
In particular, $\nu\in\{+,-\}$ selects one of the two signed
remainders. Each is the full remainder, not an arbitrary
well-factorable sequence. Formula~\textup{(W)} includes the
smooth-only term. Every prime-tuple term keeps its smooth divisor. The smooth-only term
retains its divisor when there is no tail.
The estimate below applies separately to either signed remainder.
When the directed lower-sieve inequality of Iwaniec
\cite[Theorem~4]{AA37} is used, it takes $\nu=-$.

\subsection{Proof of the interior signed mean value}

\begin{proof}[Proof of Theorem~\ref{thm:aa37-interior}]
The support and coefficient bounds in the theorem are proved in
(B.2), using Iwaniec's cubic label conditions for every vector in
\textup{(W)}. In particular, with $e_Q=(1+\sigma+\tau)\theta$,
\[
 q<D^{1+\sigma+\tau}=N^{e_Q}<N^{\Theta(u)},
 \qquad (q,2N)=1,
\]
including the smooth-only term. The choice
$\sigma\le(\Theta(u)-\theta)/200$ gives the strict second inequality.
Take the source $\mu_D=\mathbf1_{\mathbb P}\mathbf1_{I_N}$ of (B.1).
By (B.3), $p>q$ for every source prime $p$ and every supported
modulus $q$. The coprimality $(a,q)=1$ is assumed. So the natural
comparator in (B.1) is exactly $\pi_{I_N}/\varphi(q)$, and
\[
 \Delta_D(q;a)=\pi_{I_N}(q,a)-\frac{\pi_{I_N}}{\varphi(q)}.
\]
Before applying the dispersion theorem, (B.5)--(B.6) replace
$w_N^\nu$ by
$\widehat w_N^\nu(q)=\mathbf1_{(q,a)=1}\widetilde w_N^\nu(q)$,
where $\widetilde w_N^\nu$ has the distinctness masks removed.
By the hypothesis on the support of $w_N^\nu$, the difference
$w_N^\nu-\widehat w_N^\nu$ is a collision contribution. By (B.5),
its complete signed discrepancy is $O_\sigma(N\log N/Y)$.
The coprimality projection $\mathbf1_{(q,a)=1}$ factors across the
three slots in (4.13), as shown after (B.6); this includes the
smooth-only term. Thus the retained coefficient has the exact
slot form (3.5). All its moduli are coprime to $a$.
Appendix~\ref{app:proof-aa37} extends the support envelope
$q<N^{e_Q}$ of (B.2) to $\widetilde w_N^\nu$ and $\widehat w_N^\nu$.
For each source scale $X>N/(2(\log N)^B)$, (B.3) then places this entire
retained support below $X^{0.6}$. The residue bound gives
$a/X<(2/3)(\log N)^{B+1}$. Hence C4 holds with
$K_{\rm scale}=B+1$. This includes $a\le X^{1+\sigma/20}$ for large
$N$. On the doubled complementary support, (B.28) gives
$a-mn\ge\varepsilon N/2$. It uses $a\ge N$ and the upper source cut. Since $a\ll N\log N$ and $mn\ll N$ there, also
$|mn-a|\le X^3$ for large $N$. These are both parts of the
nonresonance condition (3.10).

After the collision projection, the prime-power contribution is
(B.8). The logarithm-freezing error is (B.18). The source-boundary
error is (B.19). Lemma~\ref{lem:aa37-typeII-lengths} proves
the selected-interval lengths and integer endpoint convention in
(B.22)--(B.27). Lemma~\ref{lem:aa37-pack} supplies the slot inequalities, while
(B.20)--(B.21) classify every retained box. The internal Type-I and
Type-II bounds, including the box and arithmetic costs, give (B.31).
The integer choice
$C_{\rm box}\ge\lceil2A+32\rceil$
makes the
logarithm-freezing and boundary errors in (B.18)--(B.19)
$O(N(\log N)^{-A})$. The aggregate exponent
$A_{\rm agg}=A$ here is the theorem convention; the separate
$A_{\rm agg}=A+3$ bookkeeping used for (3.95) requires
$C_{\rm box}\ge\lceil2A+38\rceil$ and is outside Theorem~\ref{thm:main}.
$A'=A+24C_{\rm box}+31$ and $A_{\rm work}=A'+4$ absorb the box count.
Appendix~\ref{app:proof-aa37} verifies the box-family bound (H7), that
is (3.91) with $K_0=24C_{\rm box}+26$. The aggregate form of Theorem~\ref{thm:main}, namely (3.92),
bounds the Type-II boxes. For each fixed admissible $\sigma>0$, the mask and proper-power
terms have fixed power savings.
Proposition~\ref{prop:aa37-prime-source} combines these terms over
the complete support of \textup{(W)}. This gives the assertion for
each fixed $A,B,\varepsilon,u,\theta,\sigma$ in the stated ranges.
\end{proof}

\subsection{Proof of the logarithmic-integral comparator}

\begin{proof}[Proof of Corollary~\ref{cor:aa37-interior-li}]
Put $L=\log N$. The coefficient bound in
Theorem~\ref{thm:aa37-interior} and
$\sum_{q\le x}1/\varphi(q)\ll\log(2x)$ give
$\sum_q|w_N^\nu(q)|/\varphi(q)\ll_{\sigma,u,\theta} L$.
The prime number theorem at the two endpoints of $I_N$ gives
\[
 \pi_{I_N}=\operatorname{Li}((1-\varepsilon)N)
              -\operatorname{Li}(N/L^B)
              +O_A(NL^{-A-1})
\]
for every fixed $A>0$. Strict endpoint conventions contribute at most
two primes. Multiply this error by the reciprocal-totient sum and apply
\eqref{eq:aa37-interior-mean}.
\end{proof}

\appendix
\section{Source checks}\label{app:finite-source-checks}

The source input used for the coefficients in formula~\textup{(W)}
is \cite[Theorem~4]{AA37}. Its hypotheses are \cite[(1)--(2)]{AA37}.
Its upper inequality permits $2\le z\le D$;
its lower inequality permits $2\le z\le D^{1/2}$.  Fix $0<u\le 1/500$ and $1/2<\theta<\Theta(u)$ as in
Theorem~\ref{thm:aa37-interior}. Then $D=N^\theta$ and $z=N^u$ satisfy
$2u<\theta$.  Thus both ranges apply for all sufficiently large $N$.
The choice $\sigma\le\sqrt{u/(2\theta)}$ also gives
$Y=D^{\sigma^2}\le N^{u/2}<z$.  The mesh is
$\epsilon_s=\sigma$ and $\delta_s=\sigma^9=\tau$. The
smooth-prime cutoff is $Y$, as in
\eqref{p1:eq:2.28}--\eqref{p1:eq:2.33b}.

Theorem~4's printed argument list of $\lambda_d^\nu$ contains the
mesh parameter and box tuple. It contains neither the sequence nor the
cutoff.
The construction in \cite[(23)--(26)]{AA37} collects universal
Rosser and fundamental-lemma coefficients.  The sequence enters through
$r(\mathcal A,q)$.  At each cutoff $z$, we use the coefficients and
prime-tuple support supplied by that construction. The coefficient
family is independent of the sequence. It may depend on $z$.
The signed prime-source estimate is proved
separately in Appendix~\ref{app:proof-aa37}.

The $O$-constant in the main-term error of \cite[Theorem~4]{AA37}
is not given numerically.

The proof of Theorem~\ref{thm:aa37-interior} uses no finite computation.

\section{Prime-source estimate for Iwaniec's remainder weights}
\label{app:proof-aa37}
\begingroup
\setcounter{lemma}{0}
\renewcommand{\thelemma}{B.\arabic{lemma}}

This section proves the source estimate used in Section~5. Equations
(B.2)--(B.6) control the coefficients and distinctness masks.
Equations (B.7)--(B.19) treat prime detection, logarithm freezing and
boundary boxes. Equations (B.20)--(B.32) bound the retained boxes.

Fix the parameters in Theorem~\ref{thm:aa37-interior} and put
\[
 \begin{gathered}
 \tau=\sigma^9,\quad
 e_Q=(1+\sigma+\tau)\theta<\Theta(u)\le\frac{153}{256},\\
 D=N^{\theta},\quad Y=D^{\sigma^2}.
 \end{gathered}
\]
Put $s=\log D/\log z=\theta/u$, $z=N^u$, $L=\log N$,
$H=N/L^B$, and use the width $C_{\rm box}$ in
Theorem~\ref{thm:aa37-interior}. Write
$\mathcal S_n(q;a)=\mathbf1_{n\equiv a\ (q)}
-\mathbf1_{(n,q)=1}/\varphi(q)$.
All thresholds below may depend on the fixed parameters. For a product
shell $[X,2X)$, set $L_X=\log X$. The source measure is
\[
 \begin{aligned}
 \mu_D(n)&=\mathbf1_{\mathbb P}(n)
          \mathbf1_{H<n<(1-\varepsilon)N},\\
 \Delta_D(q;a)&=\sum_n\mu_D(n)
 \left(\mathbf1_{n\equiv a\ (q)}
       -\frac{\mathbf1_{(n,q)=1}}{\varphi(q)}\right).
 \end{aligned}
 \tag{B.1}
\]

\begin{lemma}[Modulus support and coefficient size]\label{lem:aa37-support}
For either sign in \textup{(W)},
\[
 \operatorname{supp}w_N^\nu\subset
 \{q<N^{e_Q}:(q,2N)=1\},\qquad
 \sup_q|w_N^\nu(q)|\le C_{\sigma,u,\theta}.
 \tag{B.2}
\]
If $[X,2X)$ meets $(H,(1-\varepsilon)N)$, then, for all sufficiently large $N$,
\[
 q<N^{e_Q}\le X^{3/5}
 \quad(q\in\operatorname{supp}w_N^\nu),\qquad H>N^{e_Q}.
 \tag{B.3}
\]
\end{lemma}

\emph{Proof.} By the definition of $\mathcal D^\nu(D)$ in \eqref{p1:eq:2.29},
every label vector in (W) satisfies the cubic label conditions of Iwaniec's
sieve remainder \cite[Theorem~4]{AA37}. Write
$s_i=\log_ND_i$ for a label vector
$\mathbf D=(D_1,\ldots,D_k)$. The cubic conditions in Iwaniec imply
$\sum_i s_i<\theta$. For the plus sign and odd $k$, use the last odd-rank condition.
For even $k$, use the condition at rank $k-1$
and $s_k\le s_{k-1}$. For the minus sign, if $k=1$ then
$s_1<\theta/2$. If $k$ is even, use the last even-rank condition.
If $k\ge3$ is odd, use rank $k-1$ and $s_k\le s_{k-1}$. Thus every
term of (W) satisfies
\[
 q<N^{\sigma \theta+(1+\tau)\sum_i s_i}
   <N^{\theta(1+\sigma+\tau)}=N^{e_Q}.
\]
All prime factors in $d$ and in the tail avoid $2N$.

Each grid coordinate has the form $D^{\sigma^2(1+\tau)^h}$ and is
less than $D^{1/2}$. Thus there are at most
$1+2\tau^{-1}\log(\sigma^{-2})$ choices per position, since
$\log(1+\tau)\ge\tau/2$. Also $k<\sigma^{-2}$ by the
last-rank cubic condition and $s_i\ge\sigma^2\theta$. With
$\tau=\sigma^9$ and $0<\sigma\le10^{-10}$, the total number of
vectors of all lengths is at most
\[
 \sum_{0\le k<\sigma^{-2}}
  \bigl(1+2\sigma^{-9}\log(\sigma^{-2})\bigr)^k
 <\exp(8\sigma^{-3}).
\]
The logarithm of the left side is at most
$\log(1+\sigma^{-2})+
\sigma^{-2}\log(1+4\sigma^{-9}\log(1/\sigma))$.
With $t=1/\sigma\ge10^{10}$, this is at most
$(3+11t^2)\log t<8t^3$, using $\log t\le t/2$.
For fixed $q$, the primes below $Y$
determine the smooth divisor $d$ uniquely. The remaining prime factors
can be assigned to the $k$ labelled positions in at most $k!$ ways.
The Iwaniec coefficients have absolute value at most $1$.
Summing over the displayed finite vector family proves the
multiplicity bound, including the empty-tail term.

If the shell meets the source interval, then $X>H/2$ and
$\log X>L-B\log L-\log2$. Impose
\[
 (3/5-e_Q)L\ge(3/5)(B\log L+\log2).
 \tag{B.4}
\]
It gives $e_QL\le(3/5)\log X$, proving the first claim
of (B.3). The second follows from $N^{1-e_Q}/L^B\to\infty$.
\hfill$\square$

Here $p>H>N^{e_Q}$, and the theorem assumes $(a,q)=1$ on the
coefficient support. So the natural comparator in (B.1) is
$\pi_{I_N}/\varphi(q)$. Thus (B.1) is the exact prime-source
discrepancy for this source.

\subsection*{Removing the distinctness masks}

Let $\widetilde w_N^\nu$ be obtained from (W) by replacing each
$K_{\mathbf D}$ by $1$, and put
$\widehat w_N^\nu(q)=\mathbf1_{(q,a)=1}\widetilde w_N^\nu(q)$.
The counting argument in Lemma~\ref{lem:aa37-support} also gives
$\sup_q|\widetilde w_N^\nu(q)|+\sup_q|\widehat w_N^\nu(q)|
\ll_{\sigma,u,\theta}1$. Its termwise product bound applies to every
unmasked term. So $\widetilde w_N^\nu$ and $\widehat w_N^\nu$ are
also supported on $\{q<N^{e_Q}:(q,2N)=1\}$, and (B.3) applies to them.
Their support need not lie in $\operatorname{supp}w_N^\nu$.
The theorem assumes $w_N^\nu(q)=0$ when $(q,a)>1$. Hence
$w_N^\nu-\widehat w_N^\nu$ is supported on the terms removed by a
distinctness mask. This holds even when an unrestricted collision
modulus is not coprime to $a$. Every such modulus is divisible by
$p^2$ for some tail prime $p\ge Y$. The coefficient difference is
$O_\sigma(1)$. Put $Q_{\rm supp}=N^{e_Q}$.
\[
 \sum_{\substack{q\le Q_{\rm supp}\\p^2\mid q\text{ for some prime }p\ge Y}}
 \left(\frac Nq+1+\frac N{\varphi(q)}\right)
 \ll \frac{N\log(2Q_{\rm supp})}{Y-1}.
 \tag{B.5}
\]
For the proof, write $q=kp^2$, use
$\varphi(kp^2)\ge\varphi(k)\varphi(p^2)$, and sum
$1/k$ and $1/\varphi(k)$ up to $Q_{\rm supp}$. Then use
$\sum_{p\ge Y}1/(p(p-1))\ll1/(Y-1)$. The left side bounds the
discrepancy of each removed term. Its progression count is at
most $N/q+1$. Consequently
\[
 \sum_q w_N^\nu(q)\Delta_D(q;a)
 =\sum_q\widehat w_N^\nu(q)\Delta_D(q;a)
   +O_\sigma(NL/Y).
 \tag{B.6}
\]
As $Y=N^{\theta\sigma^2}$, this error has a fixed power saving for each
fixed $\sigma>0$.
For a fixed Iwaniec vector, the exact factorization (4.13) survives
the projection. Replace each slot factor $\gamma_i(m)$ by
$\gamma_i(m)\mathbf1_{(m,a)=1}$. Here
$(m_1m_2m_3,a)=1$ if and only if each $(m_i,a)=1$.
The bounds $|\gamma_i|\le1$ are unchanged. Equivalently, the smooth
divisor and each tail prime are restricted to avoid $a$ within their
existing summation ranges. Thus the projected coefficient still has
the labelled form (3.4). Every modulus to which the dispersion
theorem is applied is coprime to $2a$.

\subsection*{Prime detection and the rough Heath--Brown identity}

The shells meeting the source satisfy $H/2<X<N$. Set
\[
 z_0=z_0(X)=X^{1/(\log\log X)^3},\qquad U=2X^{1/3}.
\]
For sufficiently large $N$, every $n\in[H,(1-\varepsilon)N]$ has $n\ge z_0$.
Write $\Lambda_{z_0}(n)=\Lambda(n)\mathbf1_{P^-(n)\ge z_0}$.
For every integer $n\ge z_0$, and so on the whole source range,
\[
 \mathbf1_{\mathbb P}(n)=\frac{\Lambda_{z_0}(n)}{\log n}
 -\sum_{k\ge2}\frac1k\mathbf1_{n=p^k,\ p\ge z_0}.
 \tag{B.7}
\]
The identity follows by checking primes, prime powers, and all other
integers $n\ge z_0$. It fails for primes below $z_0$, which lie
outside the source range. The proper powers contribute to the weighted discrepancy
\[
 O_\sigma(N^{1/2+o(1)}),
 \tag{B.8}
\]
in total. There are $O(N^{1/2}\log N)$ such integers in the source
range $H<n<(1-\varepsilon)N$. The theorem's range
$N\le a\le NL/3$ gives $a\ne n$ there. Each contributes
at most $O_\sigma(d(a-n)+L)$, where $d(m)$ is the divisor function.

For $n\le 8X$, the order-three rough Heath--Brown identity is
\[
 \Lambda_{z_0}(n)=\sum_{j=1}^3(-1)^{j-1}\binom3j
 \sum_{\substack{m_1\cdots m_jn_1\cdots n_j=n\\
                   m_i\le U,\ m_i,n_i\in\mathcal R_{z_0}}}
 \mu_{z_0}(m_1)\cdots\mu_{z_0}(m_j)\log n_1.
 \tag{B.9}
\]
Here $\mathcal R_{z_0}$ consists of $1$ and the $z_0$-rough integers,
and $\mu_z(m)=\mu(m)\mathbf1_{P^-(m)\ge z}$.
To verify (B.9), put
\[
 \zeta_z(s)=\sum_{r\in\mathcal R_z}r^{-s},\qquad
 M_U(s)=\sum_{m\le U}\mu_z(m)m^{-s}.
\]
The coefficients of $\zeta_zM_U$ vanish except at $1$ through degree
$U$. Thus all coefficients of $(1-\zeta_zM_U)^3$ at $n\le U^3$ vanish. Multiply the expansion of $1-(1-\zeta_zM_U)^3$ by
$-\zeta_z'/\zeta_z$ to obtain (B.9). Every variable is rough; this
includes the logarithmic variable $n_1$.

The tested indices for shell $X$ and pattern $j$ are
\[
 T_{X,j}=\left\{
 \begin{aligned}
 \mathbf v&=(m_1,\ldots,m_j,n_1,\ldots,n_j)
             \in\mathcal R_{z_0}^{2j}:\\
 &X\le\prod_i v_i<2X,\quad
   H<\prod_i v_i<(1-\varepsilon)N,\\
 &m_i\le U
 \end{aligned}\right\}.
 \tag{B.10}
\]
There are at most six variables. Their coefficient is
\[
 c_j(\mathbf v)=(-1)^{j-1}\binom3j
 \mu_{z_0}(m_1)\cdots\mu_{z_0}(m_j)\log n_1.
\]
The shell, pattern, and tuple are part of the index. The identity's
range $n\le8X$ is automatic on $T_{X,j}$.

For every nonnegative $G$ and each $j\in\{1,2,3\}$,
\[
 \sum_X\sum_{\mathbf v\in T_{X,j}}G\!\left(\prod_i v_i\right)
 \le\sum_{n<N}\tau_{2j}(n)G(n)
 \le\sum_{n<N}\tau_6(n)G(n).
 \tag{B.11}
\]
The product determines its shell. Roughness restrictions can be
dropped in this upper bound. Consequently the sum over all three
patterns is at most $3\sum_{n<N}\tau_6(n)G(n)$. The factor $3$
is absorbed in later implied constants.

For $n<N$, (B.2) and its extension to $\widehat w_N^\nu$ after (B.5)
give the pointwise majorant
\[
\left|\sum_q\widehat w_N^\nu(q)\mathcal S_n(q;a)\right|
 \le C_\sigma F_a(n),\qquad
 F_a(n)=d(a-n)+2(1+L).
 \tag*{}
\]
The congruence term counts $q\mid a-n$. The natural-comparator term
uses $\sum_{q\le N^{e_Q}}1/\varphi(q)\ll L$. On the prime source,
(B.1) equals $\sum_n\mu_D(n)\mathcal S_n(q;a)$. Apply (B.7) to this
identity while keeping $\mathcal S_n(q;a)$ attached to the tested
integer. Each HB tuple then contributes the natural PSDT discrepancy
$\mathcal S_{\prod_i v_i}(q;a)$. This holds also when its product
shares factors with $q$. The proper-power terms use the same discrepancy and
obey (B.8).

\subsection*{Equal-ratio boxes and faces}

For each product shell put
\[
 \delta_X=L_X^{-C_{\rm box}},\qquad
 J_X=\left\lceil\frac{\log2}{\log(1+\delta_X)}\right\rceil,
 \qquad \Delta_X=2^{1/J_X}-1.
\]
For sufficiently large $N$, $\delta_X\le1/3$. The equal-ratio
partition divides every dyadic coordinate interval $[2^a,2^{a+1})$
into the $J_X$ cells
\[
 [2^a(1+\Delta_X)^h,\ 2^a(1+\Delta_X)^{h+1}),
 \qquad 0\le h<J_X.
\]
It covers the full interval, including its last cell. Also
\[
 \frac{\delta_X}{3}\le\delta_X(1-2\delta_X)
 \le\Delta_X\le\delta_X.
 \tag{B.12}
\]
For the lower bound, set $\ell=\log(1+\delta_X)$. Since
$J_X<\log2/\ell+1$,
\[
 \log(1+\Delta_X)>\frac{\ell}{1+\ell/\log2}
 \ge\ell(1-\ell/\log2).
\]
Use $\delta_X-\delta_X^2/2\le\ell\le\delta_X$ and
$1/2+1/\log2<2$. This proves (B.12).

The mesh is imposed on the ambient integer axes before roughness is
tested. Keep $\{1\}$ as an atom. A product box is a Cartesian product
of at most six cells. Let $V_i$ be its lower endpoints and put
$e_i=\log_N V_i$; the atom has $V_i=1$. Define
$\lambda_{\mathcal B}=\sum_i\log V_i$. Its tested-side faces are
\[
 n=H,\quad n=(1-\varepsilon)N,\quad n=X,\quad n=2X,\quad m_i=2X^{1/3},
 \tag{B.13}
\]
where $n=\prod_i v_i$ and $m_i$ ranges over the Möbius variables.
The first two are the source cuts. The next two are product-shell
edges, and the last are the HB caps. Roughness is part of the
coefficient, not a face. A box is called boundary if one predicate in
(B.13) is not constant on it. Every other box is either empty or a
full Cartesian sum.

\begin{lemma}[Lattice thin-shell count]\label{lem:aa37-thin-shell}
Let $s\ge2$,
$a\in\mathbb Z^s\setminus\{0\}$ with $|a_i|\le2$, and
$J=\{i:a_i\ne0\}$. Let $1\le Y_0\le N$, $0<\eta_{\rm dist}\le1/4$, and let
$\mathcal L$ be $\gamma$-separated. Put
$K_\gamma=(1-e^{-\gamma/2})^{-1}$. Then
\[
\begin{split}
 \#\{\mathbf v\in\mathbb Z_{\ge1}^s:&\ \prod_i v_i\le N,
   \operatorname{dist}(\sum_i a_i\log v_i,\mathcal L)\le\eta_{\rm dist},
   \max_{i\in J}v_i\ge Y_0\}\\
 &\le s\left[4\eta_{\rm dist} K_\gamma N(1+L)^{s-1}
   +\left(\frac{4L}{\gamma}+2\right)
      \frac N{Y_0}(1+L)^{s-2}\right].
\end{split}
 \tag{B.14}
\]
For a single face, take $K_\gamma=1$ and replace
$4L/\gamma+2$ by $1$.
\end{lemma}

\emph{Proof.} Assign a tuple to the least active coordinate $i$ with
$v_i\ge Y_0$. Fix the other coordinates, with product $P$, and put
$Z=N/P$. The nonempty fibre has $P\le N/Y_0$. Let $q=|a_i|\in\{1,2\}$. The admissible values of $v_i$ lie in
intervals of relative length at most $e^{2\eta_{\rm dist}/q}-1\le4\eta_{\rm dist}$.
Their logarithmic left endpoints have spacing at least $\gamma/q$.
After clipping to $[Y_0,Z]$ and merging overlaps, there are at most
$2L/\gamma+2$ components. The separated left endpoints give total
length at most $4\eta_{\rm dist} K_\gamma Z$. Counting integer endpoints adds
at most $2L/\gamma+2$. For a singleton face there is one component.
Since $2L/\gamma+2\le4L/\gamma+2$, this gives the coefficient
$4L/\gamma+2$ in (B.14).
Sum over the fixed coordinates using
\[
 \sum_{P\le N}\frac{\tau_{s-1}(P)}P\le(1+L)^{s-1},\qquad
 \sum_{P\le N/Y_0}\tau_{s-1}(P)
 \le\frac N{Y_0}(1+L)^{s-2}.
\]
This proves (B.14). \hfill$\square$

\subsection*{Boundary and logarithm-freezing estimates}

For every shell meeting the source,
\[
 L-B\log L-\log2<\log X<L.
 \tag{B.15}
\]
Thus, since $C_{\rm box}$ is fixed, for every $\epsilon>0$,
$(\log X)^{-C_{\rm box}}\le(1+\epsilon)L^{-C_{\rm box}}$
once $N$ is large. Also $\log n$ and each nonempty box lower-corner
logarithm are $(1+o(1))L$ on the source.

A boundary box lies within $6\delta_X$ in logarithmic distance of a
face in (B.13). The source and product-shell faces give two singleton
families and the lattice $(\log2)\mathbb Z$. The cap faces give
$\log2+(\log2/3)\mathbb Z$. In each boundary tuple, an active
coordinate is at least $N^{1/7}$ for large $N$. For a product face,
the largest of at most six factors is at least
$(H/2)^{1/6}$. On a cap face the capped factor is at least
$(H/4)^{1/3}$. Apply Lemma~\ref{lem:aa37-thin-shell} with
$s\le6$, $Y_0=N^{1/7}$, and
$\eta_{\rm dist}\ll L^{-C_{\rm box}}$. The number of face families is
finite and independent of the number of boxes. Let $\mathcal U$ be the
set of boundary indices $(X,j,\mathbf v)$ across shells and patterns
for which $\mathbf v\in T_{X,j}$ lies in a boundary box. It excludes
ambient Cartesian points outside the source-restricted set (B.10).
The number of boundary indices is therefore
\[
 \#\mathcal U\ll N L^{5-C_{\rm box}}+N^{6/7}L^5
 \ll N L^{5-C_{\rm box}}.
 \tag{B.16}
\]

The elementary divisor estimates
\[
 \begin{aligned}
 \sum_{n<N}\tau_6(n)&\le N(1+L)^5,\\
 \sum_{n<N}\tau_6(n)F_a(n)^2&\ll N(1+L)^{26},\\
 \sum_{n<N}\tau_6(n)F_a(n)&\ll N(1+L)^{19}
 \end{aligned}
 \tag{B.17}
\]
follow from $\tau_6^2\le\tau_{36}$,
$d^2\le\tau_4$, $d^4\le\tau_{16}$, the divisor-sum bound
$\sum_{n\le x}\tau_k(n)\le x(1+\log x)^{k-1}$, and
Cauchy--Schwarz. For the last bound, apply Shiu's short-interval
estimate \cite{Shiu80} to $d(m)^2\le\tau_4(m)$:
$\sum_{a-N<m\le a}d(m)^2\ll N(1+L)^3$ uniformly for
$N\le a\le NL/3$. For $a\le2N$ the global divisor-sum bound
suffices. For $a>2N$ the interval length $N$ exceeds $a^{2/3}$
for all sufficiently large $N$. Combining this with
$\sum_{n<N}\tau_6(n)^2\ll N(1+L)^{35}$ gives the exponent $19$.
The bound for $F_a^2$ follows by summing $d(m)^4$ up to $a\le NL/3$.
In particular, the LOG part from
replacing
$1/\log n$ by the lower-corner value $1/\lambda_{\mathcal B}$ on a
box satisfies
\[
 \left|E_{\rm LOG}\right|\ll_\sigma NL^{18-C_{\rm box}}.
 \tag{B.18}
\]
Here $0\le\log n-\lambda_{\mathcal B}\le6\delta_X$,
$|c_j(\mathbf v)|\le3L$, and
$|1/\log n-1/\lambda_{\mathcal B}|\ll L^{-C_{\rm box}-2}$.
Sum the resulting $O(L^{-C_{\rm box}-1})F_a(n)$ over the tuples using
(B.17). A tuple belongs to exactly one box, so there is no box-count
factor.

For the boundary part, $|c_j(\mathbf v)|/\lambda_{\mathcal B}
\le3+o(1)$. Cauchy--Schwarz, (B.16), and the second estimate in (B.17)
give
\[
 |E_{\rm boundary}|\ll_\sigma NL^{16-C_{\rm box}/2}.
 \tag{B.19}
\]

\subsection*{Classification and interval lengths}

The thresholds in (B.4), together with
\[
 \sigma L\ge202(B\log L+\log2),\qquad
 \sigma L\ge202(\log2+6),
 \tag{B.20}
\]
ensure $\omega_X:=\log_N(N/X)\le\sigma/202$ and
$\kappa_X:=(\log2+6\log(1+\Delta_X))/L\le\sigma/202$. In this
appendix $\kappa_X$ plays the role of the $\kappa$ in (3.98).
For a box meeting $T_{X,j}$,
\[
 \begin{gathered}
 1-\omega_X-\kappa_X\le\sum_i e_i\le1-\omega_X+\kappa_X,\\
 e_i\le\frac{1-\omega_X}{3}+\kappa_X
 \quad\text{on each M\"obius variable}.
 \end{gathered}
 \tag{B.21}
\]
Here $\prod_i V_i\ge(\prod_i v_i)/(1+\Delta_X)^6$ and
$\prod_i V_i\le\prod_i v_i<2X$. The M\"obius bound follows from
$V_i\le m_i\le2X^{1/3}$. The corner error is independent of
the lower-cut exponent $B$.

Use the window
\[
 \mathcal W=[60\sigma,\,1/3+14\sigma].
\]
A box is Type II if some $e_i\in\mathcal W$. Select the largest such
$e_i$, with ties broken by the fixed variable order. Suppose no exponent
lies in $\mathcal W$. A box is Type I-1 when it has exactly one large
variable and that variable is not a M\"obius variable. It is Type I-2
when it has exactly two large variables and neither is a M\"obius
variable. There are no other nonboundary boxes. In fact, a large
M\"obius variable would imply
 $14\sigma\le\kappa_X-\omega_X/3\le\sigma/202$. If all variables were
small, (B.21) would give
 $1-\sigma/101\le360\sigma$. If three were large, it
would give $1+42\sigma\le1+\sigma/202$. Each is impossible.

\begin{lemma}[Lengths on retained Type-II boxes]
\label{lem:aa37-typeII-lengths}
Let $\mathcal B$
be a nonboundary Type-II box. Its selected cell is
\[
 [T,r_XT),\qquad r_X=1+\Delta_X,
\]
Let $M$ be the product of the other lower endpoints. At least one
complementary cell is nonunit, since the source product exceeds
$H>N^{1/3+15\sigma}$ eventually. There are $1\le r\le5$ such cells.
Set
\[
 I_T=(T,r_XT],\quad
 N_I=\lfloor r_XT\rfloor-\lfloor T\rfloor,\quad
 M_I=M(r_X^r-1).
 \tag{B.22}
\]
Then, for sufficiently large $N$,
\[
 N_I\asymp T L_X^{-C_{\rm box}},\qquad
 M_I\asymp (X/T)L_X^{-C_{\rm box}},
 \tag{B.23}
\]
The PSDT interval $I_T$ is contained in $(T,2T]$. It has exactly
the same selected integer coefficients as the half-open cell
$[T,r_XT)$.
\end{lemma}

\emph{Proof.} The real length of the selected cell is
$h_N=T\Delta_X$. Both
\[
 \lceil r_XT\rceil-\lceil T\rceil
 \quad\hbox{and}\quad
 N_I
\]
differ from $h_N$ by less than $1$. The only integer endpoints at
which the half-open cell and $I_T$ can differ are powers of two.
Every internal mesh endpoint is $2^a2^{h/J_X}$ with
$0<h<J_X$, and is irrational. Write $h/J_X=u/v$ in lowest terms with
$v>1$. If $2^{u/v}$ were rational, its $2$-adic valuation would give
$u=v\,v_2(2^{u/v})$, contradicting $(u,v)=1$. At powers of two the
selected rough coefficient vanishes, since $z_0>2$. Thus the two
integer conventions agree after multiplication by the coefficient.

Put
\[
 \sigma L\ge2^{130}(C_{\rm box}+1)(\log L+1).
 \tag{B.24}
\]
Then $h_N\ge T/(3L_X^{C_{\rm box}})>2$, so
$N_I\ge T/(6L_X^{C_{\rm box}})$ and
$T/N_I\le6L_X^{C_{\rm box}}\le6L^{C_{\rm box}}$. Also
$\prod_i V_i\asymp X$, hence $M\asymp X/T$ and
$M_I\asymp L_X^{-C_{\rm box}}M$, proving (B.23).
Write $x=\log_NX=1-\omega_X$, $\xi=\log_NT$, and
$\Sigma_2=\log_NS$, $\Sigma_3=\log_NR$ for the slot scales in
Lemma~\ref{lem:aa37-pack}. The scale $Q_{\rm slot}$ is the first-slot
scale in that lemma; $Q_{\rm supp}$ is the modulus cutoff in (B.5).
The inequalities in (4.9) and $M_I\asymp XT^{-1}L_X^{-C_{\rm box}}$
give
\[
\begin{aligned}
 \log_N(Q_{\rm slot}/T)
 &\le-\left(1+\frac{7271}{256}\right)\sigma<-29\sigma x,\\
 2\xi+\Sigma_2+2\Sigma_3&\le1-60\sigma,\\
 \log_N\frac{M_I X^{-89\sigma}}{Q_{\rm slot}RS}
 &\ge x-2\xi-\Sigma_2-\Sigma_3
   +\left(1+\frac{7271}{256}\right)\sigma
   -89\sigma x-o(1)\\
 &\ge \left(\frac{7271}{256}-28\right)\sigma
   -\omega_X+\Sigma_3+89\sigma\omega_X-o(1)>0.
\end{aligned}
\tag{B.25}
\]
Here $x=1-\omega_X$, $\omega_X\le\sigma/202$, and
$\log L_X/L=o(1)$ as $N\to\infty$ for fixed $\sigma,C_{\rm box}$.
The terms $o(1)$ in (B.25) are taken in this sense. For fixed $\sigma$
they are eventually smaller than $\sigma/1000$. Thus
$Q_{\rm slot}<TX^{-29\sigma}$ and
$Q_{\rm slot}RS<M_IX^{-89\sigma}$.
It follows that
\[
 N_I\ge Q_{\rm slot}X^{\sigma/3},\qquad M_I\ge X^{1/2}
 \tag{B.26}
\]
for all sufficiently large $N$. For the first inequality, the ratio
is at least
\[
 (6L_X^{C_{\rm box}})^{-1}X^{(29-1/3)\sigma};
\]
for the second, use the bound $T\le X^{1/3+15\sigma}$ of (B.30) in (B.23).

Let $m_0=M(1+r_X^r)/2$, and let $\phi_0$ be the fixed smooth cutoff in
(3.9). Then
\[
 \psi_I(m)=\phi_0\!\left(\frac{m-m_0}{M_I}\right)
\]
equals $1$ on $[M,Mr_X^r]$ and is supported on
$[M-M_I/2,Mr_X^r+M_I/2]$. It majorizes the full complementary
interval without changing its length parameter.

Since $T\ge N^{60\sigma}$ and $L_X\le L$, (B.24) gives
\[
 T^{2^{-126}}\ge12L_X^{C_{\rm box}}.
\]
Together with $N_I\ge T/(6L_X^{C_{\rm box}})$ this proves the
integer-length condition
\[
 N_I\ge(2T)^{1-2^{-126}}.
 \tag{B.27}
\]
This is the selected interval length hypothesis (LEN) in Section~3.2.
\hfill$\square$

The source cuts are constant on each retained box by the definition
of a nonboundary box. Let $n_0$ be a tested product in a retained
Type-II box, and let $T,M,r$ be as in
Lemma~\ref{lem:aa37-typeII-lengths}. Every selected integer is at most
$r_XT$. The complementary cutoff $\psi_I$ is supported below
$M[1+\tfrac32(r_X^r-1)]$. As $TM\le n_0$ and $1\le r\le5$,
their product is at most
\[
 n_0r_X[1+\tfrac32(r_X^r-1)]
 \le n_0(1+11\delta_X)
 \le(1-\varepsilon/2)N
\]
for sufficiently large $N$, using $\Delta_X\le\delta_X\to0$.
Thus on this enlarged support the upper cut leaves
\[
 a-mn\ge\varepsilon N/2\ge X(\log X)^{-C_{\rm box}}.
 \tag{B.28}
\]
The final inequality follows from $X<N$ and
$\delta_X=(\log X)^{-C_{\rm box}}<\varepsilon/2$.
Thus the resonant class is absent. Also
$a\le X^{1+\sigma/20}$ for all sufficiently large $N$.

For each fixed Iwaniec vector, Lemma~\ref{lem:aa37-pack} and the exact slot identity
(4.13) give the three-factor coefficient required by (3.5).
The smooth-only term has all three prime slots empty. Every tail factor
is at least $Y=N^{\theta\sigma^2}>N^{\ell_0\sigma^2}=u_0$.
This is the roughness threshold $u_0$ of (3.2), which bounds the tail
cutoff $u_L$ in C1. The smooth divisor is attached to the first
slot. The Section~4 packing lemma verifies the spectral inequality
(3.7) directly, as displayed in (4.9). Separately, (B.25)--(B.26)
verify the modulus and interval inequalities in (3.6); (B.27) gives
the integer-length condition (LEN), and (B.30) gives the selected-window
bounds used below. Also
\[
 \log_XN=\frac1{1-\omega_X}\le1+\frac{\sigma}{20},
\]
and $\log_XT\le1/3+15\sigma$ by the first inequality in (B.30) below.
The selected window also gives $X^{60\sigma}\le T$.

\subsection*{The interior estimate}

Let $E_{\rm int}$ be the sum of the HB contributions on retained
Type-I-1, Type-I-2, and Type-II boxes. It is taken after the finite
Iwaniec vector sum.
The arguments above prove
\begin{multline*}
 \left|\sum_qw_N^\nu(q)\Delta_D(q;a)\right|
 \\
 \le C_\sigma\bigl(
 NL/Y+N^{1/2+o(1)}+NL^{18-C_{\rm box}}\\
       +NL^{16-C_{\rm box}/2}\bigr)+|E_{\rm int}|.
 \tag{B.29}
\end{multline*}
The first two errors in (B.29) have fixed power savings. The choice
$C_{\rm box}\ge\lceil2A+32\rceil$ in Theorem~\ref{thm:aa37-interior}
makes the last two
$O(NL^{-A})$. Both estimates are over the union of tuples.

The tested measure $\mu_D$ is the prime indicator on the interval
of (B.1). It does not depend on the Iwaniec vector. The sieve factors
belong to the modulus coefficient in (W). The prime-detection and
rough Heath--Brown identities are (B.7)--(B.10). Every
variable, including the logarithmic one, is $z_0$-rough. The
prime-power terms are estimated before the box partition.

The exponent constraints (B.20)--(B.21) imply that the residual cells
are empty by the three contradictions above. The Type-II selected exponent
obeys
\[
 \begin{aligned}
 \log_XT&\le\frac{1/3+14\sigma}{1-\omega_X}
 \le1/3+15\sigma,\\
 \frac{1-\omega_X}{2}-\log_NT
 &\ge\frac16-(15+1/202)\sigma>0.
 \end{aligned}
 \tag{B.30}
\]
The second line records $T<X^{1/2}$, the Type-II upper-length condition;
the first line supplies the sharper bound used in (B.23).
By (B.2), its extension to $\widetilde w_N^\nu$ and $\widehat w_N^\nu$
after (B.5), and the parameter choice in Theorem~\ref{thm:aa37-interior},
$q<N^{e_Q}$ with $e_Q<\Theta(u)\le\Theta(0)=153/256<0.6$.
For a Type-I-1 box, (3.100) with $e_Q<153/256$ gives
$e_*-e_Q-(1-\omega_X)/10>0.3023$.
Take a Type-I-2 box. At most four of its at most six variables are
small, each with $e_i<60\sigma$. Their product $m$ has upper endpoint at
most $N^{240\sigma}(1+\Delta_X)^4\le N^{240\sigma+\kappa_X}$, so
$\log_NM\le240\sigma+\kappa_X$. Since $\log_NX=1-\omega_X$, $q<N^{e_Q}$,
$\kappa_X,\omega_X\le\sigma/202$ by (B.20), and
$0<\sigma\le10^{-10}$,
\[
 \begin{aligned}
 \log_XM&\le\frac{240\sigma+\kappa_X}{1-\omega_X}
 \le\frac{(240+1/202)\sigma}{1-\sigma/202}
 <2.401\cdot10^{-8},\\
 \log_Xq&<\frac{e_Q}{1-\omega_X}
 <\frac{153/256}{1-\sigma/202}<0.5977.
 \end{aligned}
\]
Also $N_1,N_2\ge X^{1/3}$ by the
large-variable definition, and $N_1N_2M\le2^{30}X$ by the
Heath--Brown product support and its dyadic scale convention.
Hence (3.101) holds, and with (B.30) and (3.100) the hypotheses
(3.99)--(3.101) of the Type-I lemmas hold.

The packing lemmas give, for each Iwaniec vector, the exact
three-slot coefficient and the packing inequalities (3.6)--(3.7).
They include the empty-tail term and retain repeated-prime
multiplicity before the mask error (B.6). The selected rough unit,
M\"obius and normalized-log coefficients are the three classes in
(3.8). The proof of (3.25) in Section~3 applies to each class, with
every auxiliary coprimality parameter $d$ and every subinterval.
Its bound is $T\tau(d)L^{-A_{\rm SW}}$ for every fixed
$A_{\rm SW}>0$. Its application to $\alpha^{(r)}$ and the small-norm
case in (3.47)--(3.48) give the norm-normalized BFI passage. The
complementary short-interval second moment and the remaining
dispersion steps are those of Theorem~\ref{thm:main}. On these boxes,
(B.22)--(B.28) verify their interval, smoothing and nonresonance data.
In the local theorem take $K_{\rm scale}=B+1$. In particular
$X>N/(2L^B)$ and $a\le NL/3$ imply
$a\le(2/3)XL^{B+1}\le X^{1+\sigma/20}$ for large $N$. The
comparator is the natural PSDT comparator on every HB product. No
coprimality of a composite product is asserted.

For the chosen width, the global box count in (3.91) has
$K_0=24C_{\rm box}+26$. The source labels, Iwaniec vectors and prefix
cells number $O_\sigma(1)$, and the HB coefficients are $O(1)$. There
are at most six variables. So there are $O(L^6)$ dyadic tuples, within
the $O(L^{24})$ of \S\ref{sec:pi-global}, and each dyadic tuple contains
$O(L^{6C_{\rm box}})$ short boxes, within the $O(L^{24C_{\rm box}})$ there.
This verifies (H7) for the prime source. Choose the local Type-II precision
$A'=K_0+A+5=A+24C_{\rm box}+31$ and
$A_{\rm work}=A'+4$, as in (3.92).
Then Theorem~\ref{thm:main} with (H7), that is (3.92), bounds the sum
over all Type-II boxes. The bound is $O(NL^{-A-5})$ plus a fixed
power-saving term. Lemmas~\ref{lem:typeI1}
and \ref{lem:typeI2} give $O(NL^{-A})$ for Type-I-1 and Type-I-2
when requested with $A_i=K_I+A+5$ and
$K_I=6(C_{\rm box}+1)+2$. Thus
\[
 |E_{\rm int}|\ll NL^{-A}
       +N^{1-c_\sigma}L^{K_0+C},
 \qquad c_\sigma=\ell_0\sigma^2/4>0.
 \tag{B.31}
\]
Here $C=C(\sigma,C_{\rm box},A',K_{\rm scale},\mathcal C_{\rm src})$
is the logarithmic exponent from (3.13).
Here the local spectral, restoration, endpoint and slot-localization
costs are included once in Theorem~\ref{thm:main}.

\begin{proposition}[Iwaniec prime-source estimate]\label{prop:aa37-prime-source}
Take the parameters of Theorem~\ref{thm:aa37-interior}, either Iwaniec
sign $\nu$, and a sufficiently large even $N$. Let $w_N^\nu$ be
the exact coefficient (W). Let $a$ satisfy the residue conditions of
Theorem~\ref{thm:aa37-interior}, and let $\Delta_D$ be (B.1).
First split off the LOG correction algebraically. Then every remaining
source summand is assigned exactly once to one class. The classes are
mask, proper-power, boundary, Type-II, Type-I-1 and Type-I-2.
Every retained box has the scale and coefficient properties stated
above, and
\[
 \left|\sum_qw_N^\nu(q)\Delta_D(q;a)\right|
 \ll_{\sigma,A,u,\theta,B,\varepsilon,C_{\rm box}} NL^{-A}.
 \tag{B.32}
\]
\end{proposition}

\emph{Proof.} The Iwaniec mask partition and support are (B.2)--(B.6).
The explicit identities (B.7)--(B.11) partition the remaining
tested measure into proper powers and HB tuples. The equal-ratio mesh
puts each tuple in exactly one box. A face-straddling box belongs to
the boundary class. On every other box both source predicates are constant. The
exponent-polytope argument above puts every nonempty such box in
exactly one of II, I-1 and I-2. The displayed tie rule selects
one II factor. LOG is split off algebraically before that assignment.
The errors are (B.6), (B.8), (B.18), (B.19) and (B.31). Their sum is
$O_{\sigma,A,u,\theta,B,\varepsilon,C_{\rm box}}(NL^{-A})$. \hfill$\square$

\endgroup


\section*{Use of AI tools}
The author devised the method of proof.
Codex (OpenAI) and Claude Code (Anthropic) are the AI tools used.
They assisted with derivation checks, code, typesetting, and review.
The author checked every proof step, table, and reference.
The author takes full responsibility for the content of the paper.

\section*{Statements and declarations}
\textbf{Competing interests.} The author declares no competing interests.

\textbf{Funding.} No external funding was received.

\textbf{Data availability.} No external data or numerical computations are used in the proof. Appendix~A records analytic source checks and their scope.

\begin{thebibliography}{99}
\bibitem{Lichtman23} J. D. Lichtman,
Primes in arithmetic progressions to large moduli, and Goldbach beyond
the square-root barrier,
\href{https://arxiv.org/abs/2309.08522v1}{arXiv:2309.08522v1}, 2023
(with an appendix by S. Drappeau and J. D. Lichtman).
\bibitem{Bombieri65} E. Bombieri, On the large sieve,
\emph{Mathematika} \textbf{12} (1965), no.~2, 201--225.
DOI: \href{https://doi.org/10.1112/S0025579300005313}{10.1112/S0025579300005313}.
\bibitem{Vinogradov65} A. I. Vinogradov,
The density hypothesis for Dirichlet $L$-series (in Russian),
\emph{Izv. Akad. Nauk SSSR Ser. Mat.} \textbf{29} (1965), no.~4,
903--934.
MR0197414.
\bibitem{Shiu80} P. Shiu,
A Brun--Titchmarsh theorem for multiplicative functions,
\emph{J. Reine Angew. Math.} \textbf{313} (1980), 161--170.
DOI: \href{https://doi.org/10.1515/crll.1980.313.161}{10.1515/crll.1980.313.161}.
\bibitem{DI82} J.-M. Deshouillers and H. Iwaniec,
Kloosterman sums and Fourier coefficients of cusp forms,
\emph{Invent. Math.} \textbf{70} (1982), no.~2, 219--288.
DOI: \href{https://doi.org/10.1007/BF01390728}{10.1007/BF01390728}.
\bibitem{BFI86} E. Bombieri, J. B. Friedlander and H. Iwaniec,
Primes in arithmetic progressions to large moduli,
\emph{Acta Math.} \textbf{156} (1986), 203--251.
DOI: \href{https://doi.org/10.1007/BF02399204}{10.1007/BF02399204}.
\bibitem{MaynardII} J. Maynard,
Primes in arithmetic progressions to large moduli II: Well-factorable estimates,
\emph{Mem. Amer. Math. Soc.} \textbf{306} (2025), no.~1543.
DOI: \href{https://doi.org/10.1090/memo/1543}{10.1090/memo/1543};
the theorem numbering cited here is from
\href{https://arxiv.org/abs/2006.07088v1}{arXiv:2006.07088v1}.
\bibitem{Pascadi25} A. Pascadi,
On the exponents of distribution of primes and smooth numbers,
\href{https://arxiv.org/abs/2505.00653v2}{arXiv:2505.00653v2}, 2025.
\bibitem{Drappeau} S. Drappeau,
Sums of Kloosterman sums in arithmetic progressions, and the error term
in the dispersion method,
\emph{Proc. Lond. Math. Soc.} \textbf{114} (2017), no.~4, 684--732.
DOI: \href{https://doi.org/10.1112/plms.12022}{10.1112/plms.12022};
\href{https://arxiv.org/abs/1504.05549v4}{arXiv:1504.05549v4}.
\bibitem{AA37} H. Iwaniec, A new form of the error term in the linear sieve,
\emph{Acta Arith.} \textbf{37} (1980), 307--320.
DOI: \href{https://doi.org/10.4064/aa-37-1-307-320}{10.4064/aa-37-1-307-320}.
\bibitem{AA36} H. Iwaniec, Rosser's sieve,
\emph{Acta Arith.} \textbf{36} (1980), no.~2, 171--202.
DOI: \href{https://doi.org/10.4064/aa-36-2-171-202}{10.4064/aa-36-2-171-202}.
\bibitem{HB82} D. R. Heath-Brown,
Prime numbers in short intervals and a generalized Vaughan identity,
\emph{Canad. J. Math.} \textbf{34} (1982), no.~6, 1365--1377.
DOI: \href{https://doi.org/10.4153/CJM-1982-095-9}{10.4153/CJM-1982-095-9}.
\bibitem{KS03} H. H. Kim,
Functoriality for the exterior square of $GL_4$ and the symmetric fourth of $GL_2$,
\emph{J. Amer. Math. Soc.} \textbf{16} (2003), 139--183;
Appendix~2 by H. H. Kim and P. Sarnak, pp.~175--181.
DOI: \href{https://doi.org/10.1090/S0894-0347-02-00410-1}{10.1090/S0894-0347-02-00410-1}.
\bibitem{Deligne74} P. Deligne, La conjecture de Weil.~I,
\emph{Publ. Math. IH\'ES} \textbf{43} (1974), 273--307.
DOI: \href{https://doi.org/10.1007/BF02684373}{10.1007/BF02684373}.
\bibitem{Weil48} A. Weil, On some exponential sums,
\emph{Proc. Nat. Acad. Sci. USA} \textbf{34} (1948), 204--207.
DOI: \href{https://doi.org/10.1073/pnas.34.5.204}{10.1073/pnas.34.5.204}.
\bibitem{Estermann61} T. Estermann,
On Kloosterman's sum,
\emph{Mathematika} \textbf{8} (1961), no.~1, 83--86.
DOI: \href{https://doi.org/10.1112/S0025579300002187}{10.1112/S0025579300002187}.
\bibitem{IK04} H. Iwaniec and E. Kowalski,
\emph{Analytic Number Theory}, AMS Colloquium Publications \textbf{53},
AMS, Providence, RI, 2004.
DOI: \href{https://doi.org/10.1090/coll/053}{10.1090/coll/053}.
\bibitem{FI78} J. Friedlander and H. Iwaniec,
On Bombieri's asymptotic sieve,
\emph{Ann. Scuola Norm. Sup. Pisa Cl. Sci. (4)} \textbf{5} (1978), 719--756.
\bibitem{FouvryIwaniec80} E. Fouvry and H. Iwaniec,
On a theorem of Bombieri--Vinogradov type,
\emph{Mathematika} \textbf{27} (1980), no.~2, 135--152.
DOI: \href{https://doi.org/10.1112/S0025579300010032}{10.1112/S0025579300010032}.
\bibitem{FouvryIwaniec83} E. Fouvry and H. Iwaniec,
Primes in arithmetic progressions,
\emph{Acta Arith.} \textbf{42} (1983), no.~2, 197--218.
DOI: \href{https://doi.org/10.4064/aa-42-2-197-218}{10.4064/aa-42-2-197-218}.
\bibitem{MaynardI} J. Maynard,
Primes in arithmetic progressions to large moduli I: Fixed residue classes,
\emph{Mem. Amer. Math. Soc.} \textbf{306} (2025), no.~1542.
DOI: \href{https://doi.org/10.1090/memo/1542}{10.1090/memo/1542}.
\bibitem{MaynardIII} J. Maynard,
Primes in arithmetic progressions to large moduli III: Uniform residue classes,
\emph{Mem. Amer. Math. Soc.} \textbf{306} (2025), no.~1544.
DOI: \href{https://doi.org/10.1090/memo/1544}{10.1090/memo/1544}.
\bibitem{DPR23} S. Drappeau, K. Pratt and M. Radziwi{\l}{\l},
One-level density estimates for Dirichlet $L$-functions with extended support,
\emph{Algebra Number Theory} \textbf{17} (2023), no.~4, 805--830.
DOI: \href{https://doi.org/10.2140/ant.2023.17.805}{10.2140/ant.2023.17.805}.
\bibitem{ABL21} E. Assing, V. Blomer and J. Li,
Uniform Titchmarsh divisor problems,
\emph{Adv. Math.} \textbf{393} (2021), article 108076.
DOI: \href{https://doi.org/10.1016/j.aim.2021.108076}{10.1016/j.aim.2021.108076}.
\bibitem{Lichtman21} J. D. Lichtman,
A modification of the linear sieve, and the count of twin primes,
\emph{Algebra Number Theory} \textbf{19} (2025), no.~1, 1--38.
DOI: \href{https://doi.org/10.2140/ant.2025.19.1}{10.2140/ant.2025.19.1};
\href{https://arxiv.org/abs/2109.02851v2}{arXiv:2109.02851v2}.
\bibitem{Kuznetsov81} N. V. Kuznetsov,
Petersson's conjecture for cusp forms of weight zero and Linnik's conjecture.
Sums of Kloosterman sums,
Original Russian version: \emph{Mat. Sb. (N.S.)} \textbf{111(153)}
(1980), no.~3, 334--383.
English transl., \emph{Math. USSR-Sb.} \textbf{39} (1981), no.~3, 299--342.
DOI: \href{https://doi.org/10.1070/SM1981v039n03ABEH001518}{10.1070/SM1981v039n03ABEH001518}.
\bibitem{Bruggeman78} R. W. Bruggeman,
Fourier coefficients of cusp forms,
\emph{Invent. Math.} \textbf{45} (1978), no.~1, 1--18.
DOI: \href{https://doi.org/10.1007/BF01406220}{10.1007/BF01406220}.
\bibitem{Linnik63} Yu. V. Linnik,
The dispersion method in binary additive problems,
translated by S. Schuur, Translations of Mathematical Monographs \textbf{4},
American Mathematical Society, Providence, RI, 1963.
Reprint record, DOI: \href{https://doi.org/10.1090/mmono/004}{10.1090/mmono/004}.
\bibitem{DLMF} NIST Digital Library of Mathematical Functions,
\S10.43(v), equations (10.43.30)--(10.43.31),
\url{https://dlmf.nist.gov/10.43.v}.
\bibitem{OperaCribro} J. Friedlander and H. Iwaniec,
\emph{Opera de Cribro}, AMS Colloquium Publications \textbf{57}, 2010.
DOI: \href{https://doi.org/10.1090/coll/057}{10.1090/coll/057}.
\end{thebibliography}
\end{document}
